% =====================================================================
%  Einstein Dynamics from Renewal Memory: Ricci Reconstruction,
%  de Sitter Geometry, and Structural Dark-Energy Constraints
%  (ATMP-oriented reorganization; theorem content and proofs retained)
% =====================================================================
\documentclass[11pt]{article}
\usepackage[T1]{fontenc}
\usepackage[margin=1.05in]{geometry}
\usepackage{amsmath,amssymb,amsthm,mathtools}
\usepackage{booktabs,array}
\usepackage{enumitem}
\usepackage[colorlinks=true,linkcolor=blue,citecolor=blue,urlcolor=blue]{hyperref}
\usepackage[nameinlink,capitalise]{cleveref}
\usepackage{comment}
\setlength{\emergencystretch}{3em}
\hbadness=10000
\raggedbottom

\newtheorem{theorem}{Theorem}[section]
\newtheorem{maintheorem}[theorem]{Main Theorem}
\newtheorem{proposition}[theorem]{Proposition}
\newtheorem{lemma}[theorem]{Lemma}
\newtheorem{corollary}[theorem]{Corollary}
\theoremstyle{definition}
\newtheorem{definition}[theorem]{Definition}
\newtheorem{hypothesis}[theorem]{Hypothesis}
\newtheorem{construction}[theorem]{Construction}
\newtheorem{principle}[theorem]{Principle}
\theoremstyle{remark}
\newtheorem{remark}[theorem]{Remark}
\crefname{maintheorem}{Main Theorem}{Main Theorems}
\Crefname{maintheorem}{Main Theorem}{Main Theorems}
\crefname{hypothesis}{Hypothesis}{Hypotheses}
\Crefname{hypothesis}{Hypothesis}{Hypotheses}
\crefname{principle}{Principle}{Principles}
\Crefname{principle}{Principle}{Principles}

\newcommand{\R}{\mathbb R}
\newcommand{\C}{\mathbb C}
\newcommand{\Z}{\mathbb Z}
\newcommand{\N}{\mathbb N}
\newcommand{\Ztwo}{\mathbb Z/2}
\newcommand{\A}{\mathcal A}
\newcommand{\B}{\mathcal B}
\newcommand{\Hh}{\mathcal H}
\newcommand{\K}{\mathcal K}
\newcommand{\M}{\mathcal M}
\newcommand{\W}{\mathcal W}
\newcommand{\CP}{\mathrm{CP}}
\newcommand{\eps}{\varepsilon}
\newcommand{\Tr}{\operatorname{Tr}}
\newcommand{\tr}{\operatorname{tr}}
\newcommand{\id}{\mathrm{id}}
\newcommand{\spec}{\operatorname{spec}}
\newcommand{\Vol}{\operatorname{Vol}}
\newcommand{\Area}{\operatorname{Area}}
\newcommand{\Ric}{\operatorname{Ric}}
\newcommand{\Scal}{\operatorname{Scal}}
\newcommand{\dS}{\mathrm{dS}}
\newcommand{\AdS}{\mathrm{AdS}}
\newcommand{\Tren}{T^{\mathrm{ren}}}
\newcommand{\J}{\mathcal J}
\newcommand{\Dirac}{\mathsf D}
\newcommand{\Lcal}{\mathcal L}
\newcommand{\Scr}{\operatorname{Scr}}
\newcommand{\bridge}{\mathrm{bridge}}
\newcommand{\hor}{\mathrm{hor}}
\newcommand{\bulk}{\mathrm{bulk}}
\newcommand{\rad}{\mathrm{rad}}
\newcommand{\grav}{\mathrm{grav}}
\newcommand{\eff}{\mathrm{eff}}
\newcommand{\ren}{\mathrm{ren}}
\newcommand{\GH}{\mathrm{GH}}
\newcommand{\etaScr}{\eta_{\mathrm{scr}}}
\newcommand{\Rmicro}{\mathcal R_{\mathrm{micro}}}
\newcommand{\Rcurv}{\mathcal R_{\mathrm{curv}}}
\newcommand{\Rcal}{\mathcal R_{\mathrm{cal}}}
\newcommand{\Rshear}{\mathcal R_{\mathrm{shear}}}
\newcommand{\Mren}{\mathcal M_{\mathrm{ren}}}
\newcommand{\mods}{\Delta}
\newcommand{\Kmod}{K_{\mathrm{mod}}}
\newcommand{\Lcurv}{L_{\mathrm{curv}}}
\newcommand{\Riem}{\operatorname{Riem}}
\newcommand{\SU}{\mathrm{SU}}
\newcommand{\Spin}{\mathrm{Spin}}
\newcommand{\Gren}{G_{\mathrm{ren}}}
\newcommand{\su}{\mathfrak{su}(2)}
\newcommand{\Ad}{\operatorname{Ad}}
\newcommand{\LC}{\mathrm{LC}}
\newcommand{\mom}{\mathrm{mom}}
\newcommand{\wideg}{\widehat g}
\newcommand{\rstar}{\rho_*}
\newcommand{\Lop}{\mathsf L}
\DeclareMathOperator{\Op}{Op}
\DeclareMathOperator{\dive}{div}
\DeclareMathOperator{\supp}{supp}
\DeclareMathOperator{\Pexp}{\mathcal P\!\exp}
\newcommand{\CV}{\mathfrak C_{\mathrm{CV}}}
\newcommand{\kren}{\kappa_{\ren}}
\newcommand{\Bstep}{B_{\mathrm{step}}}
\newcommand{\Bphys}{B_{\mathrm{phys}}}
\newcommand{\bartau}{\bar\tau}
\newcommand{\Qhat}{\widehat Q}
\newcommand{\Perr}{\mathcal P}
\newcommand{\Pphys}{\mathcal P_{\mathrm{phys}}}
\newcommand{\hscr}{h_{\mathrm{scr}}}
\newcommand{\Scond}{S_{\mathrm{cond}}}
\newcommand{\Npartial}{\nu_{\partial}}
\newcommand{\Aren}{A_{\ren}}
\newcommand{\Weyl}{C}
\newcommand{\St}{S}
\newcommand{\uu}{\mathbf u}
\newcommand{\kR}{\kappa_R}
\newcommand{\kS}{\kappa_S}
\newcommand{\kC}{\kappa_C}
\newcommand{\Pip}{\Pi_+}
\newcommand{\Ruu}{R_{uu}}
\newcommand{\Rth}{\mathcal R^{(3)}}
\newcommand{\OpW}{\operatorname{Op}^{\mathrm W}_\eps}
\DeclareMathOperator{\Symop}{Sym}
\DeclareMathOperator{\spn}{span}

\title{\textbf{Einstein Dynamics from Renewal Memory:\\
Ricci Reconstruction, de Sitter Geometry, and Structural Dark-Energy Constraints}}
\author{A.~P\'elissier}
\date{}

\begin{document}
\maketitle


\begin{abstract}
We develop a semiclassical theory in which Einstein dynamics emerges from renewal memory. Starting from a Lorentzian signed renewal geometry and its continuum Dirac limit, we show that an isolated massive band organizes successive geometric response orders: the large-deviation phase yields the proper-time action, while the Gaussian prefactor reconstructs the Ricci tensor from calibrated causal diamonds. This WKB–Ricci closure holds on a broad admissible class without assuming reversibility, detailed balance, or isotropy.
%
A single tilted matrix-renewal pressure controls the microscopic current, metric covariance, screen entropy, and higher cumulants. Predictive-screen thermodynamics then converts the Ricci response into the Einstein equation and fixes Newton’s constant from the screen entropy density. The homogeneous renewal current selects a spatially flat de Sitter branch and yields structural dark-energy constraints: reversible fractional relaxation cannot cross the phantom divide, while irreversible exchange is subject to explicit current, entropy-production, one-crossing, growth-suppression, and fluctuation-symmetry conditions.
%
The first finite-resolution correction is a four-derivative effective action determined by renewal cumulants; for an isotropic compound-Poisson channel, its Wilson coefficients lie on an explicit nontrivial ray. Thus Einstein dynamics, de Sitter geometry, dark-energy constraints, and higher-curvature corrections arise as successive response orders of one renewal process.
%
A \texttt{Lean} formalization is available at
\url{https://github.com/AI-MathPhys/renewal-geometry}.
\end{abstract}




\setcounter{tocdepth}{2}
\tableofcontents

% =====================================================================
\section{Introduction}
\label{sec:intro}
% =====================================================================

General relativity begins with a Lorentzian manifold and a metric whose dynamics are governed by the Einstein equation.  A central problem in mathematical physics is to understand whether these structures can instead arise as collective properties of a more primitive system.  Recovering only a causal cone or an effective metric is not sufficient.  A viable construction must also explain the emergence of proper-time propagation, local curvature, gravitational dynamics, the normalization of Newton's constant, and the first corrections beyond the Einstein regime.  The purpose of this paper is to show that these ingredients can arise, under explicit semiclassical and thermodynamic hypotheses, as successive response orders of a single renewal process.

Thermodynamic derivations provide an important precedent.  Jacobson showed that the Einstein equation follows from imposing the Clausius relation on all local Rindler horizons, with the Unruh temperature and an entropy proportional to area \cite{Jacobson,Unruh1976}.  Later developments connected this argument with entanglement equilibrium, non-equilibrium entropy production, and holographic or modular descriptions of local gravitational dynamics \cite{ElingGuedensJacobson,Padmanabhan,Verlinde2011,Lashkari,FaulknerEtAl,Jacobson2016}.  These results demonstrate that the Einstein tensor is strongly constrained by local thermodynamic consistency.  They do not, however, derive the Lorentzian metric, the causal diamonds on which the local response is measured, or the microscopic degrees of freedom whose entropy fixes the gravitational coupling.  Those structures are normally assumed before the thermodynamic argument begins.

Other approaches address different parts of the same problem.  Causal-set theory reconstructs curvature from order-theoretic observables, including nonlocal d'Alembertian operators and interval-chain statistics \cite{BombelliLeeMeyerSorkin,Surya,BenincasaDowker,RoySinhaSurya}.  Dynamical triangulations and group-field theories seek continuum geometry through collective phases of discrete building blocks \cite{AmbjornJurkiewiczLoll,Oriti}.  Noncommutative geometry encodes metric information spectrally, although the treatment of Lorentzian signature requires additional causal, Krein, or modular structure \cite{Connes,ChamseddineConnes,FrancoEckstein,Strohmaier}.  Loop quantum gravity and spin-foam models obtain discrete geometric spectra and horizon-state counting \cite{Barbero,Ashtekar1986,RovelliSmolin,AshtekarLewandowski,EPRL,BarrettAsymptotics,Immirzi,DomagalaLewandowski,Meissner,AshtekarBaezKrasnov}.  Effective field theory, independently of the microscopic completion, organizes corrections to Einstein gravity in a curvature expansion \cite{Stelle1977,Donoghue1994,Burgess2004}.  Stochastic gravity provides a related open-system description, but it begins with a prescribed semiclassical spacetime and studies fluctuations around it \cite{HuVerdaguer}.  The present construction differs in logical order: the renewal dynamics first determines the propagation law and its characteristic geometry, the fluctuation response of that propagation reconstructs curvature, and only then is the gravitational field equation assembled.

The microscopic object considered here is a renewal memory.  Its histories are sequences of regenerative events whose letters act on a finite internal observable algebra through completely positive channels.  Histories with the same future action are identified in a predictive quotient, in analogy with future-equivalence constructions in computational mechanics \cite{ShaliziCrutchfield}.  Renewal theory ordinarily describes repeated loss of memory within an already given temporal or geometric setting \cite{Feller}.  Here renewal is placed one level below that setting and is treated as the substrate from which temporal and geometric observables emerge.

The kinematic part of this programme was developed in the Lorentzian renewal-spacetime construction of \cite{PelissierLorentzian}.  There the positive predictive quotient carries dimension, volume, and divisibility data, while a signed anti-renewal cocycle supplies a Krein fundamental symmetry and a Lorentzian Clifford structure.  The corresponding discrete signed Dirac operators converge to a continuum Lorentzian Dirac operator, with strong-resolvent convergence globally and norm-resolvent control on every fixed physical momentum band.  That result establishes the microscopic origin of the Lorentzian principal symbol, but it does not derive a gravitational action, a curvature response, an Einstein equation, or a cosmological branch.  The present paper supplies this dynamical layer.

The first step is to show that the foundational operator limit survives the addition of the primitive reset term and passes to a uniformly isolated massive band.  The band has constant multiplicity and must therefore be treated by a matrix-valued semiclassical construction.  Its principal eigenvalue is nevertheless scalar, so the leading Hamilton--Jacobi problem is well defined.  We prove that the associated large-deviation phase is a proper-time-deficit action.  At constant renewal rate its minimizing characteristics are the timelike geodesics of the emergent metric; for slowly varying rates the deviation from the frozen geometric flow is quantitatively controlled.  The matrix Berry transport acts unitarily on the band and disappears from determinant and Hilbert--Schmidt modulus observables.

Curvature appears at the next response order.  The Gaussian factor in the renewal propagator is the van~Vleck--Morette determinant of the characteristic geometry.  By integrating its flat-normalized response over calibrated causal diamonds with two independent profiles, we reconstruct both the longitudinal Ricci contraction and the scalar curvature.  Varying the timelike axis then recovers the full Ricci tensor.  The reconstruction remains valid on a broad admissible Calder'on--Vaillancourt class.  No reversibility, detailed-balance, or isotropy assumption is required: the projected subprincipal drift separates into anti-Hermitian geometric transport, a removable gradient contribution, and an irreversible scalar whose reflection-averaged contribution is subleading in the small-diamond limit.

The geometric response is converted into a field equation through predictive-screen thermodynamics.  Primitive renewal transfer channels give a boundary-local entropy law, with the area density determined by a stationary conditional-entropy rate.  For quasi-free renewal-screen scaling limits, causal propagation and wedge isotony imply the half-sided modular structure that identifies the modular flow with the local geometric boost and fixes the Unruh normalization.  Combining this modular first law with the Ricci response and the Raychaudhuri area variation yields the local Einstein equation.  Newton's constant is not inserted independently: it is the inverse of four times the predictive-screen entropy density.  In the resolved predictive-pure specialization, the quantum screen entropy equals the entropy of the predictive record; for the uniform $(K_4)$ cell this entropy is $(\log 3)$.

The microscopic gravitational observables are organized by one tilted matrix-renewal pressure.  Its derivatives with respect to distinct tilts give the signed deficiency current, the Green--Kubo metric covariance, the predictive-screen entropy response, and the higher connected cumulants.  These quantities have a common statistical origin but are not numerically identified with one another.  In particular, the homogeneous cosmological variable is the signed current per physical renewal depth, rather than the dimensionless increment per regenerative event.  Keeping this clock conversion explicit is essential for relating the microscopic transfer process to the continuum Hubble scale.

In the homogeneous pure-deficiency sector, the physical renewal current fixes a positive cosmological constant and selects the expanding de~Sitter branch.  Retaining the FLRW curvature index in the Hamiltonian constraint shows that the spatially flat slicing is selected rather than assumed.  In four spacetime dimensions this selection is preserved by the complete local curvature-squared response throughout the controlled derivative regime.  The resulting local geometry has positive constant curvature.  Its identification with complete global de~Sitter space requires a separate predictive-completion principle, which is stated explicitly and not folded into the local theorem.

The same homogeneous current also leads to a sharp distinction between equilibrium vacuum energy and dynamical dark energy.  If the deficiency enters only through a vacuum tensor proportional to the metric, the contracted Bianchi identity forces it to be constant.  A rolling deficiency therefore belongs to an effective-fluid or open-exchange description, not to a second source-free vacuum branch.  At the critical renewal clock, the effective homogeneous mode obeys a fractional relaxation law of Caputo order (1/3).  Convex reversible relaxation is monotone and cannot cross the phantom divide.  An irreversible cycle current can open a transient reconstructed phantom interval only after exceeding an explicit restoring threshold, and the thermodynamic uncertainty relation imposes a nonzero entropy-production floor.

For the linearized hereditary response, we isolate a spectral condition under which the crossing topology becomes independent of amplitudes and cosmological fitting parameters.  When the positive exchange modes have strictly faster relaxation rates than the negative restoring modes, the derivative of the homogeneous current has at most one zero.  Any crossing is then necessarily directed from the phantom to the nonphantom side.  Positive matter-to-deficiency exchange also gives an exact comoving-matter deficit.  Under a minimal smooth and geodesic perturbation closure, a comparison theorem converts this deficit into suppression of linear growth at fixed expansion history.  When the microscopic current satisfies single-affinity local detailed balance, the tilted pressure obeys a Gallavotti--Cohen symmetry, leading to a fluctuation theorem for accumulated expansion and an exact hierarchy relating its mean, variance, skewness, higher cumulants, and entropy production.  These are structural predictions; the crossing redshift and the present dark-energy density still require microscopic clock and exchange calibration.

Beyond Einstein order, the finite renewal resolution produces a local four-derivative effective action.  Its operator basis is fixed by covariance, parity, field content, and the four-dimensional Gauss--Bonnet identity, while its coefficients are determined by normalized renewal cumulants.  For a specified isotropic compound-Poisson channel, we compute the curvature-squared volume response explicitly and obtain a nontrivial Wilson ray with a positive scalar-curvature coefficient and a smaller Weyl component.  The same expansion defines a dimensionless control parameter comparing the renewal scale with the local curvature radius.  Its loss of smallness marks the trans-renewal boundary at which no finite-order continuum truncation is predictive.  This boundary is a statement about the failure of the derivative expansion, not a proposed bounce or ultraviolet completion.

The construction therefore has a simple organizing principle.  The Lorentzian operator limit supplies the causal principal symbol.  The large-deviation phase supplies proper-time propagation.  The Gaussian response supplies Ricci curvature.  Predictive-screen thermodynamics supplies the Einstein equation and Newton normalization.  The homogeneous signed current supplies the de~Sitter and dark-energy sectors.  Higher renewal cumulants supply the first finite-resolution corrections.  These layers are not separate models joined phenomenologically; they are different derivatives and asymptotic orders of one renewal transfer process.

The paper is organized accordingly.  \Cref{sec:kinematics} introduces the matrix-renewal pressure, transfers the Lorentzian operator limit to the isolated massive band, constructs the matrix WKB propagator, identifies its proper-time action, and proves the admissible-channel subprincipal closure.  \Cref{sec:einstein-reconstruction} reconstructs Ricci curvature from causal diamonds, derives the predictive-screen entropy law and modular calibration, and assembles the local Einstein equation.  \Cref{sec:desitter} studies homogeneous renewal cosmology, including the de~Sitter branch, spatial-flatness selection, higher-curvature protection, global completion, perturbative stability, and horizon thermodynamics.  \Cref{sec:qg} develops the finite-resolution effective action, the cumulant-to-Wilson map, the explicit isotropic Wilson ray, Wald entropy, and the validity boundary of the derivative expansion.  \Cref{sec:dark-energy} treats dynamical renewal dark energy, including reversible no-crossing, irreversible exchange thresholds, the spectral one-crossing theorem, matter depletion, growth comparison, and expansion-current fluctuation relations.  \Cref{sec:interpretation} summarizes the logical status of the construction and the remaining microscopic closure problems.


% =====================================================================
\section{Renewal transfer operators and semiclassical propagation}
\label{sec:kinematics}
% =====================================================================

This section isolates the operator-theoretic and semiclassical input to the
gravity construction.  We first record the physical-depth Perron dictionary,
then transfer the foundational Lorentzian endpoint to the isolated massive
renewal band, construct its matrix WKB propagator, identify the proper-time
action and causal cone, and close the subprincipal analysis on the admissible
channel class.  Curvature reconstruction and the thermodynamic field equation
are deferred to \Cref{sec:einstein-reconstruction}.

\subsection{Matrix-renewal pressure and physical-depth observables}
\label{subsec:matrix-renewal-dictionary}

The gravitational variables used below are most transparent when the renewal
depth is kept inside the microscopic transfer object.  We therefore record the
matrix-renewal pressure theorem imported from the finite current--cycle layer
and spell out the normalization needed in the Einstein sector
\cite{PelissierCurrentsK4Arxiv}.

\begin{definition}[Tilted matrix-renewal kernel]
\label{def:tilted-matrix-renewal}
Let \(V\) be a finite set of predictive states and let
\[
   Q=\{Q_{xy}(n)\}_{x,y\in V,\ n\ge0},
   \qquad Q_{xy}(n)\ge0,
\]
be an irreducible aperiodic matrix-renewal kernel, normalized so that
\(\rho(\Qhat(0,0,0,0))=1\).  Let \(b(x,y)\) be a signed deficiency cochain,
\(\xi(x,y)\in\R^d\) a displacement cochain, and
\(\alpha(x,y)\) a screen-surprisal observable.  Define
\[
 \Qhat(z,\chi,k,\vartheta)_{xy}
 :=
 \sum_{n\ge0}Q_{xy}(n)
 \exp\{-zn+\chi b(x,y)+ik\cdot\xi(x,y)+\vartheta\alpha(x,y)\}.
\]
Write
\[
 F(z,\chi,k,\vartheta)
 :=
 \log\rho\!\left(\Qhat(z,\chi,k,\vartheta)\right),
 \qquad
 \bartau:=-\partial_zF(0,0,0,0).
\]
The number \(\bartau\) is the mean renewal depth in the Perron stationary
state.
\end{definition}

\begin{theorem}[Matrix-renewal Perron dictionary]
\label{thm:matrix-renewal-dictionary}
Assume that the entries of \(\Qhat\) are finite and analytic in a neighbourhood
of the origin and that \(0<\bartau<\infty\).  Then there is a unique real
analytic physical-depth pressure
\[
   \Pphys(\chi,k,\vartheta)
\]
near the origin such that
\[
 \rho\!\left(
 \Qhat(\Pphys(\chi,k,\vartheta),\chi,k,\vartheta)
 \right)=1,
 \qquad
 \Pphys(0,0,0)=0.
\]
Let
\[
 \Bstep:=\partial_\chi F(0),\qquad
 s_{\rm step}:=\partial_\vartheta F(0),
 \qquad
 C^{\rm step}_{ab}:=-\partial_{k_a}\partial_{k_b}F(0)
\]
after centering the displacement observable.  Then
\[
 \boxed{\;
 \Bphys:=\partial_\chi\Pphys(0)
       =\frac{\Bstep}{\bartau},
 \qquad
 s_{\rm phys}:=\partial_\vartheta\Pphys(0)
       =\frac{s_{\rm step}}{\bartau}. \;}
\]
Moreover
\[
 C^{\rm phys}_{ab}
 :=
 -\partial_{k_a}\partial_{k_b}\Pphys(0)
\]
is the physical-depth Green--Kubo covariance, and all higher derivatives of
\(\Pphys\) are the corresponding physical-depth connected cumulants.  Thus
current, metric covariance, entropy response, and higher homogenisation
tensors are different components of one analytic Perron jet; the theorem does
not identify their numerical values with one another.
\end{theorem}

\begin{proof}
At the origin the positive matrix \(\Qhat\) has a simple leading eigenvalue by
Perron--Frobenius theory.  Analytic perturbation theory therefore makes that
eigenvalue, and hence \(F\), analytic near the origin
\cite{Kato,EvansHoeghKrohn}.  Since
\(\partial_zF(0)=-\bartau\ne0\), the analytic implicit-function theorem applied
to
\[
   F(z,\chi,k,\vartheta)=0
\]
gives the unique analytic solution \(z=\Pphys(\chi,k,\vartheta)\).

Differentiating the identity
\(F(\Pphys(\chi,k,\vartheta),\chi,k,\vartheta)=0\) at the origin gives
\[
 -\bartau\,\partial_\chi\Pphys(0)+\partial_\chi F(0)=0,
 \qquad
 -\bartau\,\partial_\vartheta\Pphys(0)
       +\partial_\vartheta F(0)=0,
\]
which proves the first two formulas.  Twice differentiating in \(k\), after
centering so that \(\partial_k\Pphys(0)=0\), gives
\[
 -\bartau\,\partial_{k_a}\partial_{k_b}\Pphys(0)
 +\partial_{k_a}\partial_{k_b}F(0)=0
\]
when depth and displacement factor at the centered level; in the general
matrix-renewal case the repeated implicit differentiation adds precisely the
depth--displacement connected terms.  The resulting Hessian is the
Green--Kubo covariance per physical depth.  This is the standard tilted
operator characterization of additive-functional cumulants
\cite{Nagaev,KipnisVaradhan}.  Repeating the differentiation yields the higher
connected cumulants.
\end{proof}

\begin{corollary}[Embedded-chain specialization and metric clock]
\label{cor:embedded-renewal-dictionary}
Suppose
\[
   Q_{xy}(n)=\Pi_{xy}\psi(n)
\]
with finite mean \(\bartau\), and suppose the centered displacement and screen
observables depend only on the embedded jump.  Then
\[
 \Qhat(z,\chi,k,\vartheta)
 =\widehat\psi(z)L(\chi,k,\vartheta),
\]
and
\[
 \Bphys=\frac{\Bstep}{\bartau},
 \qquad
 C^{\rm phys}_{ab}=\frac{C^{\rm step}_{ab}}{\bartau},
 \qquad
 s_{\rm phys}=\frac{s_{\rm step}}{\bartau}.
\]
Defining the per-renewal metric length and the signed-cone speed by
\[
 \eps^2:=\frac1d\tr C^{\rm step},
 \qquad
 c_{\ren}:=\frac{\eps}{\bartau},
\]
one has
\[
 \boxed{\quad
 B=\Bphys=\frac{\Bstep c_{\ren}}{\eps}.
 \quad}
\]
\end{corollary}

\begin{proof}
Factorization gives
\[
 F(z,\chi,k,\vartheta)
 =\log\widehat\psi(z)
  +\log\rho(L(\chi,k,\vartheta)).
\]
The first-derivative formulas follow from
\Cref{thm:matrix-renewal-dictionary}.  Centering removes the first \(k\)
derivative, so the second implicit derivative is simply divided by
\(\bartau=-\partial_z\log\widehat\psi(0)\).  The final identity is the
definition \(c_{\ren}=\eps/\bartau\).
\end{proof}

\begin{remark}[Three quantities that must not be conflated]
\label{rem:three-renewal-rates}
The primitive reset rate \(\lambda\), the physical signed deficiency rate
\(B=\Bphys\), and the inverse mean renewal depth \(\bartau^{-1}\) are logically
different data.  They coincide only in special Poisson or one-lapse
normalizations.  All cosmological formulas in this paper use \(B=\Bphys\);
\(\Bstep\) is the signed increment per renewal event.
\end{remark}



\subsection{Lorentzian renewal input and the positive-sector obstruction}

The starting point is the Lorentzian operator geometry obtained in the
foundational construction.  Since the present paper is concerned with dynamics
rather than with reconstructing the kinematic layer, we first isolate exactly
which parts of that earlier result will be used.  The point is not merely that a
Lorentzian principal symbol exists: it is that the discrete signed renewal
operators already converge to it in the operator topologies appropriate to the
low-energy limit.  We then show that this convergence survives the addition of
the primitive reset sector that produces the massive renewal band.

\begin{theorem}[Imported Lorentzian renewal endpoint and operator limit
\cite{PelissierLorentzian}]
\label{thm:imported-lorentzian}
The signed renewal-spacetime construction supplies a Krein Hilbert space
$(\K,\J)$, a signed renewal Dirac operator $\Dirac_{\rm ren}$, a time element
$T$, and a Clifford principal symbol.  For mesh size $h_R\downarrow0$, its
finite-difference signed Dirac operators converge on a common Schwartz core to
the continuum Lorentzian Dirac operator.  The associated self-adjoint spatial
Hamiltonians $H_R$ converge in strong resolvent sense to
\[
 H_g=c\,e_a^i(x)\alpha^a(-i\partial_i)+H_{\rm ord},
\]
with strong convergence of propagators and spectral projections; on every fixed
physical momentum band $P_{\Lambda_0}$ the resolvents converge in norm, with the
second-order consistency estimate of the symmetric renewal differences.  The
same conclusion holds locally for slowly varying reset statistics in the
adiabatic curved regime.  The principal symbol is
\[
        h_g(x,\theta)=c\,e_a^i(x)\theta_i\alpha^a,
\]
and determines the spatial cometric
$g_x=\sum_a e_a(x)\otimes e_a(x)$ and the unit-lapse Lorentzian metric
$\mathbf g=-dt^2+g_{ij}(x)\,dx^idx^j$.  The foundational theorem does not
identify the ordering term $H_{\rm ord}$ with the Levi--Civita spin connection
and does not add the primitive reset/mass sector used below.
\end{theorem}

The theorem is imported rather than reproved.  It is stronger than a merely
formal symbol assignment: the Lorentzian Dirac principal symbol used below is
already the continuum limit of discrete signed renewal operators.  What must be
proved in the present paper is the stability of that limit under the bounded
primitive reset term and its passage to the gap-isolated band governing the
renewal large deviations.

\begin{theorem}[Foundational-to-dynamical transfer]
\label{thm:foundational-transfer}
Let $C_R$ denote the coarse-graining maps of
\Cref{thm:imported-lorentzian}, and fix a physical momentum band
$P_{\Lambda_0}$.  Let $\lambda_R\in C_b^3(X)$ satisfy
$\lambda_R\to\lambda$ in $C^3$, with
$\lambda_R,\lambda\ge\lambda_{\min}>0$, and let $\beta$ be a Hermitian
involution anticommuting with the spatial Clifford matrices.  Define the
microscopic massive renewal operators and their continuum limit by
\[
 \mathcal M_R=C_RH_RC_R^*+\lambda_R(\beta-\mathbf 1),
 \qquad
 \mathcal M=c\,e_a^i(x)\alpha^a(-i\partial_i)
             +\lambda(x)(\beta-\mathbf 1).
\]
Then, on $P_{\Lambda_0}L^2$, the bounded reset perturbation preserves the
foundational norm-resolvent limit:
\[
 \bigl\|P_{\Lambda_0}\bigl[(\mathcal M_R-z)^{-1}
       -(\mathcal M-z)^{-1}\bigr]P_{\Lambda_0}\bigr\|
 \longrightarrow0,
\]
locally uniformly for $z$ in compact subsets of the common resolvent set.  The
continuum principal symbol
\[
 M(x,\theta)=c\,e_a^i(x)\theta_i\alpha^a
             +\lambda(x)(\beta-\mathbf 1)
\]
has the two gap-separated branches
\[
 \sigma_\pm(x,\theta)=\pm\sqrt{\lambda(x)^2+c^2g_x(\theta,\theta)}-\lambda(x),
 \qquad
 \sigma_+-\sigma_-\ge2\lambda_{\min},
\]
with smooth Riesz projectors
\[
 \Pi_\pm(x,\theta)=\frac12\left(
 \mathbf 1\pm\frac{c\,e_a^i\theta_i\alpha^a+\lambda\beta}
 {\sqrt{\lambda^2+c^2g_x(\theta,\theta)}}\right).
\]
Consequently the lower-band Hamiltonian
$h(x,\theta)=\sqrt{\lambda^2+c^2g_x(\theta,\theta)}-\lambda$
and the principal-symbol, gap, and projector components of
\textup{(W1)--(W3)} are inherited from the foundational microscopic limit
rather than postulated.

If, in addition, microscopic order-$\eps$ corrections $B_{R,\eps}$ satisfy,
for a scale separation $h_R^2=o(\eps)$,
\[
 \left\|P_{\Lambda_0}\left(
 C_RB_{R,\eps}C_R^*-\operatorname{Op}^{\rm W}_\eps(\ell_1)
 \right)P_{\Lambda_0}\right\|_{H^1_\eps\to L^2}=o(1),
\]
then the full first-order semiclassical realization \textup{(W1)} follows on
every fixed physical momentum band.  After projection to the constant-rank band,
the geometric Berry terms are anti-Hermitian and hence invisible to the
unitary-invariant modulus response.  The remaining scalar calibration is
\[
  \kappa_{\ren}:=\frac1r\Re\Tr(\Pi_+\ell_1\Pi_+).
\]
In half-density normalization one has \(\kappa_{\ren}=0\) for the reversible
isotropic reset class of \Cref{hyp:reversible-isotropic}, as proved in
\Cref{thm:drift-cancellation}.  For a general admissible channel
\(\kappa_{\ren}\) need not vanish; the quantitative order-\(\eps\) symbol
remainder and the subleading character of the projected irreversible scalar
drift are established on the full admissible class in
\Cref{thm:unconditional-bare-remainder,thm:unconditional-wkb-ricci}.
\end{theorem}

\begin{proof}
On a fixed momentum band, \Cref{thm:imported-lorentzian} gives norm-resolvent
convergence of $C_RH_RC_R^*$ to its continuum Dirac Hamiltonian.  Multiplication
by $\lambda_R(\beta-\mathbf 1)$ is uniformly bounded and converges in operator
norm to $\lambda(\beta-\mathbf 1)$.  Stability of norm-resolvent convergence under
norm-convergent bounded self-adjoint perturbations, equivalently the second
resolvent identity and a Neumann argument, gives the first conclusion.  The
Clifford relations imply
\[
 \bigl(c\,e_a^i\theta_i\alpha^a+\lambda\beta\bigr)^2
 =\bigl(\lambda^2+c^2g_x(\theta,\theta)\bigr)\mathbf 1,
\]
which gives the branches, their uniform gap, and the displayed projectors.
Norm-resolvent convergence across an isolated contour implies convergence of the
corresponding Riesz projectors, so the isolated massive band is stable under the
microscopic limit.  The additional estimate on $B_{R,\eps}$, together with
$h_R^2=o(\eps)$, promotes the principal-symbol convergence to the stated
first-order semiclassical expansion.  The final assertion follows after projecting the order-\(\eps\) symbol to the
band: anti-Hermitian geometric transport drops from the determinant modulus,
while the normalized real trace of the reset contribution is the sole scalar
amplitude correction.  \Cref{thm:drift-cancellation} shows that this trace
vanishes for reversible isotropic resets; the admissible-class analysis below
shows that a nonzero irreversible trace is curvature-free at the relevant
order and contributes only subleading reflection-averaged error.
\end{proof}

\begin{remark}[Sharp scope of the transfer]
\label{rem:transfer-scope}
Full norm-resolvent convergence of the unprojected symmetric lattice operator is
neither used nor claimed; it is obstructed by high-momentum doublers, exactly as
explained in \cite{PelissierLorentzian}.  Strong-resolvent convergence remains
the global statement, while the fixed-momentum-band norm estimate is the correct
input for the low-energy semiclassical window.  The transfer theorem closes the microscopic origin of the principal massive
symbol and its isolated band.  It does not by itself compute the complete
order-\(\eps\) Lindblad/ordering symbol.  Nevertheless,
\Cref{subsec:unconditional-microscopic-closure} proves that on the admissible
channel class this symbol has the canonical geometric/gradient/irreversible
decomposition and that its scalar part is subleading in the observable used for
Ricci reconstruction.  Thus no additional zero-trace bridge remains for the
local renewal-to-Einstein theorem, although a full microscopic identification
of every subprincipal matrix entry remains model-dependent.
\end{remark}

Having established the positive transfer from the foundational Dirac limit to
the dynamical band, we next explain why the signed/Krein passage is not an
optional reformulation.  The Einstein argument below requires smooth null
generators, local causal diamonds, and a Lorentzian world function.  The next
two elementary observations show that neither of the positive-sector limits
realized by the renewal construction supplies this structure.

\begin{lemma}[Ellipticity of the positive renewal endpoint]
\label{lem:cp-elliptic}
For the positive continuum endpoint of \cite{PelissierLorentzian}, the principal
cometric has the frame-square form
$g_x^+(\theta,\theta)=\sum_a\langle e_a(x),\theta\rangle^2\ge0$.  If the local
frame spans, its real characteristic variety is empty away from the zero
covector; thus this positive endpoint is elliptic and carries no nonzero real
Lorentzian null directions.
\end{lemma}

\begin{proof}
The frame-square formula is part of the imported endpoint; it vanishes only if
every $\langle e_a,\theta\rangle$ vanishes, which for a spanning frame forces
$\theta=0$.
\end{proof}

\begin{lemma}[A polyhedral predictive cone is not exactly a Lorentz quadric]
\label{lem:poly-not-quad}
Let $\mathcal C\subset\R^{d+1}$ be a polyhedral cone with a genuine facet or edge
away from its apex.  Then $\mathcal C$ cannot equal the null set
$\{v:Q(v,v)=0\}$ of a nondegenerate Lorentzian quadratic form.
\end{lemma}

\begin{proof}
Away from the apex a nondegenerate Lorentzian null cone is a smooth real
algebraic hypersurface, while a genuine polyhedral cone is a finite union of
planar facets meeting along nonsmooth edges; equality would force a smooth
irreducible quadric to contain an open planar facet and its adjacent edge,
impossible for a nondegenerate form.  This excludes exact equality, not
smoothing or a Lorentz--Finsler interpretation.
\end{proof}

\begin{theorem}[Positive-sector obstruction for the realized renewal limits]
\label{thm:positive-obstruction}
In either positive-sector limit realized by the construction the Einstein
assembly is unavailable: \textup{(a)} the elliptic endpoint of
\Cref{lem:cp-elliptic} has no nonzero real null family on which the
Raychaudhuri, local-wedge, and null-reconstruction arguments can act; \textup{(b)}
the unrounded deterministic predictive limit has a polyhedral cone, which by
\Cref{lem:poly-not-quad} is not exactly a Lorentzian null cone and does not
supply the world function and van Vleck response used below.  Hence the
signed/Krein endpoint is necessary for the specific metric-Lorentzian assembly
proved here.
\end{theorem}

\begin{proof}
Part (a) is the absence of nonzero null covectors at the elliptic endpoint; part
(b) is that the Ricci-diamond and Clausius steps use the world function,
Levi--Civita connection, and smooth quadratic null generators of a Lorentzian
metric.  The theorem excludes these steps in the two realized positive limits, not
in every conceivable continuum completion.
\end{proof}

The local symbol of the frozen generator is therefore
$M_x(\theta)=c\langle e_a(x),\theta\rangle\alpha^a+\lambda(x)(\beta-1)$, with
branches $\sigma_\pm(x,\theta)=\pm\sqrt{\lambda(x)^2+c^2g_x(\theta,\theta)}-\lambda(x)$.
The obstruction theorem tells us why this signed symbol is necessary; the next
subsection explains what dynamical information is contained in its isolated
lower band.  In particular, the phase of the corresponding kernel will provide
the macroscopic action, while its leading transport amplitude will later encode
curvature.

\subsection{Matrix WKB propagation and the proper-time action}

We now pass from operator convergence to semiclassical propagation.  Two points
are essential.  First, in the physically selected case $d=3$ the massive Dirac
band has constant multiplicity two and must not be replaced by a fictitious
one-dimensional eigenspace.  Second, its Berry connection has both spacetime and
momentum-space components.  The correct object is therefore a matrix-valued band
parametrix.  The scalar observables used below are determinant or
Hilbert--Schmidt responses, for which the unitary Berry holonomy cancels.

\begin{hypothesis}[Slow homogeneity and gap]
\label{hyp:SH}
The fields $e_a,\lambda$ are $C^3$ with bounded derivatives on the patch,
$g_x=\sum_a e_a\otimes e_a$ is uniformly nondegenerate, and
$\lambda(x)\ge\lambda_{\min}>0$.  The renewal semigroup is considered in
semiclassical scaling
$\partial_T\psi_\eps=\eps^{-1}H_\eps\psi_\eps$ with
$H_\eps=\eps c\,e_a(x)\cdot\nabla\,\alpha^a+\lambda(x)(\beta-1)$.
\end{hypothesis}

\begin{hypothesis}[Constant-multiplicity renewal band]
\label{hyp:wkb}
In addition to \Cref{hyp:SH}, on a relatively compact patch $X$ with
$\mathcal H=L^2(X,d\mu_g;\C^m)$ assume the coarse-grained renewal semigroup is
generated by a sectorial family $\mathcal L_\eps$ with:
\begin{enumerate}[label=\textup{(W\arabic*)},leftmargin=2.9em]
\item \emph{semiclassical realization:} after conjugation by the coarse-graining
map, $\mathcal L_\eps$ is a first-order semiclassical $\Psi$DO with symbol
$\ell(x,\theta;\eps)=M_x(\theta)+\eps\ell_1(x,\theta)+O(\eps^2)$ in $S^1_{1,0}$,
uniformly with three derivatives, and the microscopic transfer semigroups
converge to $e^{T\mathcal L_\eps/\eps}$ in graph norm on compact time intervals;
\item \emph{gap-isolated band of constant rank:} the tilted symbol has a real
band eigenvalue
$-h(x,\theta)=\lambda(x)-\sqrt{\lambda(x)^2+c^2g_x(\theta,\theta)}$,
separated by a gap $\delta>0$, with a $C^3$ Riesz projector $\Pi_+$ of constant
rank $r$;
\item \emph{matrix adiabatic reduction:} bounded intertwiners reduce
$\mathcal L_\eps$ on the band to
$-\operatorname{Op}_\eps(h)\mathbf 1_r+
 \eps\operatorname{Op}_\eps(b)+\eps^2R_\eps$,
$\|R_\eps\|_{H^1\to L^2}\le C$, with complementary band
$O(e^{-\delta T/\eps})$;
\item \emph{nonfocal minimizer:} endpoints are joined by a unique nondegenerate
minimizer of the Legendre action of $h$, with no conjugate point before $T$;
\item \emph{controlled reset transport:} writing the projected transport as
$\nabla^{\rm band}=\nabla^{\LC}+\nabla^{\mom}+\mathcal A_{\ren}$, the
first-order reset drift is either calibrated explicitly or removed by a
stationary-state ground-state transform.  Any surviving real scalar potential
$V$ is retained in the flat-normalized residual rather than identified with the
van Vleck transport.
\end{enumerate}
\Cref{thm:foundational-transfer} supplies the microscopic origin of the
principal massive symbol and the gap.  On the admissible channel class of
\Cref{def:admissible-renewal-channel}, the required projected symbol convergence
and its curvature-level consequence are proved in
\Cref{thm:unconditional-bare-remainder,thm:unconditional-wkb-ricci}.
Condition \textup{(W5)} therefore records the transport bookkeeping needed by
the abstract parametrix; it is not an additional zero-trace hypothesis for the
local Ricci reconstruction.
\end{hypothesis}

\begin{lemma}[Band symbol and Legendre dual]
\label{lem:band-legendre}
Under \Cref{hyp:SH,hyp:wkb},
$h_x(\theta)=\sqrt{\lambda(x)^2+c^2g_x(\theta,\theta)}-\lambda(x)$ is smooth and
strictly convex in $\theta$, with Legendre transform
\[
 L_x(v)=\lambda(x)\left(1-\sqrt{1-g_x(v,v)/c^2}\right)
\]
for $g_x(v,v)<c^2$ and $+\infty$ outside the spectral velocity cone.
\end{lemma}

\begin{proof}
Differentiation gives
$\partial_\theta h=c^2g^{-1}\theta/(\lambda^2+c^2g^{-1}(\theta,\theta))^{1/2}$.
The Hessian is positive definite on the open cone because
$\lambda\ge\lambda_{\min}>0$.  Solving this relation for $\theta$ and substituting
in $\langle\theta,v\rangle-h$ yields the displayed formula; no finite dual point
exists when $g(v,v)\ge c^2$.
\end{proof}

\begin{theorem}[Constant-multiplicity band parametrix]
\label{thm:matrix-wkb}
Assume \Cref{hyp:wkb}.  On every compact nonfocal set and
$T\in[T_0,T_1]\Subset(0,\infty)$, the band kernel has the matrix WKB form
\[
 K^+_\eps(p,q;T)=e^{-S_T(p,q)/\eps}\,
 J_T(p,q)^{1/2}\,U_\gamma(p,q)
 \bigl(\mathbf 1_r+O_{C^0}(\eps)\bigr)
 +O(e^{-\delta T/(2\eps)}),
\]
where $S_T=\inf_\gamma\int_0^T L_\gamma(\dot\gamma)dt$ solves
$\partial_TS+h(x,\nabla_xS)=0$, $J_T^{1/2}$ is the scalar mixed-endpoint
Jacobian of the Hamilton flow, and
$U_\gamma=\mathcal P\exp(-\int_\gamma\mathcal A)$ is the parallel transport of
the projected band connection.
\end{theorem}

\begin{proof}
The constant-rank Riesz projector admits a superadiabatic refinement with
$[\mathcal L_\eps,\Pi_{+,\eps}]=O(\eps^\infty)$, and the spectral gap suppresses
the complementary band exponentially.  Conjugation by a local orthonormal band
frame gives a matrix operator with scalar principal symbol $-h\mathbf 1_r$ and
matrix subprincipal symbol $b$.  The multicomponent WKB construction then yields
the scalar Hamilton--Jacobi equation, the mixed-endpoint Jacobian, and the
non-abelian transport equation
$\dot U_\gamma=-\mathcal A(\dot x,\dot\theta)U_\gamma$.
This is the standard constant-multiplicity band parametrix
\cite{LittlejohnFlynn,WilczekZee,EmmrichWeinstein,Teufel}; the gap and nonfocal
Hessian bounds give the stated uniform remainders.
\end{proof}

\begin{lemma}[Anti-Hermiticity of geometric Berry transport]
\label{lem:berry-antiherm}
For any $C^1$ orthonormal frame $u_1,\ldots,u_r$ of the band, the Berry matrix
$\mathcal A^{B}_{ab}=\langle u_a,du_b\rangle$ is anti-Hermitian.  Hence its
holonomy lies in $U(r)$.
\end{lemma}

\begin{proof}
Differentiating $\langle u_a,u_b\rangle=\delta_{ab}$ gives
$\langle du_a,u_b\rangle+\langle u_a,du_b\rangle=0$, which is precisely
$(\mathcal A^B)^*=-\mathcal A^B$ \cite{Berry,Simon,WilczekZee}.
\end{proof}

\begin{theorem}[Band-transport decomposition]
\label{thm:transport-decomp}
For the massive Dirac projector, the projected transport decomposes as
\[
 \nabla^{\rm band}=\nabla^{\LC}+\nabla^{\mom}+\mathcal A_{\ren}.
\]
Here $\nabla^{\LC}$ is the spacetime-frame contribution, $\nabla^{\mom}$ is the
momentum-space Berry connection, and $\mathcal A_{\ren}$ is the projected
reset/Lindblad subprincipal term.  The first two terms are anti-Hermitian.  Thus
only $\kappa_{\ren}=r^{-1}\Re\Tr\mathcal A_{\ren}$ can change a scalar modulus
response.
\end{theorem}

\begin{proof}
Split the exterior differential in phase space as $d=d_x+d_\theta$ in
$\Pi_+d\Pi_+$.  The frame-dependent $d_x$ contribution is the usual spinorial
connection induced by the moving orthonormal frame, whereas the $d_\theta$
contribution is the generally nonzero momentum-space Berry connection
\cite{LittlejohnFlynn,EmmrichWeinstein}.  Both arise from orthogonal projection of
an orthonormal band frame and are anti-Hermitian by
\Cref{lem:berry-antiherm}.  The remaining projected order-$\eps$ symbol is
$\mathcal A_{\ren}$, and taking $r^{-1}\Re\Tr$ annihilates the two geometric
Berry terms.
\end{proof}

\begin{proposition}[Unitary-invariant scalar band response]
\label{prop:invariant-response}
Define
\[
 \Phi_{\det}:=\frac1r\log|\det K^+_\eps|,
 \qquad
 \Phi_{\rm HS}:=\frac12\log\left(\frac1r\|K^+_\eps\|_{\rm HS}^2\right).
\]
Then
\[
 \Phi_\bullet=-\frac{S_T}{\eps}+\frac12\log J_T
 -\int_\gamma(\kappa_{\ren}+V)+O(\eps),
 \qquad \bullet\in\{\det,{\rm HS}\},
\]
whenever the Hermitian part of $\mathcal A_{\ren}$ is scalar on the band; for the
determinant formula no scalarity assumption is needed.  In particular, all
Levi--Civita and momentum-space Berry holonomies cancel from the modulus.
\end{proposition}

\begin{proof}
Taking determinants in \Cref{thm:matrix-wkb} gives
$|\det K^+_\eps|=e^{-rS_T/\eps}J_T^{r/2}|\det U_\gamma|(1+O(\eps))$.
Jacobi's formula for the path-ordered exponential yields
$\log|\det U_\gamma|=-\int_\gamma\Re\Tr\mathcal A$; by
\Cref{thm:transport-decomp} this is
$-r\int_\gamma(\kappa_{\ren}+V)$ after separating the first-order drift
from the real scalar potential.  For the Hilbert--Schmidt functional, unitary
geometric holonomy preserves the Hilbert--Schmidt norm, and a scalar Hermitian
reset drift factors out in the same way.
\end{proof}

\begin{theorem}[Renewal WKB--van Vleck expansion]
\label{thm:wkb}
Under \Cref{hyp:SH,hyp:wkb}, removal of the first-order reset drift identifies
$J_T$ with the van Vleck--Morette determinant of the characteristic geometry,
denoted $\Delta_{\wideg}$, and, after flat normalization of the scalar potential,
\[
 \Phi_\bullet(p,q;T)
 =-\eps^{-1}S_T(p,q)+\tfrac12\log\Delta_{\wideg}(p,q)+O(\eps)
\]
uniformly on compact nonfocal sets.  Without drift removal and flat normalization, the additional scalar term is
$-\int_\gamma(\kappa_{\ren}+V)$.  For constant $\lambda$,
$\wideg=g$; for slowly varying $\lambda$, $\wideg$ denotes the Hamiltonian
characteristic geometry and differs from $g$ by the controlled rate-gradient
terms of \Cref{thm:adiabatic-rate} below.
\end{theorem}

\begin{proof}
The scalar mixed-endpoint Jacobian in the Hamiltonian parametrix is the
van Vleck determinant of the associated nonfocal characteristic flow.  The
invariant-response formula removes all anti-Hermitian transport.  After the drift has been removed and the constant part of $V$ subtracted by flat
normalization, no leading scalar weight remains; the nonconstant part of $V$ is
recorded in the deficiency residual.
For constant rate the Euler--Lagrange trajectories are metric geodesics by
\Cref{thm:constant-rate-geodesics}, so the characteristic determinant is
$\Delta_g$.
\end{proof}

\begin{corollary}[Derived local large deviations and Gaussian response]
\label{hyp:LD}
For every sufficiently small nonfocal diamond, the invariant band response
satisfies the local large-deviation principle with rate $S_T$ and Gaussian
prefactor $\Delta_{\wideg}^{1/2}$, with
\[
 E_{\rm WKB}(\eps,\tau)\le C_1\eps+C_2e^{-\delta\tau/(2\eps)}
 +E_{\rm cg}(\eps,\tau)+C_3\tau\|\nabla\log\lambda\|
 +C_4\tau^{-1}E_V(D_\tau).
\]
\end{corollary}

\medskip
\noindent\textbf{From the WKB phase to proper time.}
Let
\[
 S_T(x,y)=\inf_\gamma\int_0^T
 \lambda(\gamma)\left(1-\sqrt{1-g_\gamma(\dot\gamma,\dot\gamma)/c^2}\right)dt.
\]
The rate field must be treated explicitly: exact metric geodesics occur at constant
rate, while a slowly varying rate produces a controlled force.

\begin{lemma}[Conserved energy at constant rate]
\label{lem:constant-rate-energy}
If $\lambda$ is constant, the autonomous Lagrangian energy
$E=v^i\partial_{v^i}L-L=\lambda\phi(g(v,v))$ is conserved, where
$\phi(s)=(1-s/c^2)^{-1/2}-1$ is strictly increasing on $[0,c^2)$.
\end{lemma}

\begin{proof}
The Beltrami identity gives conservation.  Direct differentiation of
$f(s)=1-\sqrt{1-s/c^2}$ yields
$2sf'(s)-f(s)=(1-s/c^2)^{-1/2}-1$, whose derivative is positive.
\end{proof}

\begin{theorem}[Constant-rate geometric exactness]
\label{thm:constant-rate-geodesics}
If $\nabla\lambda=0$, every minimizer of $S_T$ has constant $g$-speed and is a
$g$-geodesic.  Together with the unit-lapse time direction it is a timelike
geodesic of $\mathbf g=-dt^2+g$.
\end{theorem}

\begin{proof}
By \Cref{lem:constant-rate-energy}, strict monotonicity of $\phi$ makes
$g(\dot\gamma,\dot\gamma)$ constant.  Hence $f'(g(\dot\gamma,\dot\gamma))$ is a
nonzero constant, and the Euler--Lagrange equation reduces after division to
\[
 \frac d{dt}(g_{ij}\dot\gamma^j)
 =\frac12\partial_i g_{k\ell}\dot\gamma^k\dot\gamma^\ell,
\]
the lowered geodesic equation.  The time component is affine
\cite{Arnold}.
\end{proof}

\begin{theorem}[Adiabatic rate-gradient control]
\label{thm:adiabatic-rate}
Let $\gamma_{\ren}$ be a renewal minimizer and $\gamma_0$ the characteristic of
the Hamiltonian with $\lambda$ frozen at the diamond centre.  Below the first
conjugate point and with velocities bounded away from $c$,
\[
 \|\gamma_{\ren}-\gamma_0\|_{C^1([0,\tau])}
 \le C\tau^2\|\nabla\log\lambda\|_{C^0(D_\tau)}.
\]
The induced phase and logarithmic Jacobian errors are respectively
$O(\tau^2\|\nabla\log\lambda\|)$ and
$O(\tau\|\nabla\log\lambda\|)$.
\end{theorem}

\begin{proof}
The Hamiltonian vector fields of $h$ and its frozen-rate version differ by
$O(\nabla\log\lambda)$ on the relevant compact phase-space set.  The difference
of the fixed-endpoint characteristics satisfies the linearized Hamilton--Jacobi
boundary problem with this source.  Below the first conjugate point its Jacobi
Green operator has norm $O(\tau^2)$, giving the $C^1$ estimate by variation of
constants and Gronwall.  Smooth dependence of the action and mixed-endpoint
Jacobian on the flow gives the two stated errors \cite{Hartman}.
\end{proof}

\begin{corollary}[Quadratic optical metric]
\label{cor:optical-metric}
In the near-coincidence regime $g(v,v)\ll c^2$,
\[
 L_x(v)=\frac12\widetilde g_x(v,v)+O(g(v,v)^2/c^4),
 \qquad \widetilde g:=\frac{\lambda}{c^2}g.
\]
Thus the leading spatial characteristics are geodesics of the optical metric
$\widetilde g$.  If the renewal calibration extends this rescaling conformally
to a spacetime metric $\widehat{\mathbf g}=e^{2\omega}\mathbf g$ with
$\omega=\tfrac12\log(\lambda/\lambda_0)$, then
\[
 \Ric[\widehat{\mathbf g}]_{\mu\nu}=\Ric[\mathbf g]_{\mu\nu}
 -(n-2)(\nabla_\mu\nabla_\nu\omega-\nabla_\mu\omega\nabla_\nu\omega)
 -(\Box\omega+(n-2)|\nabla\omega|^2)\mathbf g_{\mu\nu}.
\]
Hence the conversion from optical to background Ricci is explicit and belongs to
the rate-gradient/deficiency sector.
\end{corollary}

\begin{proof}
Taylor expansion gives
$1-\sqrt{1-s/c^2}=s/(2c^2)+s^2/(8c^4)+O(s^3)$, proving the optical metric
statement.  The final identity is the standard conformal transformation law for
the Ricci tensor \cite{WaldBook,Besse}.
\end{proof}

\begin{theorem}[Renewal proper-time-deficit action]
\label{thm:proper-time}
For endpoints joined by a unique nondegenerate minimizing characteristic,
\[
 \eps\,\Phi_\bullet(x,y;T)\longrightarrow -S_T(x,y).
\]
At constant rate the minimizers are timelike geodesics of $\mathbf g$ by
\Cref{thm:constant-rate-geodesics}; at slowly varying rate they obey the exact
Hamilton equations of $h$ and satisfy the estimate of
\Cref{thm:adiabatic-rate}.  No claim of exact $g$-geodesicity is made when
$\nabla\lambda\ne0$.
\end{theorem}

\begin{proof}
The phase statement follows from \Cref{thm:matrix-wkb} and
\Cref{prop:invariant-response}; the characteristic statements are
\Cref{thm:constant-rate-geodesics,thm:adiabatic-rate}.
\end{proof}

\begin{theorem}[Causal propagation]
\label{thm:causal-propagation}
The band kernel is exponentially suppressed whenever the endpoint lies outside
the spectral velocity cone, while inside it the Hamilton--Jacobi reachable set is
generated by the characteristics of $h$.
\end{theorem}

\begin{proof}
The Legendre transform is infinite outside the cone and finite inside it.  The
Lax--Oleinik representation therefore gives exponential suppression outside and
the characteristic reachable set inside.
\end{proof}


\subsection{An explicit massive-Dirac renewal model}
\label{subsec:nonempty-model}

The preceding parametrix is formulated for an abstract constant-rank renewal
band. To show that its assumptions are mutually compatible, and to separate
the structural hypotheses from the remaining microscopic estimate, we now
exhibit a concrete class of slowly varying massive-Dirac renewal generators.
The discussion proceeds in three stages: first the spectral gap and adiabatic
reduction, then an exact velocity-jump realization of the principal band, and
finally the treatment of the projected subprincipal drift.

\begin{construction}[Slowly varying massive-Dirac renewal model]
\label{constr:massive-dirac-model}
Let $(X,g)$ be a relatively compact patch with $C^3$ orthonormal coframe $e^a$
and $C^3$ rate field $\lambda\ge\lambda_{\min}>0$.  Let $\C^m$ carry an
irreducible Hermitian Clifford representation and let $\mathcal E$ be a primitive
unital completely positive reset channel with faithful invariant state.  Set
\[
 M(x,\theta)=c\,e^a_i(x)\theta^i\alpha^a+\lambda(x)(\beta-\mathbf 1),
 \qquad
 \mathcal L_\eps=\operatorname{Op}^{\rm W}_\eps(M)
 +\eps\operatorname{Op}^{\rm W}_\eps(\ell_1).
\]
\end{construction}

\begin{proposition}[Explicit gap and constant-rank projector]
\label{prop:model-W2}
Writing $N=c\,e^a_i\theta^i\alpha^a+\lambda\beta$, one has
$N^2=(\lambda^2+c^2g_x(\theta,\theta))\mathbf 1$.  Hence
\[
 \sigma_\pm=\pm|N|-\lambda,
 \qquad
 \Pi_\pm=\tfrac12(\mathbf 1\pm N/|N|),
\]
with gap at least $2\lambda_{\min}$ and constant rank $r=m/2$.  In $d=3$,
$r=2$; the band is not simple, but its principal eigenvalue is scalar on the
band.
\end{proposition}

\begin{proof}
The Clifford anticommutation relations cancel the cross terms in $N^2$.  Since
$|N|\ge\lambda_{\min}$, functional calculus gives $C^3$ projectors with bounded
derivatives and the claimed uniform gap.
\end{proof}

The explicit projector provides the spectral data required by the band
construction. The next proposition promotes those pointwise data to an
operator-level reduction and isolates the exponentially suppressed
complementary sector.

\begin{proposition}[Matrix reduction and complementary decay]
\label{prop:model-W3}
The uniform gap and $C^3$ projector give bounded superadiabatic intertwiners for
which the on-band generator is
$-\operatorname{Op}_\eps(h)\mathbf 1_r+
 \eps\operatorname{Op}_\eps(b)+\eps^2R_\eps$ and the complementary band is
$O(e^{-2\lambda_{\min}T/\eps})$.
\end{proposition}

\begin{proof}
This is the constant-multiplicity adiabatic reduction of
\cite{Nenciu,Teufel,EmmrichWeinstein}; the scalarity of the principal band
eigenvalue leaves only matrix transport at subprincipal order.
\end{proof}

The adiabatic reduction controls the internal band dynamics, but the WKB
kernel also requires a nondegenerate classical trajectory. Strict convexity of
the same Hamiltonian supplies the needed local noncaustic window.

\begin{lemma}[Noncaustic window]
\label{lem:model-W4}
There exists $T_*>0$ such that for $T<T_*$ every admissible pair of endpoints is
joined by a unique nondegenerate minimizing characteristic.
\end{lemma}

\begin{proof}
Strict convexity of $h$ and continuous dependence of the linearized Hamilton
flow give a positive lower bound on the first conjugate time on every relatively
compact phase-space set.
\end{proof}

These ingredients now close the abstract hypotheses at principal order.
Collecting them identifies precisely what has been verified and what remains to
be checked for a bare microscopic reset channel.

\begin{theorem}[Explicit model verification and admissible-class completion]
\label{thm:model-verifies}
For \Cref{constr:massive-dirac-model}, conditions \textup{(W2)--(W4)} hold and
\textup{(W1)} holds by construction.  The geometric Berry transport is unitary
and therefore harmless for the determinant and Hilbert--Schmidt modulus
responses.  On the reversible isotropic subclass,
\[
  \kren=\frac1r\Re\Tr\mathcal A_{\ren}=0
\]
after the stationary-state ground-state transform.  On a general admissible
channel, \(\kren\) need not vanish; nevertheless its entire contribution to the
flat-normalized, reflection-symmetric Ricci tomography is the subleading term
\(O(\tau\|\kren+V\|_{C^2})\).  Hence vanishing of the trace-Hermitian
subprincipal term is sufficient for the simplest worked model but is not
required for the local Einstein assembly.
\end{theorem}

\begin{proof}
The gap, constant-rank projector, adiabatic reduction, and noncaustic window are
\Cref{prop:model-W2,prop:model-W3,lem:model-W4}.  Anti-Hermiticity of the
Levi--Civita and momentum-space Berry connections follows from
\Cref{lem:berry-antiherm,thm:transport-decomp}, so they disappear from the
unitary-invariant scalar response of \Cref{prop:invariant-response}.  For the
reversible isotropic subclass, \Cref{thm:drift-cancellation} removes the
gradient drift and leaves no irreversible component, hence \(\kren=0\).
For arbitrary admissible channels, the canonical decomposition and
flat-normalized estimate are
\Cref{lem:canonical-drift-split,lem:line-integral,thm:unconditional-ricci};
they place the surviving scalar in
\(O(\tau\|\kren+V\|_{C^2})\).  This proves the final assertion.
\end{proof}



The preceding model verifies consistency of the semiclassical assumptions,
but it does not yet explain why the relativistic square-root band should arise
from renewal statistics. The following elementary velocity-jump process gives
that band exactly, rather than only as an effective ansatz.

\begin{construction}[Discrete velocity-jump renewal process]
\label{constr:velocity-jump}
At each renewal epoch the predictive state records one of two ballistic
velocities.  In one spatial dimension the tilted continuous-time generator is
\[
 L_\theta=
 \begin{pmatrix}
 c\theta-\lambda&\lambda\\
 \lambda&-c\theta-\lambda
 \end{pmatrix}.
\]
In dimension $d$, the velocity observable is replaced by
$c\alpha^i\theta_i$, where the Hermitian matrices $\alpha^i$ satisfy
$\{\alpha^i,\alpha^j\}=2g^{ij}\mathbf 1$.
\end{construction}

\begin{theorem}[Exact renewal origin of the massive band]
\label{thm:exact-renewal-band}
The Perron eigenvalue of the tilted generator in
\Cref{constr:velocity-jump} is
\[
 h(x,\theta)=\sqrt{\lambda(x)^2+c^2g^{ij}(x)\theta_i\theta_j}-\lambda(x).
\]
Thus the Hamiltonian used in the renewal WKB construction is not merely an
asymptotic ansatz: it is the exact log-growth rate of a discrete velocity-jump
renewal model.  In the small-momentum regime
$h=\frac{c^2}{2\lambda}g^{ij}\theta_i\theta_j+O(|\theta|^4)$, while for large
momentum $h\sim c|\theta|_g$, recovering respectively the diffusive and causal
limits.
\end{theorem}

\begin{proof}
The eigenvalues of the one-dimensional matrix are
$-\lambda\pm\sqrt{\lambda^2+c^2\theta^2}$; the dominant eigenvalue is the
displayed $h$.  In the Clifford realization,
$(c\alpha^i\theta_i)^2=c^2g^{ij}\theta_i\theta_j\mathbf1$, so the same spectral
calculation gives the $d$-dimensional formula.  The two asymptotic expansions are
immediate.  The underlying telegraph process is classical \cite{GoldsteinKac,Kac}.
Its relativistic and spinorial use is also established: Gaveau, Jacobson, Kac,
and Schulman derived a Dirac propagator from Poisson-distributed velocity
reversals, while Ord and McKeon--Ord developed correlated-walk realizations in
$1+1$ dimensions \cite{GaveauJacobsonKacSchulman,Ord1983,McKeonOrd}.  The claim
here is therefore not the telegraph--Dirac correspondence itself, but its
placement inside the predictive renewal quotient, the identification of its
rate with the primitive reset scale, and its transfer to the curved isolated
band used in the subsequent WKB analysis.
\end{proof}

The exact band fixes the principal symbol. We next turn to the subprincipal
amplitude, where an uncontrolled Hermitian drift would contaminate the van
Vleck response. Reversibility and isotropy provide a natural class in which
that drift can be identified and removed.

\begin{hypothesis}[Reversible isotropic reset class]
\label{hyp:reversible-isotropic}
The primitive reset channel is quantum detailed-balanced with respect to its
faithful invariant state $\rstar$, the spatial jump kernel is isotropic with
vanishing odd cumulants, the discrete derivative is central, and the continuum
generator is obtained by Weyl/midpoint quantization with $h_R^2=o(\eps)$.
\end{hypothesis}

The band-reduced subprincipal symbol has the standard space-adiabatic form
\begin{equation}
 b_{\rm eff}=\Pi_+\ell_1\Pi_+
 +\frac{i}{2}\Pi_+\{M,\Pi_+\}\Pi_+ + b_{\rm quant}.
 \label{eq:projected-subprincipal}
\end{equation}

\begin{lemma}[Geometric and quantization terms]
\label{lem:projected-geo-quant}
The normalized real trace of the geometric no-name term in
\eqref{eq:projected-subprincipal} vanishes, while the remaining geometric
transport is the Liouville half-density term generating
$\Delta_{\wideg}^{1/2}$.  For Weyl/midpoint quantization and central
differences, the order-$\eps$ quantization drift also vanishes when
$h_R^2=o(\eps)$.
\end{lemma}

\begin{proof}
The no-name term is the projected Berry-curvature contribution and has purely
imaginary trace \cite{LittlejohnFlynn,EmmrichWeinstein,PST}.  The scalar
Liouville term is $\tfrac12\dive H_h$ and integrates to the mixed-endpoint
Jacobian.  Weyl quantization introduces no additional kinetic subprincipal drift,
and central differences have an even $O(h_R^2)$ consistency error
\cite{DimassiSjostrand,HormanderIII}.
\end{proof}

\begin{lemma}[Reversible reset drift is a gradient]
\label{lem:gradient-reset-drift}
For the class of \Cref{hyp:reversible-isotropic}, the scalar part of the
band-projected first moment has the form
\[
 v=D\nabla\log\rstar^{1/2},
\]
where $D$ is the band diffusion tensor.  Isotropy and vanishing odd cumulants
exclude any additional nongradient first moment.
\end{lemma}

\begin{proof}
Quantum detailed balance is reversibility in the symmetric KMS inner product and
reduces on the spatial marginal to the reversible-jump identity.  The first
moment of a reversible jump kernel is the gradient drift associated with its
stationary density \cite{KossakowskiFGV,Alicki1976,Cipriani,KipnisLandim}.
Parity removes the remaining odd component.
\end{proof}

The geometric and quantization contributions are therefore harmless, while
the remaining reset drift is a stationary gradient. The ground-state transform
then gives the precise amplitude normalization needed by the curvature
reconstruction.

\begin{theorem}[Projected drift removal and residual reset potential]
\label{thm:drift-cancellation}
Assume \Cref{hyp:reversible-isotropic} and the order-$\eps$ symbol expansion of
\Cref{thm:foundational-transfer}.  In the $\rstar$-adapted half-density gauge,
the ground-state transform removes the first-order reset drift and leaves only a
real scalar potential,
\[
 \rstar^{-1/2}(\Pi_+\ell_1^{\rm res}\Pi_+)\rstar^{1/2}=V,
 \qquad
 V=\frac14\langle\nabla\log\rstar,D\nabla\log\rstar\rangle
 +\frac12\dive(D\nabla\log\rstar^{1/2}).
\]
Consequently the van Vleck prefactor is unrescaled at the reconstruction order.
The spatially constant part of $V$ cancels in the flat-normalized partition
function, while its nonconstant part is retained as the deficiency-sector error
$E_V(D_\tau)$.
\end{theorem}

\begin{proof}
By \Cref{lem:gradient-reset-drift}, the projected reset generator is a reversible
diffusion with gradient drift.  The standard Doob ground-state transform by
$\rstar^{1/2}$ eliminates its first-order term and produces the displayed
potential \cite{Doob,Pinsky}.  \Cref{lem:projected-geo-quant} removes the
geometric and discretization drifts.  Flat normalization subtracts the constant
part of the integrated potential; the remainder has order
$O((\nabla\log\rstar)^2,\nabla^2\log\rstar)$ and therefore belongs to the
rate-gradient/deficiency sector.
\end{proof}

\begin{remark}[Irreversible resets]
Without detailed balance the reset drift need not be a gradient and the invariant
response may acquire a genuine factor $\exp(-\int_\gamma\kappa_{\ren})$.
The reversible theorem removes only the drift; it does not assert the false
pointwise identity that every Hermitian scalar subprincipal term vanishes.
\end{remark}


To make the preceding cancellation theorem fully concrete, we now give a
primitive internal--spatial channel for which the invariant state, drift, and
residual potential can all be written explicitly. This example will also
support the uniform discrete remainder estimate below.

\begin{construction}[Worked reversible internal--spatial reset channel]
\label{constr:worked-reset-channel}
On $\mathfrak h=\C^2\otimes L^2(\R^d)$ let
$\mathcal G=\mathcal G_{\rm int}\otimes\id+\id\otimes\mathcal G_{\rm sp}$,
where
\[
 L_-=\sqrt\gamma\,\sigma_-,\qquad
 L_+=\sqrt{\gamma e^{-\beta\omega}}\,\sigma_+,
 \qquad
 \mathcal G_{\rm sp}=D(\Delta-\nabla\phi\cdot\nabla).
\]
Here $D>0$ and $\phi$ is a confining $C^4$ potential with bounded derivatives on
the local patch.  The spatial discretization is central and the continuum
operator is Weyl/midpoint quantized.
\end{construction}

\begin{proposition}[Explicit stationary data of the worked channel]
\label{prop:worked-reset-data}
The channel of \Cref{constr:worked-reset-channel} is primitive, unital, and
quantum detailed-balanced, with
\[
 \rstar=\rho_{\rm int}\otimes\pi,\qquad
 \rho_{\rm int}=\frac1{1+e^{-\beta\omega}}
 \begin{pmatrix}1&0\\0&e^{-\beta\omega}\end{pmatrix},
 \qquad \pi(x)=Z_\phi^{-1}e^{-\phi(x)}.
\]
Its scalar spatial diffusion, reversible drift, and ground-state potential are
\[
 D_{ij}=D\delta_{ij},\qquad
 v=D\nabla\log\pi^{1/2}=-\frac D2\nabla\phi,
 \qquad
 V=\frac D4|\nabla\phi|^2-\frac D2\Delta\phi.
\]
For $\phi(x)=\frac\alpha2|x|^2$ one obtains
$v(x)=-\frac{D\alpha}{2}x$ and
$V(x)=\frac{D\alpha^2}{4}|x|^2-\frac{D\alpha d}{2}$.
\end{proposition}

\begin{proof}
The two internal jump rates satisfy the Gibbs detailed-balance ratio, and direct
substitution gives $\mathcal G_{\rm int,*}(\rho_{\rm int})=0$.  The spatial
operator is reversible in $L^2(\pi dx)$ and has invariant density
$\pi\propto e^{-\phi}$.  Its first moment is
$D\nabla\log\pi^{1/2}$, and the ground-state transform
$e^{\phi/2}\mathcal G_{\rm sp}e^{-\phi/2}=D\Delta-V$ gives the displayed
potential.  Confinement and the qubit relaxation gap give primitivity.
\end{proof}

\begin{lemma}[Moment-bounded symbol regularity]
\label{lem:worked-channel-regularity}
Suppose the reset kernel $w(x,d\xi)$ has uniformly bounded moments through order
three and uniformly bounded symbol derivatives through the order required by
Calder\'on--Vaillancourt.  Then
\[
 \sigma_\eps(C_R\mathcal L_{R,\eps}C_R^*)
 =M+\eps\ell_1+\eps^2\rho_\eps+O(h_R^2),
\]
where $\rho_\eps$ has uniformly bounded symbol seminorms.  The Gaussian spatial
reset of \Cref{constr:worked-reset-channel} satisfies these assumptions.
\end{lemma}

\begin{proof}
Taylor expansion of the displacement factor in the discrete Weyl symbol gives
$M+\eps\ell_1$ with a Lagrange remainder controlled by the third moment.  Central
differences contribute an even $O(h_R^2)$ consistency error.  Bounded derivatives
of the kernel and the internal coefficients give the stated uniform symbol
seminorms.
\end{proof}

The explicit stationary data determine the first two terms of the symbol.
The remaining issue is quantitative: one must control the difference between
the bare discrete generator and its semiclassical truncation in the operator
norm used by the transfer theorem.

\begin{theorem}[Bare-discrete subprincipal remainder for the worked class]
\label{thm:bare-discrete-remainder}
Under \Cref{lem:worked-channel-regularity},
\[
 \left\|C_R\mathcal L_{R,\eps}C_R^*
 -\operatorname{Op}^{\rm W}_\eps(M)
 -\eps\operatorname{Op}^{\rm W}_\eps(\ell_1)\right\|_{H^1_\eps\to L^2}
 \le C(\eps^2+h_R^2).
\]
In particular this holds for \Cref{constr:worked-reset-channel}.  Together with
\Cref{thm:drift-cancellation}, the microscopic WKB--Ricci chain is closed for
this moment-bounded reversible isotropic class.
\end{theorem}

\begin{proof}
By \Cref{lem:worked-channel-regularity} the difference is
$\operatorname{Op}^{\rm W}_\eps(\eps^2\rho_\eps)+\mathcal R_h$.
The semiclassical Calder\'on--Vaillancourt theorem gives
$\|\operatorname{Op}^{\rm W}_\eps(\eps^2\rho_\eps)\|_{H^1_\eps\to L^2}
\le C\eps^2$, while the central-difference consistency estimate gives
$\|\mathcal R_h\|_{H^1_\eps\to L^2}\le Ch_R^2$
\cite{DimassiSjostrand,Zworski}.  The uniform band gap keeps the Riesz projector
and its derivatives bounded, so projection does not alter the estimate.
\end{proof}

\subsection{Admissible channels and subprincipal closure}
\label{subsec:unconditional-microscopic-closure}

The reversible isotropic construction above is a transparent worked subclass, but it is not required for the curvature reconstruction.  The following results remove reversibility, detailed balance, and isotropy from the bare-microscopic closure.

\subsubsection{Admissible renewal channels}
\label{sec:class}
% =====================================================================

\begin{definition}[Admissible renewal channel]
\label{def:admissible-renewal-channel}
Let $(X,g)$ be a relatively compact patch with $C^3$ orthonormal coframe $e^a$
and $C^3$ rate field $\lambda\ge\lambda_{\min}>0$.  A coarse-grained renewal
generator $C_R\Lcal_{R,\eps}C_R^*$ is \emph{admissible} if its primitive unital
completely positive reset channel $\mathcal E_x$ has a faithful invariant state
$\rstar(x)>0$ and:
\begin{enumerate}[leftmargin=2.6em]
\item[\textup{(A1)}] \emph{(moments)} the spatial jump kernel $w(x,d\xi)$ has
moments $M_n^{a_1\cdots a_n}(x)=\int\xi^{a_1}\cdots\xi^{a_n}\,w(x,d\xi)$ that are
$C^2$ in $x$ and uniformly bounded for $n\le 3$, and the resulting Weyl symbol has
seminorms bounded to the order required by the (semiclassical)
Calder\'on--Vaillancourt theorem;
\item[\textup{(A2)}] \emph{(gap)} the tilted principal symbol
$M(x,\theta)=c\,e^a_i\theta^i\alpha^a+\lambda(\beta-\mathbf1)$ has the two
square-root branches with uniform gap $\sigma_+-\sigma_-\ge2\lambda_{\min}$ and
$C^3$ Riesz projector $\Pip$ of constant rank $r$;
\item[\textup{(A3)}] \emph{(consistency)} the spatial difference operator is
central, with even $O(h_R^2)$ consistency, under Weyl/midpoint quantization with
$h_R^2=o(\eps)$.
\end{enumerate}
We write $\CV$ for the class of admissible channels.  \emph{No reversibility,
detailed balance, or isotropy is assumed.}
\end{definition}

\begin{remark}
\label{rem:cv-natural}
$\CV$ is essentially the largest class on which the statement can even be posed:
without (A1) the truncated symbol $M+\eps\ell_1$ is not defined, without (A2) there
is no isolated band and no scalar Hamiltonian $h$, and without (A3) the discrete
operator does not converge to its Weyl quantization.  The worked Gaussian channel
of the previous analysis, every compound-Poisson velocity-jump channel, and every
finite-range non-reversible reset with bounded moments all lie in $\CV$.
\end{remark}

% =====================================================================
\subsubsection{Uniform symbol remainder}
\label{sec:partI}
% =====================================================================

\begin{lemma}[Moment-bounded symbol regularity, reversibility-free]
\label{lem:cv-symbol-reg}
For every admissible channel,
\[
 \sigma_\eps\bigl(C_R\Lcal_{R,\eps}C_R^*\bigr)
 =M+\eps\,\ell_1+\eps^2\rho_\eps+O(h_R^2),
\]
where $\ell_1$ is the order-$\eps$ Weyl/Kramers--Moyal symbol determined by the
first and second jump moments and the internal channel, and $\rho_\eps$ has
uniformly bounded symbol seminorms.
\end{lemma}

\begin{proof}
Taylor-expand the displacement factor $e^{i\xi\cdot\theta}$ of the discrete Weyl
symbol of the reset kernel to second order; the linear and quadratic moments
$M_1,M_2$ produce $M+\eps\ell_1$, and the Lagrange remainder is controlled by the
third moment $M_3$, bounded uniformly by (A1).  The internal completely positive
generator contributes a bounded matrix symbol with the same regularity.  Central
differences contribute an even $O(h_R^2)$ consistency error by (A3).  None of
these steps uses any symmetry of $w$ beyond the boundedness of $M_{\le 3}$.
\end{proof}

\begin{theorem}[Unconditional bare-discrete remainder]
\label{thm:unconditional-bare-remainder}
For every admissible channel,
\[
 \bigl\|C_R\Lcal_{R,\eps}C_R^*
 -\OpW(M)-\eps\,\OpW(\ell_1)\bigr\|_{H^1_\eps\to L^2}\le C(\eps^2+h_R^2),
\]
with $C$ uniform over any $\CV$-subfamily with common moment and gap bounds.  In
particular the first-order semiclassical realization \textup{(W1)} holds on every
fixed physical momentum band, with no reversibility or isotropy hypothesis.
\end{theorem}

\begin{proof}
By Lemma~\ref{lem:cv-symbol-reg} the operator difference equals
$\OpW(\eps^2\rho_\eps)+\mathcal R_h$.  The semiclassical Calder\'on--Vaillancourt
theorem \cite{CalderonVaillancourt,Zworski,DimassiSjostrand} gives
$\|\OpW(\eps^2\rho_\eps)\|_{H^1_\eps\to L^2}\le C\eps^2$ from the bounded symbol
seminorms of $\rho_\eps$; the central-difference consistency estimate gives
$\|\mathcal R_h\|_{H^1_\eps\to L^2}\le Ch_R^2$.  By (A2) the Riesz projector $\Pip$
and its derivatives are bounded uniformly, so band projection does not enlarge the
estimate.  Reversibility and isotropy are nowhere invoked.
\end{proof}

\begin{remark}
\label{rem:partI-scope}
Theorem~\ref{thm:unconditional-bare-remainder} already upgrades the previous bare-discrete
remainder from the Gaussian worked channel to all of $\CV$, because the original
proof used only moment bounds, Calder\'on--Vaillancourt, central differences, and
the gap---all of which are exactly (A1)--(A3).  What remains is to show that the
\emph{value} of the surviving scalar $\kren$ does not contaminate the curvature
reconstruction; this is resolved by the drift decomposition and reconstruction estimates below.
\end{remark}

% =====================================================================
\subsubsection{Canonical drift decomposition}
\label{sec:partII}
% =====================================================================

We first record the structure of the projected order-$\eps$ symbol for an
arbitrary admissible channel.  Write $\Pip\ell_1\Pip$ in the half-density gauge.

\begin{lemma}[Geometric / gradient / irreversible split]
\label{lem:canonical-drift-split}
For every admissible channel there is a canonical decomposition, orthogonal in
$L^2(\rstar)$ on the band,
\[
 \Pip\,\ell_1\,\Pip
 =\mathcal A_{\mathrm{geo}}
 +\mathcal A_{\mathrm{grad}}
 +\mathcal A_{\mathrm{irr}},
\]
in which:
\begin{enumerate}[leftmargin=2.6em]
\item $\mathcal A_{\mathrm{geo}}=\nabla^{\mathrm{LC}}+\nabla^{\mathrm{mom}}$ is the
spacetime-frame plus momentum-space Berry transport, which is anti-Hermitian and
therefore obeys $\Re\Tr(\Pip\mathcal A_{\mathrm{geo}}\Pip)=0$;
\item $\mathcal A_{\mathrm{grad}}$ is the gradient part of the band first moment,
$v_{\mathrm{grad}}=D\,\nabla\phi$ with $\rstar=e^{-\phi}$ \textup{(}$D$ the band
diffusion tensor\textup{)}, removed by the ground-state transform
$e^{\phi/2}(\cdot)e^{-\phi/2}$ and producing the real scalar potential
$V=\tfrac14\langle\nabla\phi,D\nabla\phi\rangle-\tfrac12\dive(D\nabla\phi)$;
\item $\mathcal A_{\mathrm{irr}}$ is the solenoidal \textup{(}irreversible\textup{)}
part of the first moment, $\dive(\rstar v_{\mathrm{irr}})=0$, producing the real
scalar
\[
 \kren=\tfrac12\,v_{\mathrm{irr}}\!\cdot\!\nabla\phi
 =\tfrac1r\Re\Tr\bigl(\Pip\,\ell_1^{\mathrm{res}}\,\Pip\bigr).
\]
\end{enumerate}
The decomposition reduces, when $v_{\mathrm{irr}}=0$, to the reversible case of the
previous analysis \textup{(}then $\kren=0$\textup{)}.
\end{lemma}

\begin{proof}
Item~1 is the previous band-transport decomposition: differentiating the
band frame orthonormality gives anti-Hermitian $\nabla^{\mathrm{LC}}$ and
$\nabla^{\mathrm{mom}}$, whose real trace vanishes
\cite{LittlejohnFlynn,EmmrichWeinstein,PST}.  For items~2--3, the band first moment
$v$ is the drift of a (generally non-reversible) diffusion with faithful invariant
density $\rstar$.  The $L^2(\rstar)$ Helmholtz/Hodge decomposition of stationary
drifts \cite{Pavliotis,LelievreStoltz,KipnisLandim} writes uniquely
$v=v_{\mathrm{grad}}+v_{\mathrm{irr}}$ with $v_{\mathrm{grad}}=D\nabla\phi$ a
gradient and $v_{\mathrm{irr}}$ $\rstar$-divergence-free,
$\dive(\rstar v_{\mathrm{irr}})=0$.  Doob's $h$-transform by $\rstar^{1/2}$
eliminates the gradient part and yields the displayed potential $V$
\cite{Doob,Pinsky}.  The $\rstar$-divergence-free condition gives
$\dive v_{\mathrm{irr}}=v_{\mathrm{irr}}\!\cdot\!\nabla\phi$, so the residual
real scalar carried by the half-density transport is
$\tfrac12\dive v_{\mathrm{irr}}=\tfrac12v_{\mathrm{irr}}\!\cdot\!\nabla\phi=\kren$.
Orthogonality of the three pieces in $L^2(\rstar)$ is the orthogonality of
co-exact, exact, and harmonic-flux components.
\end{proof}

\begin{remark}[Consistency with the three-state ring]
\label{rem:ring}
For the irreversible three-state ring the solenoidal drift is the cycle current,
$v_{\mathrm{irr}}\propto j_{\mathrm{st}}$, and Lemma~\ref{lem:canonical-drift-split} reproduces the
previous $\kappa_{\mathrm{irr}}=C_0\,j_{\mathrm{st}}$ with the linear-response
coefficient $C_0$ now identified intrinsically as
$\tfrac12\,\partial_{j}(v_{\mathrm{irr}}\!\cdot\!\nabla\phi)$.  Thus $\kren$ need
not vanish off the reversible locus; the reconstruction estimate below shows that it is nevertheless subleading.
\end{remark}

% =====================================================================
\subsubsection{Subleading scalar drift in Ricci reconstruction}
\label{sec:partIII}
% =====================================================================

The closure rests on two structural facts about the real scalar
$\Sigma:=\kren+V$ produced by Lemma~\ref{lem:canonical-drift-split}: it is curvature-free, and it
enters the two-profile tomography only through a flat-normalized,
reflection-averaged line integral.

\begin{lemma}[Curvature-freeness of the scalar drift]
\label{lem:curvature-free}
For every admissible channel, $\Sigma=\kren+V$ is a universal local function
\[
 \Sigma(x)=\mathfrak s\bigl(e^a(x),\partial e^a(x),\lambda(x),
 \partial\lambda(x);\,\{M_n(x)\}_{n\le3}\bigr)
\]
of the order-$\le1$ jet of the frame and rate and the low-order channel moments.
It contains no explicit curvature: in Riemann normal coordinates centred at $p$,
$\Sigma$ admits the expansion
\[
 \Sigma(x)=\Sigma_0\bigl(\lambda,\partial\lambda;\{M_n\}\bigr)(x)
 +\mathfrak c_1\,[\Riem(p)]\!\cdot\!(x-p)
 +O\bigl(|x-p|^2\bigr),
\]
in which $\Sigma_0$ is independent of curvature and the curvature-linear term
$\mathfrak c_1[\Riem(p)]\!\cdot\!(x-p)$ is odd under $x-p\mapsto-(x-p)$.
\end{lemma}

\begin{proof}
By Lemma~\ref{lem:cv-symbol-reg} the symbol $\ell_1$ is the order-$\eps$ Weyl symbol,
built algebraically from $M$, $\partial_x M$, and the first/second jump moments;
$\partial_x M$ involves $\partial e^a,\partial\lambda$ but no second derivatives of
the frame.  The Riesz projector $\Pip=\tfrac12(\mathbf1+N/|N|)$ is algebraic in
$M$, hence in $(e^a,\lambda)$ alone.  Explicit curvature can enter the generator
only at order $\eps^2$, through the bounded-geometry channel coupling
$\|\mathcal E_x-\mathcal E\|_{\mathrm{cb}}\le C_0\eps^2(|\Ric_\Sigma|+|K|^2)$ of the
previous analysis, or through $\partial^2 e^a$, which lives in the order-$\eps^2$ symbol
$\ell_2$.  Therefore $\Sigma=\tfrac1r\Re\Tr(\Pip\ell_1\Pip)$ depends on the metric
only through $(e^a,\partial e^a)$ and on $\lambda$ only through
$(\lambda,\partial\lambda)$.  In normal coordinates at $p$ one has $e^a(p)=\delta$,
$\partial e^a(p)=0$ and $\partial^2 e^a(p)=O(\Riem(p))$, whence
$\partial e^a(x)=O(\Riem(p)\,(x-p))$.  Substituting into $\mathfrak s$ and expanding
gives the stated split: the term linear in $\partial e^a$ is linear in
$\Riem(p)\,(x-p)$ and odd, the curvature-free remainder is
$\Sigma_0(\lambda,\partial\lambda;\{M_n\})$, and all explicit curvature is pushed to
$O(|x-p|^2)$ or to order $\eps^2$.
\end{proof}

\begin{lemma}[Flat-normalized, reflection-averaged line integral]
\label{lem:line-integral}
Let $\rho_\tau$ be a normalized diamond profile invariant under the reflection
$x-p\mapsto-(x-p)$ and the spatial rotations fixing $u$ \textup{(}as in the
two-profile family of the previous analysis\textup{)}.  Then the scalar-drift
contribution to the flat-normalized diamond response
$\mathfrak D=\log\mathcal Z_\eps(D_\tau)/\mathcal Z_\eps(D_\tau^{\mathrm{flat}})$
obeys
\[
 \Bigl|\;\Bigl\langle\!\int_\gamma\Sigma\Bigr\rangle_{\rho_\tau}
 -\Bigl\langle\!\int_\gamma\Sigma\Bigr\rangle^{\mathrm{flat}}_{\rho_\tau}\;\Bigr|
 \le C\,\tau^3\bigl(\|\Sigma\|_{C^2(D_\tau)}+\|\Riem\|_{C^0(D_\tau)}\bigr).
\]
\end{lemma}

\begin{proof}
Parametrize a minimizing characteristic $\gamma$ by proper time $t\in[0,T]$,
$T=O(\tau)$, $|\gamma(t)-p|\le C\tau$.  By Lemma~\ref{lem:curvature-free},
\[
 \int_\gamma\Sigma
 =\Sigma_0(p)\,T
 +\nabla\Sigma_0(p)\!\cdot\!\!\int_0^T(\gamma-p)\,dt
 +\mathfrak c_1[\Riem(p)]\!\cdot\!\!\int_0^T(\gamma-p)\,dt
 +\tfrac12\!\int_0^T\!\nabla^2\Sigma(\gamma-p)^2\,dt+\cdots.
\]
The constant term $\Sigma_0(p)\,T$ is curvature-free and identical for the curved
and flat diamonds in normal coordinates (same $\lambda,\partial\lambda$ data at
$p$), so it cancels in the flat-normalized difference.  The two terms linear in
$(\gamma-p)$ are odd under the diamond reflection; averaging against the symmetric
profile $\rho_\tau$ annihilates them, and the residual curved/flat discrepancy from
geodesic bending is $O(\tau\cdot\tau\cdot\tau/L)=O(\tau^3/L)$.  The quadratic term is
bounded by $\|\nabla^2\Sigma\|\,\tau^2\cdot T=O(\|\Sigma\|_{C^2}\tau^3)$.  Collecting
the surviving contributions gives the stated bound.
\end{proof}

\begin{theorem}[Unconditional two-profile Ricci reconstruction]
\label{thm:unconditional-ricci}
Let two flat-normalized calibrated diamond measures of an admissible channel
converge in second moments to positive rotationally symmetric, reflection-symmetric
profiles $\rho_1,\rho_2$ with invertible moment matrix $M=\binom{a_1\ b_1}{a_2\ b_2}$.
Then, with $\mathfrak D_i=\log\mathcal Z_{\eps,i}(D_\tau(p,u))$,
\[
 \frac1{\tau^2}\binom{\mathfrak D_1}{\mathfrak D_2}
 =M\binom{\widehat R_{\mu\nu}u^\mu u^\nu}{\widehat R}
 +O\!\Big(\tfrac{\tau}{L^3}\Big)+O\!\Big(\tfrac{\eps}{\tau^2}\Big)
 +O\bigl(\tau\|\nabla\log\lambda\|\bigr)
 +O\bigl(\tau\,\|\Sigma\|_{C^2}\bigr),
\]
and inversion reconstructs $\widehat R_{\mu\nu}u^\mu u^\nu$ and $\widehat R$;
varying $u$ over an open timelike set reconstructs the full characteristic Ricci
tensor.  No reversibility, detailed-balance, or isotropy hypothesis is used: the
entire dependence on the projected reset drift $\Sigma=\kren+V$ is the
\emph{subleading} term $O(\tau\|\Sigma\|_{C^2})$.
\end{theorem}

\begin{proof}
By Theorem~\ref{thm:unconditional-bare-remainder} the invariant scalar band response equals the van
Vleck response up to $O(\eps)+O(h_R^2)$ for any admissible channel.  The van Vleck
half-density gives the previous continuum response
$\tau^2(a_i\widehat R_{uu}+b_i\widehat R)$ with truncation error $O(\tau^3/L^3)$.
The projected drift contributes through $-\int_\gamma\Sigma$; by
Lemma~\ref{lem:line-integral} its flat-normalized, profile-averaged value is
$O(\tau^3(\|\Sigma\|_{C^2}+\|\Riem\|_{C^0}))$, i.e.\ $O(\tau\|\Sigma\|_{C^2})$ after
division by $\tau^2$ (the $\|\Riem\|_{C^0}$ piece is absorbed in the existing
$O(\tau/L^3)$ curvature error).  The rate-gradient term is the previous analysis's
$O(\tau\|\nabla\log\lambda\|)$.  Inverting the (assumed invertible) moment matrix
separates the two curvature channels; a symmetric bilinear form is fixed by its
quadratic form on an open timelike cone.
\end{proof}

\begin{corollary}[Tightened error bookkeeping]
\label{cor:tightened}
The term recorded as $O(\|\kren\|)$ in the two-profile separation theorem
and in the validity-window estimate is in fact $O(\tau\,\|\kren+V\|_{C^2})$, which
vanishes as $\tau\to0$.  In particular it never competes with the leading curvature
signal, and the reversible-isotropic restriction may be deleted from those
statements.
\end{corollary}

\begin{proof}
Immediate from Theorem~\ref{thm:unconditional-ricci}: the reconstruction error
attributable to the projected drift is $O(\tau\|\Sigma\|_{C^2})$ on the whole of
$\CV$.
\end{proof}

% =====================================================================
\subsubsection{Closure theorem}
\label{sec:cv-edits}
% =====================================================================

\begin{theorem}[Unconditional WKB--Ricci closure on $\CV$]
\label{thm:unconditional-wkb-ricci}
For every admissible renewal channel \textup{(Definition~\ref{def:admissible-renewal-channel})}:
\begin{enumerate}[leftmargin=2.6em]
\item the displayed bare-discrete order-$\eps$ symbol-remainder estimate holds in
$H^1_\eps\to L^2$, uniformly on moment-and-gap-bounded subfamilies
\textup{(Theorem~\ref{thm:unconditional-bare-remainder})};
\item the band-projected first-order drift splits canonically into anti-Hermitian
geometric transport, a gradient part removed by the ground-state transform, and an
irreversible scalar $\kren$ \textup{(Lemma~\ref{lem:canonical-drift-split})};
\item the two-profile Ricci reconstruction holds with the entire drift dependence
relegated to the subleading residual $O(\tau\|\kren+V\|_{C^2})$
\textup{(Theorem~\ref{thm:unconditional-ricci})}.
\end{enumerate}
Consequently the local renewal-Einstein assembly is valid for the full
Calder\'on--Vaillancourt class, with no reversibility, detailed-balance, or
isotropy hypothesis.  The reversible-isotropic worked class is recovered as the
special case $\kren=0$.
\end{theorem}

\begin{proof}
Combine Theorems~\ref{thm:unconditional-bare-remainder} and \ref{thm:unconditional-ricci} with
Lemma~\ref{lem:canonical-drift-split}.  The screen-thermodynamic side of the assembly is untouched,
so the Clausius balance and the null-reconstruction step proceed verbatim with the
unconditional Ricci input.
\end{proof}

\begin{hypothesis}[Optional spin-isotropising reset]
\label{hyp:spin-iso}
The combined tilted reset channel on the band endomorphism algebra and the
internal renewal algebra is primitive, with a uniform spectral gap and a unique
faithful invariant band density.
\end{hypothesis}

\begin{proposition}[Optional Perron--spin scalarisation]
\label{thm:perron-spin}
Under \Cref{hyp:spin-iso}, the combined tilted transfer operator has a simple
leading Perron eigenvalue and uniformly gapped complementary modes.  Its
large-deviation Hamiltonian is $h(x,\theta)$.
\end{proposition}

\begin{proof}
Apply noncommutative Perron--Frobenius theory to the primitive combined channel
\cite{EvansHoeghKrohn,Schrader}.  Since the principal Dirac eigenvalue is scalar
on the rank-$r$ band, contraction onto the unique invariant density leaves the
same scalar $h$.
\end{proof}

% =====================================================================
\section{Causal-diamond reconstruction and the Einstein equation}
\label{sec:einstein-reconstruction}
% =====================================================================

The WKB phase determines the characteristic propagation law, while its Gaussian
prefactor carries the local curvature response.  We now reconstruct the Ricci
tensor from calibrated causal diamonds, derive the predictive-screen entropy
and modular normalization, and combine both sides in the local area balance
that yields the Einstein equation.

\subsection{Renewal diamonds and Ricci tomography}

The proper-time-deficit phase determines the classical characteristics, but the
phase alone does not contain the local curvature tensor.  The curvature extracted
below is that of the characteristic metric $\wideg$: it equals $g$ at constant
rate and differs from it by the explicit adiabatic rate-gradient terms above.  Curvature appears at the
next order, through the fluctuation determinant around the minimizing history.
We therefore integrate the flat-normalized van Vleck response over small causal
diamonds and ask which contractions of the Ricci tensor can be recovered from
such measurements.  Lorentzian van~Vleck transport and its relation to geodesic
focusing are standard \cite{VisserVanVleck}.  Curvature reconstruction from
discrete causal data also has a direct order-theoretic precedent: the
Benincasa--Dowker operator estimates $R$, while Roy--Sinha--Surya recover $R$ and
$R_{00}$ from chain abundances in small causal diamonds
\cite{BenincasaDowker,RoySinhaSurya}.  The observable used here is different:
it is a flat-normalized kernel determinant and its profile tomography, not an
interval-count estimator.

Fix $p$, a future unit timelike $u$, and $\tau$; let $D_\tau(p,u)$ be the causal
diamond with tips at proper time $\pm\tau/2$.

\begin{definition}[Calibrated renewal diamond measure]
\label{def:diamond-measure}
Let $\Gamma_\tau(p,u)$ be the renewal histories whose continuum image lies in
$D_\tau(p,u)$ with endpoints on the two tips up to $O(\eps)$.  The calibrated
diamond measure is
$d\mu^\eps=Z_\eps^{-1}\exp(-\eps^{-1}S_T(\gamma))\,d\nu^\eps$, with $d\nu^\eps$ the
microscopic counting measure transported to the diamond, $Z_\eps$ its
normalization, and the flat-normalized partition function
$\mathcal Z_\eps(D_\tau)=Z_\eps(D_\tau)/Z_\eps(D_\tau^{\rm flat})$.
\end{definition}

A rotationally symmetric profile greatly simplifies the second-moment tensor,
but it does not isolate $R_{\mu\nu}u^\mu u^\nu$ by itself: the scalar curvature
enters through the transverse moment.  The following lemma makes this mixing
explicit and motivates the use of two independent profiles.

\begin{lemma}[Rotationally symmetric diamond response]
\label{lem:vvm-average}
For a small geodesic diamond of the characteristic metric $\wideg$ and a
normalized profile $\rho_\tau$ invariant under the spatial rotations fixing
$u$, there are coefficients $a_\rho,b_\rho$ such that
\[
 \int_{D_\tau}\log\Delta_{\wideg}(p,q)^{1/2}\rho_\tau(q)\,dq
 =\tau^2\bigl(a_\rho\widehat R_{\mu\nu}u^\mu u^\nu
             +b_\rho\widehat R\bigr)
 +O(\tau^3/L_{\rm curv}^3).
\]
A single positive rotationally symmetric profile need not annihilate
$b_\rho$.
\end{lemma}

\begin{proof}
The Synge--DeWitt expansion for $\wideg$ gives in normal coordinates
$\log\Delta_{\wideg}^{1/2}
 =\tfrac1{12}\widehat R_{\mu\nu}\sigma^\mu\sigma^\nu
 +O(|\sigma|^3)$ \cite{Synge,DeWitt,VisserVanVleck}.  Rotational symmetry gives
$\int\sigma^\mu\sigma^\nu\rho_\tau
 =\tau^2(A_\rho u^\mu u^\nu+B_\rho\widehat h^{\mu\nu})$, and
$\widehat h^{\mu\nu}\widehat R_{\mu\nu}
 =\widehat R+\widehat R_{uu}$.  Collecting the two contractions yields the
stated coefficients.
\end{proof}

Thus a single diamond average is generally insufficient.  Two calibrated
profiles with linearly independent second moments give a $2\times2$ tomography
problem.

\begin{theorem}[Two-profile Ricci separation]
\label{thm:ricci}
Assume \Cref{hyp:SH,hyp:LD}.  If two flat-normalized calibrated diamond measures
converge in second moments to positive rotationally symmetric profiles
$\rho_1,\rho_2$, then, with
$\mathfrak D_i=\log\mathcal Z_{\eps,i}(D_\tau(p,u))$,
\[
 \frac1{\tau^2}\binom{\mathfrak D_1}{\mathfrak D_2}
 =\begin{pmatrix}a_1&b_1\\a_2&b_2\end{pmatrix}
  \binom{\widehat R_{\mu\nu}u^\mu u^\nu}{\widehat R}
 +O(\tau/L_{\rm curv}^3)+O(\eps/\tau^2)
 +O(\tau\|\nabla\log\lambda\|)
+O\!\bigl(\tau\|\kappa_{\ren}+V\|_{C^2(D_\tau)}\bigr).
\]
If the moment matrix is invertible, the responses reconstruct
$\widehat R_{\mu\nu}u^\mu u^\nu$ and $\widehat R$.  Varying $u$ over an open
timelike set reconstructs the full characteristic Ricci tensor by polarization.
\end{theorem}

\begin{proof}
\Cref{lem:vvm-average} gives the continuum response.  The matrix WKB theorem and
\Cref{prop:invariant-response} replace the microscopic response by the invariant
scalar van Vleck response.  The rate-gradient term is controlled by
\Cref{thm:adiabatic-rate}; the scalar reset contribution is the
flat-normalized, reflection-averaged line integral of \(\kren+V\), bounded by
\Cref{lem:line-integral}.  Invertibility of the moment matrix separates the two contractions, and a
symmetric bilinear form is determined by its quadratic form on an open cone.
\end{proof}


\begin{proposition}[Explicit causal-diamond tomography in $d=3$]
\label{prop:explicit-tomography}
In flat normal coordinates on the diamond let
\[
 \rho_{m,n}(\eta,z)=C_{m,n}
 \left(1-\frac{2|\eta|}{\tau}\right)_+^m
 \left[1-\frac{|z|^2}{(\tau/2-|\eta|)^2}\right]_+^n .
\]
These profiles are genuinely supported on $D_\tau$.  Their normalized moments
are
\[
 A(m)=\frac{1}{2(m+d+2)(m+d+3)},\qquad
 B(m,n)=\frac{m+d+1}{4(d+2n+2)(m+d+3)}.
\]
For $d=3$, the pair $(m_1,n_1)=(0,0)$ and $(m_2,n_2)=(0,2)$ gives
\[
 M=\begin{pmatrix}1/240&1/360\\[2pt]19/6480&1/648\end{pmatrix},
 \qquad \det M=-\frac1{583200},
 \qquad M^{-1}=\begin{pmatrix}-900&1620\\1710&-2430\end{pmatrix},
\]
and $C_{\rm rec}=\|M^{-1}\|_2\approx3498$.  Over
$m,n\in\{0,1,2,3\}$, the best-conditioned pair is $(0,3)$ versus $(3,0)$,
with $C_{\rm rec}\approx1530$.
\end{proposition}

\begin{proof}
Writing $R(\eta)=\tau/2-|\eta|$, the $z$-integration contributes the factor
$R(\eta)^d$.  Beta-function integration gives
$\langle\eta^2\rangle/\tau^2=A(m)$ and
$\langle|z|^2\rangle/(d\tau^2)=B(m,n)$.  Since
$a=(A+B)/12$ and $b=B/12$, direct substitution gives the displayed matrix,
determinant, inverse, and singular-value norm.  The finite-grid optimum follows
by evaluating the same exact formulas on the stated exponent set.
\end{proof}


\begin{proposition}[Least-squares diamond tomography]
\label{prop:least-squares-tomography}
For $K>2$ diamond profiles, let $M_K$ be the $K\times2$ matrix with rows
$(a_k,b_k)$ and reconstruct the two curvature channels by the pseudoinverse
$M_K^+=(M_K^\top M_K)^{-1}M_K^\top$.  Then
$C_{\rm rec}=\|M_K^+\|_2=1/\sigma_{\min}(M_K)$, and adding profiles cannot
increase $C_{\rm rec}$.  For the polynomial family of
\Cref{prop:explicit-tomography} in $d=3$,
\[
\begin{array}{c|c}
\text{profile grid} & C_{\rm rec}\\ \hline
\{0,\ldots,3\}^2 & \approx975\\
\{0,\ldots,6\}^2 & \approx495\\
\{0,\ldots,10\}^2 & \approx315
\end{array}
\]
so the degree-$10$ least-squares family improves the best two-profile constant
by about a factor five.
\end{proposition}

\begin{proof}
The least-squares error amplification is the operator norm of the Moore--Penrose
pseudoinverse, hence $1/\sigma_{\min}(M_K)$.  Row addition replaces
$M_K^\top M_K$ by $M_K^\top M_K+rr^\top$, a positive-semidefinite increment, so
its smallest eigenvalue cannot decrease.  The displayed values follow by direct
evaluation of the exact Beta moments in \Cref{prop:explicit-tomography}.
\end{proof}

\begin{remark}[Where the renewal theorem enters]
\label{rem:where-renewal-enters-ricci}
The local van Vleck expansion is standard.  The renewal content is that the
reset-history kernel has the matrix band saddle of \Cref{thm:matrix-wkb}, that
its unitary-invariant response removes the non-abelian Berry holonomy, and that
the remaining scalar Gaussian response converges to the characteristic van Vleck
determinant.  Two profiles are essential because one rotationally symmetric
profile does not separate the longitudinal and scalar-curvature channels.
\end{remark}

\begin{corollary}[Background-metric reconstruction at slowly varying rate]
\label{cor:background-ricci}
At constant rate, $\widehat R_{\mu\nu}=R_{\mu\nu}$.  In the adiabatic window,
\Cref{thm:adiabatic-rate} converts the characteristic Ricci response into the
background response with an additional residual bounded by
$C\tau\|\nabla\log\lambda\|$.  Under the conformal optical calibration of
\Cref{cor:optical-metric}, the conversion is the explicit conformal Ricci formula
there.
\end{corollary}

\subsection{Predictive-screen entropy}

The diamond calculation supplies the geometric side of the field equation.  To
turn it into dynamics one needs an independent thermodynamic side: a local
entropy density, a modular first law, and the normalization relating modular
flow to boosts.  We derive the area extensivity from primitive renewal transfer
channels and then state precisely which modular input is imported from the
general theory of half-sided inclusions.

A local predictive screen is a codimension-two predictive surface separating the
null halves of a small causal diamond.  The modular package of the Einstein
assembly follows from standard modular structure once the screen is stationary and
half-sided modular.  Modular nuclearity provides the established criterion by
which wedge-local constructions acquire nontrivial compact localization
\cite{BuchholzLechner}; we use it only as context for the screen-locality
assumptions and do not derive a new nuclearity theorem.  The only remaining
normalization is the conversion between
screen entropy density and Newton's constant.

\begin{definition}[Primitive matrix-product renewal screen]
\label{def:primitive-mp-screen}
Let $\mathcal E:\mathfrak b\to\mathfrak b$ be a primitive unital CP transfer map
on $\mathfrak b=M_D(\C)$ with invariant $\rho_*>0$ and
$\|\mathcal E^n-Q\|_{1\to1}\le C_{\mathcal E}e^{-n/\xi}$, $Q(a)=\rho_*(a)\mathbf 1$.
Following the finitely correlated-state construction
\cite{FannesNachtergaeleWerner,KretschmannWerner}, a matrix-product renewal strip
comes from an emission isometry $V:\C^D\to\C^D\otimes\mathcal H_{\rm cell}$ with
transfer map $\mathcal E$; a screen patch $\Sigma$ is tiled by $N_\Sigma$ cells and,
up to a finite-depth tangential circuit, is a product of strips.  The split
regulator of width $r$ cuts each strip at distance $r$ and replaces the exterior
memory by $\rho_*$, giving $\rho^{(r)}_\Sigma$.
\end{definition}

Primitivity first implies that auxiliary boundary data are forgotten
exponentially fast along each strip.  This mixing estimate is the mechanism that
localizes all nonextensive entropy contributions near the boundary of the
regulated patch.

\begin{lemma}[Exponential loss of boundary memory]
\label{lem:boundary-memory-loss}
Changing the auxiliary memory at a regulated strip endpoint changes the reduced
state of cells at distance $\ge r$ by at most $C_0e^{-r/\xi}$ in trace norm; for a
patch with $|\partial\Sigma|$ endpoints the total interior error is
$C_0|\partial\Sigma|e^{-r/\xi}$.
\end{lemma}

\begin{proof}
The reduced state beyond $r$ steps depends on the endpoint memory only through
$(\mathcal E^*)^r$; dualizing the $1\to1$ contraction gives convergence to the
invariant memory at rate $C_{\mathcal E}e^{-r/\xi}$, preserved by complete
positivity of later emissions; tensorization gives the boundary factor.
\end{proof}

Once boundary dependence is controlled, the entropy increment along a long
strip converges to a stationary conditional-entropy rate.  Summing these stable
increments produces the extensive contribution that will become the screen area
term.

\begin{lemma}[Entropy rate of a primitive renewal strip]
\label{lem:strip-entropy-rate}
The von Neumann entropy $S_n$ of the first $n$ cells satisfies
$S_n=n\eta_{\rm strip}+b+O(e^{-n/\xi_1})$ with
$\eta_{\rm strip}=\lim_n(S_n-S_{n-1})\ge0$.
\end{lemma}

\begin{proof}
The increment $S_n-S_{n-1}$ depends only on the incoming memory and one fixed
finite-dimensional channel; primitivity makes incoming memories from two boundary
conditions differ by $O(e^{-n/\xi})$, and the Alicki--Fannes continuity bound
\cite{AlickiFannes}, with the fixed cell dimension, gives
$|(S_n-S_{n-1})-\eta_{\rm strip}|\le Ce^{-n/\xi_1}$; summing yields $b$ and the
remainder.
\end{proof}

Combining exponential boundary-memory loss with the existence of the strip
entropy rate gives the desired screen law.  The theorem below separates the
universal extensive term from the regulator-dependent local counterterm and
keeps the collar error quantitative.

\begin{theorem}[Boundary-local area law for primitive renewal screens]
\label{thm:renewal-screen-area-law}
For a screen patch $\Sigma$ of \Cref{def:primitive-mp-screen} with
$\Area_g(\Sigma)=a_\eps N_\Sigma+O(\eps|\partial\Sigma|)$, there are a finite
density $\etaScr=\eta_{\rm cell}/a_\eps\ge0$ and a local boundary counterterm
$B_r(\partial\Sigma)$ such that
\[
 S(\rho^{(r)}_\Sigma)=\etaScr\Area_g(\Sigma)+B_r(\partial\Sigma)
 +O(|\partial\Sigma|e^{-r/\xi_*})+O(\eps|\partial\Sigma|),
\]
so the renormalized split entropy
$S_{\rm split}^{(r)}(\Sigma)=S(\rho^{(r)}_\Sigma)-B_r(\partial\Sigma)$ satisfies the
same area law without the counterterm, and
$S_{\rm split}^{(r)}/\Area_g\to\etaScr$ as $r/\xi_*\to\infty$.
\end{theorem}

\begin{proof}
Omitting the tangential circuit, \Cref{lem:strip-entropy-rate} summed over strips
gives the extensive term $\eta_{\rm cell}N_\Sigma$; strip-end constants form a
local boundary functional $B_r=O(|\partial\Sigma|)$, and replacing the true
exterior memory by the split memory changes each by $O(e^{-r/\xi})$
(\Cref{lem:boundary-memory-loss}).  The conditional-entropy continuity uses only
the fixed collar dimension, giving $O(|\partial\Sigma|e^{-r/\xi_*})$.  A finite-range
tangential circuit has only $O(|\partial\Sigma|)$ collar gates, absorbed into
$B_r$.  Replacing $N_\Sigma$ by $\Area_g/a_\eps$ adds $O(\eps|\partial\Sigma|)$.
\end{proof}

\begin{corollary}[Microscopic screen law and Newton normalization]
\label{cor:microscopic-screen-law}
After adding the bulk-state entropy,
$S_{\rm scr}^{(r)}(\Sigma)=\etaScr\Area_g(\Sigma)+S_{\bulk}(\Sigma)
+O(|\partial\Sigma|e^{-r/\xi_*})+O(\eps|\partial\Sigma|)$, with coupling
$\Gren=1/(4\etaScr)$; $\etaScr$ is model-dependent but computable from the
stationary conditional-entropy rate of the transfer channel.
\end{corollary}

\begin{proof}
Combine \Cref{thm:renewal-screen-area-law} with the split-regulated bulk entropy
and \Cref{prop:newton-from-screen-density}.
\end{proof}

\begin{definition}[Resolved predictive-pure screen]
\label{def:predictive-pure-screen}
Let \(Y\) be the retained branch record crossing one regenerative screen cell,
and let \(S_{\rm edge}\) be the edge causal state, the minimal sufficient
statistic for the future continuation law.  The record is \emph{resolved} when
\(S_{\rm edge}=f(Y)\) for a deterministic map \(f\).  Its screen state has a
central decomposition
\[
 \rho_\Sigma=\bigoplus_y p_y\,\rho_y^L\otimes\rho_y^R.
\]
The screen algebra is \emph{predictive pure} when each resolved causal-state
sector has a one-dimensional conditional GNS representation, equivalently
\[
   S(\rho_y^L)=S(\rho_y^R)=0
   \qquad\text{for every }y .
\]
\end{definition}

\begin{lemma}[Central screen-entropy decomposition]
\label{lem:central-screen-entropy}
For the central state of \Cref{def:predictive-pure-screen},
\[
 \boxed{\;
 S(\rho_\Sigma)
 =H(p)+\sum_y p_yS(\rho_y^L)+\sum_y p_yS(\rho_y^R).
 \;}
\]
\end{lemma}

\begin{proof}
The eigenvalues of the block
\(p_y\rho_y^L\otimes\rho_y^R\) are \(p_y\lambda_i^{(y)}\mu_j^{(y)}\),
where \(\lambda_i^{(y)}\) and \(\mu_j^{(y)}\) are the eigenvalues of the two
conditional density matrices.  Substitution into
\(-\Tr\rho_\Sigma\log\rho_\Sigma\), followed by summation over the orthogonal
central blocks, gives the Shannon term \(H(p)\) and the two averaged conditional
von Neumann entropies.  Tensor-product entropy is additive.
\end{proof}

\begin{theorem}[Predictive-pure entropy rigidity]
\label{thm:predictive-pure-entropy}
Assume the record is resolved and the screen algebra is predictive pure.  Then
the quantum entropy per cell equals the entropy of the resolved predictive
record.  For a stationary first-order regenerative record with transition
matrix \(\Pi\),
\[
 \boxed{\quad
 h_{\rm scr}^{q}
 =h(\Pi)
 =-\sum_{x,y}\pi_x\Pi_{xy}\log\Pi_{xy}.
 \quad}
\]
If the entropy is instead computed on a larger microscopic algebra, then the
only additional term is the nonnegative conditional density
\[
 \Scond
 =\mathbb E_Y\!\left[S(\rho_Y^L)+S(\rho_Y^R)\right]\ge0.
\]
It measures residual non-predictive memory or bridge leakage and is not a
negative coherent correction to the predictive-pure entropy.
\end{theorem}

\begin{proof}
Apply \Cref{lem:central-screen-entropy} to a strip of \(N\) regenerative screen
cells.  Predictive purity removes every conditional quantum entropy, so the
remaining extensive term is the Shannon entropy of the resolved record
\(Y_1,\ldots,Y_N\).  Dividing by \(N\) and taking the stationary limit gives the
Shannon entropy rate.  For a first-order stationary Markov record that rate is
\(H(Y_{n+1}\mid Y_n)\), which is the displayed expression
\cite{CoverThomas}.  On a larger algebra the same central decomposition retains
the two conditional von Neumann terms; their average is nonnegative by
positivity of entropy.
\end{proof}

\begin{corollary}[Uniform \(K_4\) screen]
\label{cor:uniform-k4-screen}
For the uniform \(K_4\) regenerative cell, each resolved state has three
equiprobable admissible successors.  Therefore
\[
 \boxed{\qquad h_{\rm scr}^{q}=h(\Pi)=\log3.\qquad}
\]
\end{corollary}

\begin{proof}
Every row of the embedded transition matrix has the three nonzero entries
\((1/3,1/3,1/3)\).  Its conditional entropy is \(\log3\), independently of the
uniform stationary vertex label.  Apply
\Cref{thm:predictive-pure-entropy}.
\end{proof}

\begin{corollary}[Predictive-pure Newton normalization]
\label{cor:k4-newton-normalization}
Assume that the exact metric area of one normalized screen cell is
\[
   a_\eps=\Npartial\,\eps^{d-1},
   \qquad \Npartial>0.
\]
Then a predictive-pure stationary screen has
\[
 \boxed{\quad
 \etaScr=\frac{h(\Pi)}{\Npartial\,\eps^{d-1}},
 \qquad
 \Gren=\frac{\Npartial\,\eps^{d-1}}{4h(\Pi)}.
 \quad}
\]
For the uniform \(K_4\) screen,
\[
 \Gren=\frac{\Npartial\,\eps^{d-1}}{4\log3}.
\]
In \(d=3\), under the normalized screen-area convention
\(\Npartial=1\),
\[
 \boxed{\qquad
 \eps^2=4\log3\,\Gren.
 \qquad}
\]
\end{corollary}

\begin{proof}
The extensive entropy of \(N_\Sigma\) predictive-pure cells is
\(h(\Pi)N_\Sigma+o(N_\Sigma)\), whereas their metric area is
\(a_\eps N_\Sigma+o(N_\Sigma)\).  Their ratio gives
\(\etaScr=h(\Pi)/a_\eps\).  The Newton formula is
\(\Gren=(4\etaScr)^{-1}\) by
\Cref{prop:newton-from-screen-density}.  Substituting the uniform \(K_4\)
entropy and then \(d=3,\Npartial=1\) gives the final identities.
\end{proof}

\begin{remark}[Scope of the normalized \(K_4\) value]
\label{rem:k4-screen-scope}
The general area-law theorem requires only \(a_\eps\asymp\eps^{d-1}\).  The
numerical identity \(\eps^2=4\log3\,\Gren\) additionally fixes the finite
screen-area coefficient to \(\Npartial=1\) and identifies the physical screen
with the resolved predictive-pure quotient.  If a larger screen algebra is
used, \(\Scond>0\) increases the entropy density; if
\(\Npartial\ne1\), the displayed factor is modified accordingly.
\end{remark}



\begin{remark}[Quadratic record-weight normalization]
\label{rem:quadratic-record-weights}
The screen area law uses the entropy of regulated predictive states and is
sufficient for the Einstein assembly.  When a downstream calculation resolves a
screen into mutually exclusive predictive records with orthogonal projectors
$\Pi_i$, the record-disintegration theorem of
\cite{PelissierBornDisintegrationArxiv} supplies the compatible probability
normalization
\[
        p_i=\frac{\|\Pi_i\psi\|_{\rm pred}^2}
        {\sum_j\|\Pi_j\psi\|_{\rm pred}^2}.
\]
This result is not an additional hypothesis in the local Einstein derivation; it
clarifies that the statistical weights used by refined predictive-screen
calculations are the quadratic weights compatible with additivity,
multiplicativity, and predictive invariance.
\end{remark}


\begin{lemma}[Screen normalization bookkeeping]
\label{lem:screen-normalization}
Let $a_\eps\asymp\eps^2$ be the metric area per screen cell,
$\eta_{\rm cell}=\etaScr a_\eps$ the dimensionless entropy per cell, and
$\etaScr$ the continuum entropy density per unit metric area.  Assume the split
collar has fixed bond dimension $D$, independent of $\eps$ and of the screen.
Then the dimensionless curvature seen by one cell is
$a_\eps\mathcal R_\Sigma$ and $a_\eps|K|^2$, while the conditional-entropy
continuity constant is uniformly $O(\log D)=O(1)$.
\end{lemma}

\begin{hypothesis}[Bounded-geometry curved screen]
\label{hyp:curved-screen}
The local screen $\Sigma$ has bounded intrinsic curvature and second fundamental
form $K$.  Its per-cell transfer channel $\mathcal E_x$ varies $C^2$ and satisfies
\[
 \|\mathcal E_x-\mathcal E\|_{\rm cb}
 \le C_0\eps^2\bigl(|\mathcal R_\Sigma(x)|+|K(x)|^2\bigr),
\]
where $\mathcal E$ is the flat primitive channel.
\end{hypothesis}

\begin{theorem}[Normalized curved-screen stability]
\label{thm:curved-screen}
Under \Cref{lem:screen-normalization,hyp:curved-screen},
\[
 \left|S_{\rm scr}(\Sigma)-\etaScr\Area_g(\Sigma)-S_{\bulk}(\Sigma)\right|
 \le C\etaScr\Area_g(\Sigma)
 \left(e^{-r/\xi}+\frac{\eps^2}{L_{\rm curv}^2}
 +\eps^2\|K\|_\infty^2\right).
\]
Thus the fractional curved-screen correction is of order
$O(\eps^2/L_{\rm curv}^2+\eps^2\|K\|_\infty^2)$.
\end{theorem}

\begin{proof}
Analytic perturbation of a gapped primitive channel makes the per-cell entropy
rate Lipschitz in the dimensionless channel perturbation
$a_\eps(|\mathcal R_\Sigma|+|K|^2)$ \cite{Kato,WolfPerez}.  The fixed collar
dimension makes the Alicki--Fannes constant uniform.  Summing over
$N_\Sigma=\Area_g/a_\eps$ cells cancels the cell-area factor; since
$\eta_{\rm cell}=\etaScr a_\eps$, the remaining absolute correction is
$C\etaScr\Area_g\,a_\eps$ times the curvature scale.  Using
$a_\eps\asymp\eps^2$ and bounded geometry gives the stated estimate, while
exponential clustering supplies the split-collar term.
\end{proof}

The area law fixes the entropy carried by the screen, but it does not by itself
identify the physical temperature or modular Hamiltonian.  For that purpose we
now place the screen inside a stationary local wedge and state the modular
properties required to identify its intrinsic flow with the emergent boost
flow.


\subsection{Modular calibration}

\begin{hypothesis}[Quasi-free renewal-screen scaling limit]
\label{hyp:quasifree-screen}
The scaling limit is a quasi-free CCR/CAR net with positive proper-time generator
$P_1\ge0$, wedge isotony, wedge duality, and a globally cyclic stationary vector.
Its one-particle dynamics obeys the causal-support propagation theorem of
\Cref{thm:causal-propagation}.
\end{hypothesis}

\begin{definition}[One-particle renewal wedge subspace]
\label{def:renewal-wedge-subspace}
For a renewal wedge $W$, define
\[
 K(W)=\overline{\{E_1f:\ f\text{ real},\ \supp f\subseteq W\}}^{\R},
\]
where $E_1$ is the one-particle field map.
\end{definition}

\begin{theorem}[Derived half-sided covariance]
\label{thm:derived-half-sided}
Under \Cref{hyp:quasifree-screen},
\[
 U_1(s)K(W)\subseteq K(W),\qquad s\ge0.
\]
Moreover the stationary vector is cyclic and separating for the wedge algebra,
so $K(W)$ is standard.
\end{theorem}

\begin{proof}
Causal propagation gives
$\supp(U_1(s)f)\subseteq J_s^+(\supp f)$.  Positive proper-time translation
maps the chosen wedge into itself, and wedge isotony therefore maps every
generator of $K(W)$ back into $K(W)$.  Closedness gives the inclusion.  Positivity
of $P_1$ yields upper-half-plane analyticity of translated matrix elements; the
standard Reeh--Schlieder argument, together with global cyclicity and wedge
duality, gives cyclicity and separatingness \cite{ReehSchlieder,StreaterWightman,BGL2}.
\end{proof}

\begin{theorem}[Half-sided modular inclusion for quasi-free renewal screens]
\label{thm:quasifree-modular}
Under \Cref{hyp:quasifree-screen}, the translated wedge algebra is a positive
half-sided modular inclusion and
\[
 \Delta_W^{it}U(s)\Delta_W^{-it}=U(e^{-2\pi t}s),
 \qquad J_WU(s)J_W=U(-s).
\]
The Borchers translation is the renewal proper-time translation, the modular
flow is the geometric boost, and the stationary state is KMS at
$T=\kappa/(2\pi)$.
\end{theorem}

\begin{proof}
\Cref{thm:derived-half-sided} makes $(K(W),U_1)$ a positive half-sided standard
pair.  Borchers--Wiesbrock theory gives the one-particle relations
\cite{Borchers,Wiesbrock,Longo,BGL}.  Quasi-free second quantization transports
these modular objects to the wedge algebra
\cite{EckmannOsterwalder,FiglioliniGuido}; restoring the boost normalization
gives the Unruh temperature.
\end{proof}


\begin{definition}[Precise modular inputs]
\label{def:precise-modular-inputs}
The quasi-free renewal realization uses: \textup{(M1)} a Rindler renewal wedge
$W$ with a future null translation $U(s)W\subseteq W$ for $s\ge0$;
\textup{(M2)} renewal microcausality and finite propagation;
\textup{(M3)} wedge duality $\A(W)'=\A(W')$; and
\textup{(M4)} global cyclicity of the stationary vector.
\end{definition}

\begin{proposition}[Concrete generalized-free renewal screen]
\label{prop:concrete-quasifree-screen}
Let the screen scaling limit be the generalized free field with two-point
function
\[
 W_1(x,y)=\int_{m_0^2}^{\infty}d\varrho(\mu)\,\Delta^+_\mu(x-y),
 \qquad m_0>0,
\]
where $d\varrho$ is positive and $\Delta^+_\mu$ is computed in the renewal
metric.  Then the transfer gap supplies $m_0>0$ and positive energy, while
microcausality, wedge duality, and Reeh--Schlieder give
\textup{(M1)--(M4)}.  Hence \Cref{thm:quasifree-modular} applies and the
stationary state has Unruh temperature $T=\kappa/(2\pi)$.
\end{proposition}

\begin{proof}
The Rindler form of the renewal wedge gives \textup{(M1)}.  Each massive
Wightman function is local, and positivity of $d\varrho$ preserves the spectrum
condition, giving \textup{(M2)}.  Massive generalized free fields satisfy wedge
duality and Reeh--Schlieder cyclicity, giving \textup{(M3)--(M4)}
\cite{Araki,Driessler,ReehSchlieder}.  The derived half-sided inclusion and
Borchers--Wiesbrock theorem then yield the claimed modular normalization.
\end{proof}

\begin{hypothesis}[Stationary modular wedge with boundary-local entropy]
\label{hyp:stationary-primitive-wedge}
For a small local wedge $\mathcal W$ with renewal algebra $\A(\mathcal W)$ assume:
\textup{(i)} a normal stationary screen state $\omega$ with cyclic separating
$\Omega$; \textup{(ii)} the translated nested wedge is a positive half-sided
modular inclusion; \textup{(iii)} either the screen is in the primitive
matrix-product class of \Cref{def:primitive-mp-screen}, or its regulated entropy
obeys a boundary-locality estimate of the same form.  For the matrix-product class
(iii) is discharged by \Cref{thm:renewal-screen-area-law}.
\end{hypothesis}

\begin{theorem}[Local renewal calibration]
\label{thm:local-calibration}
Under \Cref{hyp:stationary-primitive-wedge}: \textup{(i)} the modular flow of
$(\A(\mathcal W),\omega)$ exists and $\omega$ is KMS at inverse temperature one;
\textup{(ii)} the half-sided inclusion gives a positive-energy translation group
with $\mods^{it}U(s)\mods^{-it}=U(e^{-2\pi t}s)$, and when identified with the
emergent wedge translations the modular generator is the boost generator, with
$T=a/2\pi$; \textup{(iii)} for every differentiable family of normal states with
finite Araki relative entropy, $\delta S_{\rm rel}=\delta\langle\Kmod\rangle$,
$\Kmod=-\log\mods$; \textup{(iv)} in the matrix-product class
\Cref{thm:renewal-screen-area-law} gives the area law, and
$\Gren=1/(4\etaScr)$ gives the leading area term $\Area_g(\Sigma)/(4\Gren)$.  Thus
the KMS/Unruh/first-law package follows from the inclusion, while area extensivity
needs the separate boundary-local input.
\end{theorem}

\begin{proof}
Items (i)--(iii) are imported modular theory applied to the renewal wedge:
Tomita--Takesaki gives (i); Borchers--Wiesbrock applied to the half-sided
inclusion gives the affine relation in (ii), the geometric Unruh reading using the
identification with wedge translations
\cite{Borchers,Wiesbrock,BisognanoWichmann}; (iii) is the relative-entropy first
law, valid without a density-matrix entropy for a type-III algebra, and the
modular Hamiltonian of the local diamond is the boost-energy integral of
\cite{CasiniHuertaMyers}.  Item (iv) in the matrix-product class is
\Cref{thm:renewal-screen-area-law}; otherwise it is (iii) of the hypothesis.  The
conversion of $\etaScr$ to $\Gren$ is a normalization.
\end{proof}

The conclusion should be read in two logically distinct parts.  Modular theory
fixes the KMS and Unruh normalization once the half-sided inclusion is known,
whereas renewal transfer theory fixes the existence and microscopic value of the
area density.  Their conjunction determines the Newton normalization, but
neither part alone does so.

\begin{remark}[Normalization scope]
\label{rem:normalization-scope}
\Cref{thm:local-calibration} derives KMS thermality, the Unruh normalization, and
the first law; for primitive matrix-product screens
\Cref{thm:renewal-screen-area-law} also derives area extensivity and identifies
the density with a conditional-entropy rate.  It does not derive a universal
numerical value of that density: the $1/4$ enters through $\Gren=1/(4\etaScr)$,
while $\etaScr$ is a computable microscopic parameter.
\end{remark}

\begin{theorem}[Modular calibration equivalence; imported from Borchers--Wiesbrock]
\label{thm:modular-calibration-equivalence}
For a local renewal wedge algebra with cyclic separating $\Omega$, modular
operator $\mods$, and strongly continuous translation group $U(s)=e^{isP}$, the
following are equivalent: \textup{(i)} the modular flow acts as the geometric
boost (KMS at the local Unruh temperature); \textup{(ii)} the nested translated
inclusion is a positive half-sided modular inclusion; \textup{(iii)} the Borchers
relations $\mods^{it}U(s)\mods^{-it}=U(e^{-2\pi t}s)$, $\J U(s)\J=U(-s)$ hold with
$P\ge0$.
\end{theorem}

\begin{proof}
This is the Borchers--Wiesbrock equivalence
\cite{Borchers,Wiesbrock,ArakiZsido2005} applied to the renewal wedge; no
renewal-specific structure enters.  Borchers' theorem gives (i)$\Rightarrow$(iii);
the affine relations imply the half-sided inclusion (iii)$\Rightarrow$(ii); and
Wiesbrock's theorem reconstructs the positive-energy translations and the boost
representation from a positive half-sided inclusion (ii)$\Rightarrow$(i).
\end{proof}

\begin{proposition}[The entropy package is not an equivalence]
\label{prop:entropy-no-iff}
Geometric modular calibration does not imply finite correlation length and does
not determine the screen entropy density.  The correct equivalence is only the
modular one of \Cref{thm:modular-calibration-equivalence}; finite-correlation area
extensivity is a sufficient additional hypothesis whose area coefficient defines
$G$.
\end{proposition}

\begin{proof}
Modular data can be geometrically calibrated even in scale-invariant situations
with polynomially decaying correlations, and the ultraviolet area coefficient is
cutoff-dependent and not fixed by the modular group; the forward theorem uses
finite correlation length only for extensivity and cannot be reversed.
\end{proof}

\begin{hypothesis}[Calibration and local screen entropy]
\label{hyp:calibration}
For every small calibrated renewal wedge: \textup{(i)} the modular flow is KMS at
$T=a/2\pi$; \textup{(ii)} the screen entropy has the metric-area expansion
$S_{\rm scr}^{(r)}(\Sigma)=\Area_g(\Sigma)/4G+S_{\bulk}(\Sigma)
+O(|\partial\Sigma|e^{-r/\xi_*})+O(\eps|\partial\Sigma|)$, proved by
\Cref{thm:renewal-screen-area-law} for primitive matrix-product screens and assumed
otherwise; \textup{(iii)} the bulk first variation is
$\delta S_{\bulk}=\delta\langle K_{\rm mod}\rangle
=2\pi\int_{\mathcal H}\Tren_{\mu\nu}\chi^\mu d\Sigma^\nu+O(\mathrm{noncal.})$, $\chi$
the boost Killing field.
\end{hypothesis}

\begin{proposition}[Newton coupling from screen density]
\label{prop:newton-from-screen-density}
If calibrated screens have leading entropy density $\etaScr$, the coupling is
$\Gren=1/(4\etaScr)$.
\end{proposition}

\begin{proof}
With $S_{\rm area}=\etaScr A$, the linearized Raychaudhuri equation gives the area
variation $\etaScr R_{\mu\nu}k^\mu k^\nu$, the KMS first law gives the modular flux
$2\pi\Tren_{\mu\nu}k^\mu k^\nu$, and comparison with
$R_{\mu\nu}k^\mu k^\nu=8\pi\Gren\Tren_{\mu\nu}k^\mu k^\nu$ gives
$8\pi\Gren=2\pi/\etaScr$.
\end{proof}

\begin{proposition}[Local KMS--Clausius identity]
\label{prop:kms}
Under \Cref{hyp:calibration}, $\delta Q=T\,\delta S_{\bulk}+O(\mathrm{noncal.})$,
$T=a/2\pi$.
\end{proposition}

\begin{proof}
The KMS condition identifies the modular flow with the boost flow, the modular
first law gives $\delta S_{\bulk}=\delta\langle K_{\rm mod}\rangle$, and in the
Rindler approximation $K_{\rm mod}$ is the boost-energy integral
$2\pi\int\Tren_{\mu\nu}\chi^\mu d\Sigma^\nu$ of \cite{CasiniHuertaMyers}, so
multiplication by $T$ gives the heat flux.
\end{proof}

\subsection{Local area balance and the Einstein equation}

We have now obtained both ingredients that were missing from the kinematic
construction: the Ricci response of a renewal diamond and the thermodynamic
response of its predictive screen.  The remaining step is the local balance
argument.  We first record the geometric area variation and the elementary null
reconstruction statement, and then combine them with the modular Clausius law.

\begin{lemma}[Renewal Raychaudhuri area variation]
\label{lem:raychaudhuri}
For a calibrated local diamond with null generator $k^\mu$,
$\delta S_{\rm area}=-\tfrac1{4G}\int_{\mathcal H}R_{\mu\nu}k^\mu k^\nu\,\eta\,d\eta\,dA
+O(\theta^2+\sigma^2)+O(\eps)$.
\end{lemma}

\begin{proof}
The renewal characteristics are Lorentzian null generators
(\Cref{thm:causal-propagation}); transverse area obeys Raychaudhuri's equation
\cite{Raychaudhuri,Poisson}.  Setting expansion and shear to zero to first order at
equilibrium, integrating, and using $S_{\rm area}=A/4G$ gives the formula.
\end{proof}

The Raychaudhuri calculation gives only null contractions.  The following
linear-algebra lemma is therefore what promotes equality on every local null
generator to equality of symmetric tensors up to a scalar multiple of the
metric.

\begin{lemma}[Null reconstruction; linear algebra]
\label{lem:null-reconstruction}
In a Lorentzian vector space of dimension $\ge3$, if a symmetric $A_{\mu\nu}$ has
$A_{\mu\nu}k^\mu k^\nu=0$ for all null $k$, then $A_{\mu\nu}=fg_{\mu\nu}$.
\end{lemma}

\begin{proof}
The standard null-reconstruction fact used in Jacobson-type derivations
\cite{Jacobson}: in an orthonormal frame every $(1,n)$ with $|n|=1$ is null, so
$A_{00}+2A_{0i}n^i+A_{ij}n^in^j=0$ for all $n$; replacing $n$ by $-n$ gives
$A_{0i}=0$, and $A_{ij}n^in^j=-A_{00}$ for all $|n|=1$ gives
$A_{ij}=-A_{00}\delta_{ij}$ by polarization.
\end{proof}

We can now assemble the three layers: the renewal Gaussian response supplies
$R_{\mu\nu}$, the screen law supplies the entropy variation, and modular
calibration supplies the heat flux and its normalization.  The resulting field
equation is local, and its error term retains the distinct microscopic,
curvature, calibration, and nonlinear contributions.


\begin{proposition}[Renewal stress conservation]
\label{prop:renewal-noether}
Let $W[g,\lambda,\Phi]$ be the diffeomorphism-invariant effective renewal
functional and
$T^{\ren}_{\mu\nu}=2(\sqrt{-g})^{-1}\delta W/\delta g^{\mu\nu}$.
On the $\lambda$- and $\Phi$-equations of motion,
\[
 \nabla^\mu T^{\ren}_{\mu\nu}=0.
\]
\end{proposition}

\begin{proof}
For a compactly supported vector field $\xi$, diffeomorphism invariance gives
$0=\delta_\xi W$.  On the nonmetric equations of motion, only
$\delta g^{\mu\nu}=-(\nabla^\mu\xi^\nu+\nabla^\nu\xi^\mu)$ remains.  Integration
by parts yields
$0=\int\sqrt{-g}\,\xi^\nu\nabla^\mu T^{\ren}_{\mu\nu}$; arbitrariness of $\xi$
gives the claim \cite{Noether,WaldBook}.
\end{proof}

\begin{theorem}[Local renewal Einstein equation]
\label{thm:einstein}
Under \Cref{hyp:SH,hyp:LD,hyp:calibration}, the local renewal balance law for all
calibrated null wedges implies
$R_{\mu\nu}-\tfrac12R g_{\mu\nu}+\Lambda g_{\mu\nu}=8\pi G\,\Tren_{\mu\nu}+\mathcal R_\eps$,
with $\mathcal R_\eps=O(\eps)+O(\tau^3/L_{\rm curv}^3)+O(\theta^2+\sigma^2)+O(\mathrm{noncal.})$.
\end{theorem}

\begin{proof}
The Clausius balance $\delta Q=T\,\delta S_{\rm scr}$ gives, by
\Cref{hyp:calibration,prop:kms}, the modular flux of $\Tren_{\mu\nu}$, and by
\Cref{lem:raychaudhuri} the area part controlled by $R_{\mu\nu}k^\mu k^\nu$; hence
$(R_{\mu\nu}-8\pi G\Tren_{\mu\nu})k^\mu k^\nu=\text{residual}(k)$ for all null $k$.
\Cref{lem:null-reconstruction} gives $R_{\mu\nu}-8\pi G\Tren_{\mu\nu}=fg_{\mu\nu}$;
rewriting and taking the divergence with the contracted Bianchi identity and
\Cref{prop:renewal-noether} gives $\nabla_\nu(f-R/2)=0$, so $f-R/2=-\Lambda$ is
locally constant, yielding the displayed equation with the residual from
\Cref{hyp:LD,hyp:calibration,lem:raychaudhuri}.
\end{proof}

\subsection{Residual estimates and validity}

For later use we collect the quantitative validity statement.  Let
$\delta_{\rm cal}(D_\tau)$ be the operator norm of the difference between the
first variations of the modular Hamiltonian and the local boost-energy
functional on the chosen state perturbations, and let
$E_{\rm WKB}(\eps,\tau)$ be the flat-normalized WKB remainder after quadratic
diamond normalization.

\begin{theorem}[Validity window and stable residual]
\label{thm:validity-window}
Assume the WKB--van Vleck expansion, the curved-screen stability theorem, and
the two-profile Ricci reconstruction.  In the regime
$\eps\ll\tau\ll L_{\rm curv}$, $\delta_{\rm cal}(D_\tau)\ll1$, and
$\tau(|\theta|+|\sigma|)\ll1$, the Einstein residual satisfies
\[
\begin{split}
 \|\mathcal R_\eps\|_{C^0(D_\tau)}\le C_{\rm rec}\Big[
 &E_{\rm WKB}(\eps,\tau)
 +C_2\tfrac{\tau}{L_{\rm curv}}\|\Riem\|_{C^0(D_\tau)}
 +C_3\tfrac{\delta_{\rm cal}(D_\tau)}{\tau^2}\\
 &+C_4(\|\theta\|_{C^0}^2+\|\sigma\|_{C^0}^2)
 +C_5\tau\|\nabla\log\lambda\|_{C^0(D_\tau)}\\
 &+C_6\tau\|\kappa_{\ren}+V\|_{C^2(D_\tau)}
 +C_7\eps^2\big(L_{\rm curv}^{-2}+\|K\|_{C^0}^2\big)\Big].
\end{split}
\]
Here $C_{\rm rec}$ is the stability norm of the inverse two-profile moment
matrix.  For the explicit profiles of \Cref{prop:explicit-tomography}, one may
take $C_{\rm rec}=\|M^{-1}\|_2$ with the closed form displayed there.  On the
reversible isotropic reset class the first-order drift vanishes by
\Cref{thm:drift-cancellation}; on the quasi-free renewal-screen class,
$\delta_{\rm cal}=0$ by \Cref{thm:quasifree-modular}.
\end{theorem}

\begin{corollary}[Residual decomposition]
\label{cor:residual-decomposition}
In the validity window of \Cref{thm:validity-window},
$\mathcal R_\eps=\Rmicro+\Rcurv+\Rcal+\Rshear+\mathcal R_\lambda+\mathcal R_{\rm scr}$,
with
\(\Rmicro=O(\eps)+O\!\bigl(\tau\|\kappa_{\ren}+V\|_{C^2}\bigr)\),
$\Rcurv=O(\tau^3/L_{\rm curv}^3)$, $\Rcal=O(\delta_{\rm cal})$,
$\Rshear=O(\theta^2+\sigma^2)$,
$\mathcal R_\lambda=O(\tau\|\nabla\log\lambda\|)$, and
$\mathcal R_{\rm scr}=O(\eps^2L_{\rm curv}^{-2}+\eps^2\|K\|^2)$; the continuum equation is recovered in the joint
limit $\eps/\tau\to0$, $\tau/L_{\rm curv}\to0$, $\delta_{\rm cal}\to0$,
$(\theta,\sigma)\to0$.
\end{corollary}

\begin{proof}
The tensorial form of \Cref{thm:validity-window}; the four terms are finite
renewal scale, van Vleck truncation, calibration failure, and quadratic
Raychaudhuri terms.
\end{proof}

\begin{remark}[Why the derivation is not circular]
The metric is imported, but the field equation is not: it is forced by requiring
the same diamonds to satisfy both the renewal large-deviation curvature response
and the calibrated screen-entropy balance.  The Einstein tensor is the consistency
condition between propagation and thermodynamics.
\end{remark}

% =====================================================================
\section{Homogeneous renewal cosmology and de Sitter geometry}
\label{sec:desitter}
% =====================================================================

The local Einstein equation does not by itself select a cosmological branch.
We therefore restrict to the homogeneous sector in which the only source is the
renewal-deficiency budget.  This reduction isolates the sign and magnitude of
the effective cosmological term.  We then distinguish carefully between the
local constant-curvature theorem, the additional assumptions needed for a
global completion, and the perturbative stability and horizon consequences of
the resulting branch.

The local field equation has an undetermined integration constant $\Lambda$.  In
the homogeneous pure-deficiency sector, renewal balance fixes its sign and
magnitude.  Local positive constant curvature is a theorem of the leading field
equation; identification of the global maximal continuation with complete de
Sitter space is a separate theorem conditional on the predictive-completion
principle.

\subsection{Deficiency current and the de Sitter branch}

Homogeneity removes the local stress anisotropies and leaves a single scalar
rate $B$.  The first theorem identifies the corresponding Hubble scale.  The
following isotropy reduction then determines the full curvature tensor and
separates the de Sitter, Minkowski, and formally anti-de Sitter possibilities.

Let $B\ge0$ be the homogeneous renewal deficiency budget; in the pure-deficiency
sector $\Tren_{\mu\nu}=0$ and the only macroscopic scalar is $B$.

\begin{theorem}[Pure-deficiency Hubble relation]
\label{thm:lambda}
In the calibrated homogeneous pure-deficiency sector the balance law fixes
$H=B/d$.  Restricted to a maximally symmetric vacuum solution,
$\Lambda=\tfrac{d(d-1)}2H^2=\tfrac{(d-1)B^2}{2d}\ge0$.  For $B\ge0$, $B>0$ is the
expanding de Sitter branch and $B=0$ the Minkowski branch.
\end{theorem}

\begin{proof}
The first identity is the entropy-production balance $dH=B$.  On a maximally
symmetric vacuum branch $R_{\mu\nu}=dH^2g_{\mu\nu}$, $R=d(d+1)H^2$; substitution in
the vacuum Einstein equation gives $\Lambda=d(d-1)H^2/2$, hence the formula via
$H=B/d$.
\end{proof}


\begin{proposition}[Common-origin gravitational scalar]
\label{prop:common-origin-gravity-scalar}
Assume the embedded-chain normalization of
\Cref{cor:embedded-renewal-dictionary}, the predictive-pure screen area
\(a_\eps=\Npartial\eps^{d-1}\), and the homogeneous pure-deficiency relation
\[
   \Lambda=\frac{d-1}{2d}\Bphys^2.
\]
Then
\[
 \boxed{\quad
 \Gren\Lambda
 =
 \Npartial\,\frac{d-1}{8d}\,
 \frac{\Bstep^2c_{\ren}^2}{h(\Pi)}\,
 \eps^{d-3}.
 \quad}
\]
At the \(d=3\), unit-speed endpoint,
\[
   \Gren\Lambda
   =\Npartial\,\frac{\Bstep^2}{12h(\Pi)}.
\]
For the normalized predictive-pure uniform \(K_4\) screen
\((\Npartial=1,\ h(\Pi)=\log3)\),
\[
 \boxed{\qquad
 \Gren\Lambda=\frac{\Bstep^2}{12\log3}.
 \qquad}
\]
\end{proposition}

\begin{proof}
By \Cref{cor:k4-newton-normalization} in its general predictive-pure form,
\[
   \Gren=\frac{\Npartial\eps^{d-1}}{4h(\Pi)}.
\]
By \Cref{cor:embedded-renewal-dictionary},
\[
   \Bphys=\frac{\Bstep c_{\ren}}{\eps}.
\]
Multiplying the first identity by
\(\Lambda=(d-1)\Bphys^2/(2d)\) gives
\[
 \Gren\Lambda
 =
 \frac{\Npartial\eps^{d-1}}{4h(\Pi)}
 \frac{d-1}{2d}
 \frac{\Bstep^2c_{\ren}^2}{\eps^2},
\]
which is the first formula.  Setting \(d=3\) cancels the metric length; setting
\(c_{\ren}=1\), \(\Npartial=1\), and \(h(\Pi)=\log3\) gives the final result.
\end{proof}

\begin{remark}[Interpretation]
The cancellation of \(\eps\) at \(d=3\) is a normalization theorem, not a
prediction of the observed value of \(\Lambda\): the signed increment
\(\Bstep\) remains a state-dependent microscopic current.  The proposition
does show that, once the uniform predictive-pure \(K_4\) branch and the
unit-speed screen normalization are selected, no additional metric length
enters the dimensionless product \(\Gren\Lambda\).
\end{remark}


\begin{proposition}[Bridge displacement from the homogeneous current]
\label{prop:bridge-displacement}
Consider a calibrated homogeneous bridge coordinate $q=e^{-\chi_*}$ whose
dimensionless Perron rapidity satisfies
\[
        \sinh\chi_*=B\eps .
\]
Then, in the controlled continuum regime $B\eps\ll1$,
\[
        1-q=1-e^{-\operatorname{arcsinh}(B\eps)}
        =B\eps+O((B\eps)^2)
        =dH\eps+O((H\eps)^2).
\]
Thus the same homogeneous signed current that fixes the de Sitter Hubble rate
also fixes the leading dimensionless displacement of the bridge boundary layer
used in the early-universe construction \cite{PelissierBridge,PelissierInflationBaryogenesisArxiv}.
\end{proposition}

\begin{proof}
The first equality is the definition of the calibrated bridge coordinate.  The
Taylor expansions $\operatorname{arcsinh}x=x+O(x^3)$ and $e^{-x}=1-x+O(x^2)$ give
$1-q=B\eps+O((B\eps)^2)$.  Substituting $B=dH$ gives the final expression.
\end{proof}

\begin{corollary}[Krein-positive expansion arrow]
\label{cor:krein-positive-expansion-arrow}
For the bridge normalization of \Cref{prop:bridge-displacement},
\[
        B>0 \quad\Longleftrightarrow\quad
        \chi_*>0 \quad\Longleftrightarrow\quad
        1-e^{-2\chi_*}>0 .
\]
Hence the expanding homogeneous branch is the Krein-positive bridge sheet.  The
formally contracting branch is a branch of the same local equations, but it is
not the positive bridge continuation unless a separate trans-renewal matching
principle is supplied.
\end{corollary}

\begin{proof}
The function $\operatorname{arcsinh}$ is strictly increasing and $e^{-2x}$ is
strictly decreasing, so the three inequalities are equivalent.  The last
sentence is the interpretation of the sign of the bridge norm in the signed
Krein sector.
\end{proof}

\begin{lemma}[Isotropy reduction]
\label{lem:isotropy-reduction}
For spatially homogeneous isotropic $(M,g)$ with unit normal $u$ and spatial
$h=g+u\otimes u$, any invariant symmetric $2$-tensor has perfect-fluid form
$S_{\mu\nu}=\rho u_\mu u_\nu+p h_{\mu\nu}$; if moreover $S_{\mu\nu}=fg_{\mu\nu}$ and
$\nabla^\mu S_{\mu\nu}=0$ then $f$ is constant.
\end{lemma}

\begin{proof}
The isotropy group $\mathrm{SO}(d)$ has no invariant spatial vectors, so mixed
components vanish and the spatial block is $\propto h$; the second claim is
$\nabla^\mu(fg_{\mu\nu})=\nabla_\nu f$.
\end{proof}

\begin{theorem}[Maximally symmetric pure-deficiency branch classification]
\label{thm:vacuum-trichotomy}
For the calibrated matter-free pure-deficiency sector with $B=dH$ and the
maximally symmetric vacuum ansatz, $\Lambda=(d-1)B^2/2d\ge0$: expanding de Sitter
for $B>0$, Minkowski for $B=0$, contracting de Sitter for $B<0$.  Anti-de Sitter
is excluded within this ansatz because $\Lambda\propto B^2$.
\end{theorem}

\begin{proof}
The vacuum identities give $\Lambda$; the sign of $B$ fixes $H=B/d$ and
distinguishes expansion from contraction while $\Lambda\ge0$.  A negative
$\Lambda$ would require leaving the real maximally symmetric ansatz.
\end{proof}

\begin{corollary}[Scope of the dark-energy sign result]
\label{cor:sign-dark-energy}
Within the real maximally symmetric calibrated pure-deficiency ansatz the
cosmological term is nonnegative.  This is a branch classification, not a no-go
theorem for negative effective energy in more general sectors.
\end{corollary}

\begin{proof}
Immediate from $\Lambda=(d-1)B^2/2d$.
\end{proof}

\subsection{Spatial flatness and higher-curvature protection}
\label{subsec:spatial-flatness}

The preceding branch classification used a maximally symmetric vacuum ansatz.
To determine rather than assume the curvature of the homogeneous slices, we now
retain the FLRW curvature index throughout the Hamiltonian constraint.  This
also tests whether the four-derivative renewal layer can regenerate spatial
curvature after the leading cancellation.

\begin{definition}[General-curvature pure-deficiency ansatz]
\label{def:general-curvature-flrw}
Let
\[
 ds^2=-N(t)^2dt^2+a(t)^2d\Sigma_k^2,\qquad
 k\in\{-1,0,+1\},\qquad H=\frac{\dot a}{Na},
\]
where $d\Sigma_k^2$ has constant sectional curvature $k$.  Write
$u:=k/a^2$.  At order $\eps^2$ let $\mathcal N^{\ren}_{\mu\nu}$ denote
the metric variation of the four-derivative part of
\Cref{def:effective-action}, moved to the right-hand side of the Einstein
equation, and define
\[
 \Mren:=-\frac{2}{d(d-1)}\mathcal N^{\ren}_{00}
\]
after setting $N=1$.
\end{definition}

\begin{lemma}[FLRW Hamiltonian constraint with renewal memory]
\label{lem:flrw-renewal-constraint}
For the ansatz of \Cref{def:general-curvature-flrw}, lapse variation of the
Einstein--Hilbert--renewal action gives
\[
 H^2+\frac{k}{a^2}
 =\frac{16\pi\Gren}{d(d-1)}\rho
  +\frac{2\Lambda}{d(d-1)}+\Mren+\mathcal R_{00},
\]
where $\mathcal R_{00}$ is the normalized controlled remainder.  On the
matter-free pure-deficiency branch, using $H=B/d$ and
$2\Lambda/[d(d-1)]=H^2$, this reduces to
\begin{equation}
\label{eq:flatness-master}
  u=\Mren+\mathcal R_{00}.
\end{equation}
In particular, at the leading Einstein order $\Mren=\mathcal R_{00}=0$,
so $k=0$.
\end{lemma}

\begin{proof}
For FLRW in $d+1$ dimensions,
$G_{00}=d(d-1)(H^2+k/a^2)/2$.  Taking the $00$ component of the field equation
and dividing by $d(d-1)/2$ gives the displayed constraint, equivalently obtained
by varying the lapse before fixing $N=1$.  The pure-deficiency identities cancel
the two $H^2$ terms and give \eqref{eq:flatness-master}.  Since $a>0$ and
$k\in\{-1,0,+1\}$, $u=0$ is equivalent to $k=0$.
\end{proof}

\begin{lemma}[Four-dimensional FLRW collapse of the quadratic-curvature sector]
\label{lem:flrw-quadratic-collapse}
In four spacetime dimensions, on every FLRW background,
\[
 \left.\frac{1}{\sqrt{-g}}\frac{\delta}{\delta g^{\mu\nu}}
 \int R_{\alpha\beta}R^{\alpha\beta}\sqrt{-g}\,d^4x
 \right|_{\rm FLRW}
 =\frac13
 \left.\frac{1}{\sqrt{-g}}\frac{\delta}{\delta g^{\mu\nu}}
 \int R^2\sqrt{-g}\,d^4x\right|_{\rm FLRW},
\]
and the metric variation of $\int C^2\sqrt{-g}\,d^4x$ vanishes.
Consequently FLRW probes only the combination $c_1+c_2/3$ of the three
quadratic-curvature coefficients.
\end{lemma}

\begin{proof}
In four dimensions,
\[
 R_{\mu\nu}R^{\mu\nu}=\frac13R^2+\frac12C^2-\frac12E_4,
 \qquad
 E_4=R_{\mu\nu\rho\sigma}^2-4R_{\mu\nu}^2+R^2.
\]
The integral of $E_4$ has vanishing local metric variation, while the variation
of $\int C^2$ is proportional to the Bach tensor.  Every FLRW metric is
conformally flat, hence its Weyl tensor and Bach tensor vanish.  Varying the
identity therefore gives the result; see
\cite{Stelle1977,Besse,Starobinsky1980} for the standard quadratic-gravity and
conformally flat reductions.
\end{proof}

\begin{theorem}[Renewal selection of spatial flatness]
\label{thm:renewal-flatness}
Assume the four-dimensional pure-deficiency branch, constant renewal rate
$\lambda$, and the local order-$\eps^2$ action of
\Cref{def:effective-action}.  Absorb the constant $c_5\lambda^2R$ term into
the renormalized Newton coefficient and suppose the dimensionless Wilson
combination remains uniformly perturbative on the branch.  Then
\begin{equation}
\label{eq:curvature-locked-memory}
 \Mren=16\pi\Gren\eps^2(6c_1+2c_2)
 u(2H^2-u),
\end{equation}
while $c_3$ drops out and $c_4(\nabla\lambda)^2=0$.  Hence, neglecting only
the separately displayed controlled remainder, the constraint factorizes as
\begin{equation}
\label{eq:renewal-flatness-factorization}
 u\,[1-\Phi(u)]=0,
 \qquad
 \Phi(u)=16\pi\Gren\eps^2(6c_1+2c_2)(2H^2-u).
\end{equation}
If $\sup|\Phi|<1$ throughout the validity window---in particular when the
Wilson coefficients are $O(1)$ in renewal units and
$\chi_{\ren}:=\eps^2|\Riem|\ll1$---then $k=0$ is the unique controlled
solution.  The algebraic alternative $\Phi=1$ lies at curvature of order the
inverse renewal scale and is outside the derivative expansion.  With a nonzero
remainder, the stable conclusion is
$|k|/a^2\le(1-\sup|\Phi|)^{-1}|\mathcal R_{00}|$; exact discreteness of
$k$ then again gives $k=0$ whenever this bound is smaller than $a^{-2}$.
\end{theorem}

\begin{proof}
By \Cref{lem:flrw-quadratic-collapse}, the curvature-squared action reduces on
FLRW, modulo the topological Euler term, to
$\eps^2(c_1+c_2/3)R^2$.  The $00$ equation of an $R+\gamma R^2$ theory can be
written
\[
 3(1+2\gamma R)(H^2+u)
 =\Lambda+\frac{\gamma}{2}R^2-6\gamma H\dot R.
\]
On the constant-$H$ branch, $R=12H^2+6u$ and $\dot u=-2Hu$, so
\[
 \frac16R^2-2H\dot R-2R(H^2+u)=6u(2H^2-u).
\]
Restoring the normalization of \Cref{def:effective-action} gives
\eqref{eq:curvature-locked-memory}; the coefficient is
$6(c_1+c_2/3)=6c_1+2c_2$.  The Weyl term has zero Bach response on FLRW,
$(\nabla\lambda)^2$ vanishes for constant $\lambda$, and
$\lambda^2R$ only renormalizes the Einstein coefficient.  Substitution in
\eqref{eq:flatness-master} gives \eqref{eq:renewal-flatness-factorization}.
If $|\Phi|<1$, the second factor cannot vanish and therefore $u=0$.  The
remainder estimate follows by taking absolute values in
$u(1-\Phi)=\mathcal R_{00}$.
\end{proof}

\begin{remark}[Scope and novelty of the flatness statement]
\label{rem:flatness-scope}
The theorem selects $k=0$ only for the homogeneous calibrated pure-deficiency
branch; it is not a theorem that every renewal spacetime or every matter-filled
cosmology is spatially flat.  Its nontrivial content is radiative stability:
in $D=4$ the complete local curvature-squared renewal response is locked to an
existing factor of $k/a^2$, does not shift the flat de Sitter radius, and cannot
create curvature inside the controlled window.  The FLRW and Gauss--Bonnet
identities are standard; the renewal-specific statement is that the
finite-window operator basis of \Cref{thm:basis} obeys this protection when
inserted into the deficiency branch.
\end{remark}

\begin{corollary}[Renewal flatness attractor]
\label{cor:renewal-flatness-attractor}
For a slowly varying expanding branch, define
$\Omega_k=-k/(a^2H^2)$.  Then
\[
 \frac{d}{dt}\log|\Omega_k|=-2H-2\frac{\dot H}{H}.
\]
Thus $|\Omega_k|$ decreases whenever
$1+\dot H/H^2>0$.  With $H=B/d$, this is
$1+d\dot B/B^2>0$; for constant $B>0$,
$|\Omega_k|\propto e^{-2Ht}=e^{-2N_{\rm e}}$.
\end{corollary}

\begin{proof}
Differentiate $\log|k|-2\log a-2\log H$ and use
$\dot a/a=H$ and $H=B/d$.
\end{proof}

\subsection{Spectral control and rolling-deficiency stability}
\label{subsec:gap-and-rolling-flatness}

The preceding theorem isolates the only four-derivative combination relevant
for FLRW, $6c_1+2c_2$.  We next separate a rigorous spectral estimate from the
additional microscopic matching needed to convert it into a Wilson-coefficient
bound.  This distinction is important: the gap estimate is a theorem, whereas
the correlator-to-action map remains a matching hypothesis until the renewal
sector is integrated out directly.

\begin{hypothesis}[Renewal induced-action matching]
\label{hyp:renewal-induced-matching}
On the reversible fluctuation subspace, let $\Lop$ be the positive
self-adjoint renewal generator and $\psi$ the fluctuation vector sourcing the
quartic response.  Assume that the curvature-squared coefficients factor as
\[
 c_i=A_{\ren}\,\widehat{\mathcal C}^{(i)}_4\,\mathcal T_2,
 \qquad
 \mathcal T_2:=
 \frac{2\langle\psi,\Lop^{-3}\psi\rangle}
      {\langle\psi,\Lop^{-1}\psi\rangle},
\]
where $A_{\ren}$ is the dimensionally normalized two-derivative matching
amplitude and $\widehat{\mathcal C}^{(i)}_4$ is the corresponding normalized
fourth-cumulant projection.  The normalization of $A_{\ren}$ is fixed by the
same microscopic convention used to define $\eps$ and $\lambda$.
\end{hypothesis}

This is the renewal analogue of induced-gravity and Green--Kubo matching:
local Wilson coefficients are moments of microscopic connected response
functions \cite{Sakharov1968,Adler1982,Zee1981,KipnisVaradhan}.  We do not use
that analogy to fix the normalization silently; it is recorded explicitly in
$A_{\ren}$.

\begin{theorem}[Spectral-gap bound for the renewal response moment]
\label{thm:renewal-gap-moment}
Suppose $\Lop\ge g\mathbf 1$ on the fluctuation subspace for some $g>0$.
Then
\[
 0\le\mathcal T_2\le\frac{2}{g^2}.
\]
Equality holds precisely when the spectral measure of $\psi$ is supported on
the gap eigenspace; equivalently, the normalized response is a single
exponential of rate $g$.
\end{theorem}

\begin{proof}
Functional calculus gives $\Lop^{-2}\le g^{-2}\mathbf 1$.  Conjugating by
$\Lop^{-1/2}$ yields $\Lop^{-3}\le g^{-2}\Lop^{-1}$.  Taking the quadratic
form in $\psi$ and dividing by the positive denominator proves the estimate.
Equality in the operator inequality on $\psi$ requires its spectral measure to
be supported at $g$.
\end{proof}

\begin{corollary}[Conditional microscopic sufficient condition for flatness]
\label{cor:microscopic-flatness-bound}
Under \Cref{hyp:renewal-induced-matching,thm:renewal-gap-moment}, set
$\widehat{\mathcal C}^{\eff}_4:=6\widehat{\mathcal C}^{(R)}_4+
2\widehat{\mathcal C}^{(S)}_4$.  Then
\[
 16\pi\Gren|6c_1+2c_2|
 \le \frac{32\pi\Gren|A_{\ren}|}{g^2}
       |\widehat{\mathcal C}^{\eff}_4|.
\]
Consequently the perturbative premise of \Cref{thm:renewal-flatness} follows
from
\[
 \eps^2\frac{32\pi\Gren|A_{\ren}|}{g^2}
 |\widehat{\mathcal C}^{\eff}_4|
 \left(2H^2+\left|\frac{k}{a^2}\right|\right)<1.
\]
For a single-exponential Poisson reset channel the temporal estimate is
saturated with $g$ equal to the relaxation rate.
\end{corollary}

\begin{proof}
Insert the matching hypothesis into $6c_1+2c_2$, apply
\Cref{thm:renewal-gap-moment}, and then use the sufficient estimate for
$|\Phi|$ in \Cref{thm:renewal-flatness}.
\end{proof}

\begin{remark}[What the gap theorem does and does not remove]
\label{rem:gap-matching-scope}
The spectral theorem removes uncertainty in the temporal response moment once
the generator and its gap are known.  It does not by itself derive
\Cref{hyp:renewal-induced-matching}, fix $A_{\ren}$, or prove that a general
irreversible transfer generator obeys the same sharp constant.  A direct
cumulant expansion of the microscopic reset determinant is therefore still
needed before the Wilson-coefficient assumption can be said to have disappeared
completely.
\end{remark}

The constant-deficiency solution is also stable under a controlled roll.  Since
$k$ is a fixed topological label, the dynamical quantity is
$\Omega_k:=-k/(a^2H^2)$ rather than $k$ itself.

\begin{theorem}[Global dilution of curvature under sustained renewal acceleration]
\label{thm:global-renewal-flattening}
For every expanding FLRW solution,
\[
 \frac{d\Omega_k}{dN}
 =-2\left(1+\frac{\dot H}{H^2}\right)\Omega_k,
 \qquad N=\log a.
\]
If $\epsilon_H:=-\dot H/H^2\le\epsilon_*<1$ for all
$N\ge N_0$, then
\[
 |\Omega_k(N)|\le |\Omega_k(N_0)|
 e^{-2(1-\epsilon_*)(N-N_0)}.
\]
On a rolling deficiency branch $H=B/d+O(\eps^2)$, the leading attraction
condition is $1+d\dot B/B^2>0$.  Local diffeomorphism-invariant renewal
corrections preserve the locus $k=0$, although their $k$-independent part may
modify the on-locus evolution of $H$.
\end{theorem}

\begin{proof}
Direct differentiation of $\Omega_k=-k(aH)^{-2}$ gives the first identity.
The assumed bound implies
$d\log|\Omega_k|/dN\le-2(1-\epsilon_*)$, and integration gives the estimate.
Substitution of $H=B/d$ gives the stated leading condition.  Finally, on FLRW
every local scalar formed from the metric, homogeneous renewal fields, and
finitely many covariant derivatives is a function of
$H,\dot H,\ldots$ and $u=k/a^2$; subtracting its value at $u=0$ leaves a
factor of $u$.  Thus local corrections cannot create a curvature term on the
$k=0$ locus, although the $u^0$ part can dress the background equation.
\end{proof}

\begin{remark}[Effective deficiency equation of state]
\label{rem:deficiency-effective-eos}
The relation between a rolling budget and an equation of state depends on the
stress-energy register.  For a separately conserved effective deficiency fluid
one obtains \(w_B=-1-2\dot B/(dHB)\), reducing to
\(w_B=-1-2\dot B/B^2\) on the deficiency-dominated branch \(H=B/d\).
With matter--deficiency exchange, the physical pressure contains an additional
exchange term, whereas a noninteracting observational reconstruction still
uses the logarithmic derivative of \(\rho_B\).  These distinctions and their
crossing consequences are proved in \Cref{sec:dark-energy}.
\end{remark}

\medskip
\noindent\textbf{Transition to dynamical deficiency.}
The results above concern the selected constant-deficiency de~Sitter endpoint and
its curvature stability.  A time-dependent homogeneous deficiency requires a
separate stress-energy register and, in the genuinely interacting case, an
explicit exchange current.  The microscopic rate function, the critical
fractional relaxation law, the reversible no-crossing theorem, and the
irreversible crossing threshold are therefore developed together in the
dedicated dark-energy analysis of \Cref{sec:dark-energy}.  This separation
prevents a rolling effective fluid from being confused with the vacuum term
\(\Lambda(B)g_{\mu\nu}\), which is forced to be constant by the contracted
Bianchi identity.

\subsection{Global de Sitter completion}
\label{subsec:global-desitter-completion}

The preceding classification is local: it identifies a positive
constant-curvature geometry on each calibrated patch.  Passing from such a
local statement to complete global de Sitter space requires control of topology,
maximality, and quotient identifications.  We make that additional input
explicit rather than incorporating it silently into the local field equation.

Constant $H$ in the flat FLRW chart identifies only the expanding planar patch; a
global statement needs a continuation rule.  The foundational construction allows
$\Ztwo$ sign holonomy in the Krein--Clifford fibre and constructs no completion, so
we distinguish fibre holonomy from base topology.

\begin{principle}[Universal predictive completion]
\label{prin:universal-predictive-completion}
The physical continuum spacetime of a connected calibrated renewal phase is the
geodesically complete universal Lorentzian development of its local metric.
Internal sign holonomies classified by $H^1(G,\Ztwo)$ act on the Krein--Clifford
fibre and do not induce quotient identifications of the spacetime base.
\end{principle}

\begin{remark}[Fibre holonomy versus spacetime topology]
\label{rem:fibre-versus-base-holonomy}
The signed predictive cover may carry nontrivial $\Ztwo$ holonomy, classifying the
fundamental symmetry of the Krein fibre, not identifying spacetime points;
\Cref{prin:universal-predictive-completion} unwraps the Lorentzian base while
retaining the holonomy in the internal bundle.  This is a global postulate, not a
consequence of the local Einstein equation.
\end{remark}

\begin{lemma}[Constant-curvature tensor of the pure-deficiency vacuum]
\label{lem:pure-deficiency-constant-curvature}
For the homogeneous isotropic matter-free pure-deficiency sector with $H=B/d$,
$B>0$ and $\Lambda=d(d-1)H^2/2$, the local spacetime has constant sectional
curvature $H^2$:
$R_{\mu\nu\rho\sigma}=H^2(g_{\mu\rho}g_{\nu\sigma}-g_{\mu\sigma}g_{\nu\rho})$,
hence $R_{\mu\nu}=dH^2g_{\mu\nu}$, $R=d(d+1)H^2$.
\end{lemma}

\begin{proof}
Homogeneity and isotropy give a conformally flat FLRW metric with vanishing Weyl
tensor; the vacuum equation gives $R_{\mu\nu}=dH^2g_{\mu\nu}$, and substituting both
into the Riemann decomposition gives the constant-curvature identity.
\end{proof}

\begin{proposition}[The renewal flat slicing is a planar de Sitter chart]
\label{prop:planar-desitter-chart}
$ds^2=-dt^2+e^{2Ht}\delta_{ij}dx^idx^j$ is isometric to the expanding planar region
$X_0+X_{d+1}>0$ of the de Sitter hyperboloid
$-X_0^2+\sum_{i=1}^{d+1}X_i^2=H^{-2}$ via
$X_0=H^{-1}\sinh(Ht)+\tfrac H2e^{Ht}|x|^2$, $X_i=e^{Ht}x_i$,
$X_{d+1}=H^{-1}\cosh(Ht)-\tfrac H2e^{Ht}|x|^2$; the flat slicing is not
geodesically complete.
\end{proposition}

\begin{proof}
Direct substitution gives the hyperboloid equation and the metric, and
$X_0+X_{d+1}=H^{-1}e^{Ht}>0$ identifies the planar patch.
\end{proof}

\begin{lemma}[Positive-curvature Lorentzian space-form rigidity]
\label{lem:desitter-spaceform-rigidity}
A connected, simply connected, geodesically complete $(M,g)$ of dimension
$d+1\ge3$ with
$R_{\mu\nu\rho\sigma}=H^2(g_{\mu\rho}g_{\nu\sigma}-g_{\mu\sigma}g_{\nu\rho})$,
$H>0$, is globally isometric to complete de Sitter space $\dS_{d+1}(H)$.
\end{lemma}

\begin{proof}
This is Lorentzian space-form rigidity \cite{ONeill,Wolf}.  The curvature identity
makes $(M,g)$ a Lorentzian space form of curvature $H^2$; Cartan--Ambrose--Hicks
extends a tangent-space isometry along geodesics to a local isometry, simple
connectedness makes it path-independent, and geodesic completeness makes the
developing map a covering of the complete model.  Since $\dS_{d+1}(H)$ is simply
connected for $d\ge2$, the covering is one-sheeted.
\end{proof}

\begin{theorem}[Complete global de Sitter renewal branch]
\label{thm:complete-global-desitter}
Assume the calibrated homogeneous pure-deficiency sector with $B>0$, the local
renewal-Einstein equation, and \Cref{prin:universal-predictive-completion}.  Then
the maximal physical renewal spacetime is globally isometric to complete de Sitter
space $(M_{\ren},g_{\ren})\cong\dS_{d+1}(H)$, $H=B/d$,
$\Lambda=(d-1)B^2/2d$; the flat FLRW metric is its expanding planar chart, the
static-patch horizon has $T_{\GH}=H/2\pi=B/2\pi d$, and for primitive
matrix-product screens the regulated horizon entropy tends to
$S_{\hor}=A_{\hor}/4\Gren$.
\end{theorem}

\begin{proof}
By \Cref{lem:pure-deficiency-constant-curvature} the local metric has constant
curvature $H^2$; by \Cref{prin:universal-predictive-completion} the physical
spacetime is its complete simply connected universal development, identified with
$\dS_{d+1}(H)$ by \Cref{lem:desitter-spaceform-rigidity}.  In static coordinates the
horizon at $r=H^{-1}$ has surface gravity $H$, giving $T_{\GH}=H/2\pi$; the entropy
follows from \Cref{thm:renewal-screen-area-law} and $\Gren=(4\etaScr)^{-1}$.
\end{proof}

\begin{corollary}[Exclusion of nontrivial de Sitter quotients]
\label{cor:exclude-desitter-quotients}
The maximal physical renewal spacetime is not a nontrivial quotient
$\dS_{d+1}(H)/\Gamma$.
\end{corollary}

\begin{proof}
The universal development is simply connected; a nontrivial freely acting discrete
quotient has nontrivial fundamental group.
\end{proof}

\begin{remark}[Status of the global conclusion]
\label{rem:global-desitter-status}
Local de Sitter curvature, the planar embedding, the Hubble scale, and $\Lambda$
are theorems of the leading field equation.  Identification of the physical
maximal continuation with the complete space form uses
\Cref{prin:universal-predictive-completion} and is a separate, principle-conditional
theorem; the local result is the stronger, intrinsic statement.
\end{remark}

\subsection{Perturbative stability and no-hair}

After identifying the exact homogeneous branch, it is natural to ask whether it
is robust under small calibrated perturbations.  The statement below is an
energy-stability result within the specified closure class; it is logically
separate from both the local de Sitter theorem and the global completion
principle.

\begin{hypothesis}[Small calibrated data]
\label{hyp:small-data}
Let $(g,K,\Tren)=(g_{\dS},K_{\dS},0)+(\delta g,\delta K,\delta T)$ be small in
$H^s$, $s>d/2+2$, satisfying the renewal constraints with residual $O(\eps)$, with
the matter model obeying the dominant energy condition and renewal damping
compatible with the deficiency budget.
\end{hypothesis}

\begin{theorem}[Small-data closure and renewal no-hair]
\label{thm:nohair}
Under \Cref{hyp:small-data} the renewal-Einstein system closes locally in time and
$\|g(t)-g_{\dS}(t)\|_{H^s}+\|\Tren(t)\|_{H^{s-1}}\le Ce^{-\gamma t}(\|\delta g\|_{H^s}
+\|\delta T\|_{H^{s-1}})+O(\eps)$ for some $\gamma>0$ set by $H=B/d$.
\end{theorem}

\begin{proof}
Linearizing \Cref{thm:einstein} around de Sitter gives a hyperbolic system with
positive cosmological damping; the energy estimate
$\dot E_s\le-2\gamma E_s+CE_s^{3/2}+C\eps^2$ absorbs the nonlinear term for small
data, giving exponential decay to an $O(\eps)$ neighbourhood; constraint
propagation follows from the Bianchi identity and conservation of $\Tren$.
\end{proof}

\begin{remark}
The theorem is perturbative: the de Sitter branch is dynamically coherent under
small calibrated perturbations, not future-stable for arbitrary data.
\end{remark}

\subsection{Horizon thermodynamics}

The de Sitter solution also provides a first test of the horizon thermodynamics
derived in the previous section.  Because only stationarity and local modular
calibration are required, the same argument applies to any nonextremal
stationary renewal horizon.  We record the resulting temperature and entropy
laws and indicate, without developing it here, the interface with the separate
black-hole analysis.

\begin{corollary}[Stationary renewal horizons]
\label{cor:horizons}
For a stationary solution with nondegenerate Killing horizon $\mathcal H$, each cut
$\Sigma$ carries renewal entropy $S_{\hor}(\Sigma)=\Area(\Sigma)/4G+O(\eps^2)$ and
local KMS temperature $T_{\mathcal H}=\kappa_{\mathcal H}/2\pi$.
\end{corollary}

\begin{proof}
The metric is locally Rindler to first order near $\mathcal H$, so calibration
applies to the local wedge, giving $T=\kappa/2\pi$ and the area entropy, the first
correction belonging to the finite-resolution layer of \Cref{sec:qg}.
\end{proof}

\begin{corollary}[Two-sided horizon bridge interface]
\label{cor:two-sided}
In a stationary doubled renewal horizon sector the renewal bridge criterion applies
to the bifurcation screen: the two-sided horizon is geometrically connected exactly
when the bifurcation screen is balanced and self-averaging and the doubled state has
nonzero deck-odd Krein correlation, in which case
$I_{\rm pred}(L:R)=2S_{\hor}=2S_{\bridge}$ at leading order.
\end{corollary}

\begin{proof}
The stationary horizon supplies a calibrated bifurcation screen (\Cref{cor:horizons});
the predictive-screen bridge criterion of \cite{PelissierBridge} gives the
equivalence, and the doubled-horizon normalization gives the entropy equality.
\end{proof}

\begin{remark}[Boundary with the black-hole paper]
This identifies the stationary horizon interface; it does not derive collapse,
evaporation, radiation entropy, Page transitions, trans-renewal interiors, or
endpoint cores, which require a dynamical radiation state and belong to a separate
paper \cite{PelissierBlackHoles}.
\end{remark}

% =====================================================================
\section{Finite-resolution effective action and Wilson coefficients}
\label{sec:qg}
% =====================================================================

The Einstein equation is the leading term of the renewal asymptotic expansion,
not its endpoint.  The same finite-window expansion that produced the Gaussian
Ricci response has a next nontrivial order governed by fourth cumulants and
slow spatial variation.  This section identifies the resulting local operator
basis, separates universal symmetry constraints from channel-dependent Wilson
data, and determines the regime in which the truncation remains controlled.

We record the first correction above the leading Einstein limit as a
finite-resolution effective-action layer; no all-orders ultraviolet completion
is asserted.  The \emph{Wilson basis} is
fixed unconditionally by locality, the renewal field content, and the
four-dimensional Gauss--Bonnet identity.  For a general microscopic channel
the coefficients are functionals of its normalized fourth and mixed cumulants.
For the specified isotropic compound-Poisson volume channel, however,
\Cref{thm:triple} evaluates the curvature-squared ray explicitly:
\[
 (\kR,\kS,\kC)
 =\Gren\lambda^2\frac{Q}{D^2}
   \left(\frac53,\frac53,-\frac14\right),
 \qquad
 (c_1,c_2,c_3)
 =\Gren\lambda^2\frac{Q}{D^2}
   \left(\frac54,\frac53,-\frac14\right).
\]
This worked-channel statement does not fix the coefficients of every renewal
channel.

\subsection{Derivative expansion and four-derivative operator basis}
\label{subsec:qg-effective}

Before writing the correction explicitly, we clarify its logical status.  The
$\eps^2$ term is a Wilsonian finite-resolution contribution: covariance,
parity, and derivative order determine its admissible operator content, whereas
its coefficients retain information about the underlying renewal statistics.

\begin{proposition}[Effective status of the \(\eps^2\) layer]
\label{prop:effective-qg}
The order-$\eps^2$ correction is effective in the sense that \textup{(i)} it is the
first correction beyond the large-deviation/Gaussian Einstein saddle; \textup{(ii)}
it is local at continuum resolution and organized by a finite higher-curvature
basis; \textup{(iii)} its coefficients remember the substrate through the gap and
diamond cumulants; \textup{(iv)} it defines no all-orders ultraviolet completion,
graviton Fock space, or fundamental scattering theory.
\end{proposition}

\begin{proof}
$\eps$ measures the renewal micro-scale relative to the diamond; the leading
large-deviation term gives the proper-time action, the Gaussian prefactor the Ricci
response and Einstein saddle, and the next local finite-window contribution appears
at $\eps^2$ in the finite curvature/deficiency basis.  Stopping at this order, with
no resummation, gives the effective status.
\end{proof}

\subsubsection{Four-derivative effective action}

We now pass from this effective interpretation to the local action itself.  The
field content consists of the metric and the renewal rate, and the relevant
operators are the parity-even invariants containing four derivatives in total.
The Gauss--Bonnet identity and integration by parts reduce the apparently larger
list to a finite dynamical quotient.

The renewal partition function on a local diamond has the saddle expansion
$\log Z_\eps=-\eps^{-1}S_{\rm cl}+\tfrac12\log\det(\operatorname{Hess}S_{\rm cl})
+\eps\mathcal A_1+O(\eps^2)$; the leading piece gives Einstein--Hilbert and
cosmological terms, the next local covariant correction appears at $\eps^2$.

\begin{definition}[Renewal effective action]
\label{def:effective-action}
$S_{\eff}[g,\lambda]=\tfrac1{16\pi G}\int(R-2\Lambda)\sqrt{-g}\,d^{d+1}x
+\eps^2\int\Lcal_2(g,\lambda)\sqrt{-g}\,d^{d+1}x+O(\eps^3)$, with
$\Lcal_2=c_1R^2+c_2R_{\mu\nu}R^{\mu\nu}+c_3C_{\mu\nu\rho\sigma}C^{\mu\nu\rho\sigma}
+c_4(\nabla\lambda)^2+c_5\lambda^2R$.  The five operators are the unconditional
basis; the $c_i$ are conditional data fixed by the renewal cumulants.
\end{definition}

\begin{hypothesis}[Four-derivative sector and field content]
\label{hyp:four-derivative-sector}
The order-$\eps^2$ correction is evaluated in four dimensions for a torsion-free
metric and a parity-even scalar $\lambda$.  We keep operators of at most four
derivatives, impose $\lambda\mapsto-\lambda$, use integration by parts to remove
$\lambda\Box\lambda$ in favour of $(\nabla\lambda)^2$, assign $\lambda^4$ to the
deficiency potential, and treat $\Box R$, $\Box(\lambda^2)$ as boundary terms.
Matching is off shell in $(g,\lambda)$; the leading Einstein equation is not used.
\end{hypothesis}

\begin{theorem}[Four-dimensional channel basis and dynamical quotient]
\label{thm:basis}
Under \Cref{hyp:four-derivative-sector}, the parity-even bookkeeping basis is
$\{R^2,R_{\mu\nu}R^{\mu\nu},C^2,(\nabla\lambda)^2,\lambda^2R\}$.  With the Euler
density $E_4=R_{\mu\nu\rho\sigma}^2-4R_{\mu\nu}^2+R^2$ one has
$C^2=E_4+2R_{\mu\nu}^2-\tfrac23R^2$, so
$c_1R^2+c_2R_{\mu\nu}^2+c_3C^2=\alpha R^2+\beta R_{\mu\nu}^2+c_3E_4$ with
$\alpha=c_1-\tfrac23c_3$, $\beta=c_2+2c_3$.  The local bulk equations depend on
$(\alpha,\beta)$; this classification of the basis is unconditional, only the $c_i$
depending on the cumulant data.
\end{theorem}

\begin{proof}
For a torsion-free metric every parity-even quadratic curvature scalar is a
combination of $R_{\mu\nu\rho\sigma}^2$, $R_{\mu\nu}^2$, $R^2$ up to a divergence
\cite{Stelle1977}.  In four dimensions
$R_{\mu\nu\rho\sigma}^2=C^2+2R_{\mu\nu}^2-\tfrac13R^2$; substituting in $E_4$ gives
$E_4=C^2-2R_{\mu\nu}^2+\tfrac23R^2$, equivalent to the displayed $C^2$ identity.  The
integral of $E_4$ is topological with vanishing local variation.  The scalar-field
restrictions leave $(\nabla\lambda)^2$ and $\lambda^2R$.  No leading field equation
is used.
\end{proof}

\begin{remark}[Why field redefinitions are not used in the matching]
\label{rem:field-redefinitions}
A redefinition $g_{\mu\nu}\mapsto g_{\mu\nu}+\eps^2(aR_{\mu\nu}+bRg_{\mu\nu})$ shifts
the action by operators proportional to the leading Euler--Lagrange tensor; on shell
it can move strength between $R^2$ and $R_{\mu\nu}^2$ and remove both up to $E_4$
\cite{Stelle1977,Donoghue1994}.  Here matching is off shell and includes the renewal
stress ledger and the deficiency field, so the same redefinition would transfer
terms into matter and $\lambda$ operators; we quote coefficients in a fixed
parametrization and display the Gauss--Bonnet-invariant $(\alpha,\beta)$.
\end{remark}

\subsection{Beyond four derivatives}
\label{subsec:six-derivative-radius-shift}

The four-derivative protection of the de Sitter radius is special to four
spacetime dimensions.  Curvature locking itself is more robust: for every
local diffeomorphism-invariant finite-derivative term, the part of its FLRW
constraint that depends on $k$ is divisible by $u=k/a^2$.  At six derivatives,
however, the $u^0$ part need not vanish.

\begin{proposition}[First local radius shift]
\label{prop:six-derivative-radius-shift}
Let the order-$\eps^4$ constant-curvature reduction of the local effective
action contain
\[
 \eps^4 b_{\dS}R^3,
\]
where $b_{\dS}$ is the linear combination of cubic-curvature Wilson
coefficients evaluated on a maximally symmetric background.  Then the flat
constant-curvature equation is
\[
 R-4\Lambda=16\pi\Gren\eps^4b_{\dS}R^3,
\]
and its root continuously connected to Einstein gravity satisfies
\[
 \frac{\delta R}{4\Lambda}
 =256\pi\Gren\eps^4b_{\dS}\Lambda^2+O(\eps^8).
\]
Thus the sign of the leading radius shift is the sign of $b_{\dS}$.  The
$k$-dependent part remains curvature-locked, so this correction relocates the
de Sitter point without sourcing spatial curvature.
\end{proposition}

\begin{proof}
For
$f(R)=(R-2\Lambda)/(16\pi\Gren)+\eps^4b_{\dS}R^3$, a four-dimensional
constant-curvature solution obeys $Rf'(R)-2f(R)=0$.  This gives the displayed
cubic equation.  Writing $R=4\Lambda+\delta R$ and expanding to first order in
$\eps^4$ yields
$\delta R=16\pi\Gren\eps^4b_{\dS}(4\Lambda)^3+O(\eps^8)$.
The locking statement follows from the FLRW decomposition into a
$u$-independent part plus a polynomial whose difference from its value at
$u=0$ is divisible by $u$.
\end{proof}

\begin{remark}[Nonlocal memory is a genuine boundary of the theorem]
\label{rem:nonlocal-curvature-boundary}
The preceding divisibility argument is local.  For covariant form factors
$F(\Box)$ that are regular on the relevant de Sitter spectral modes, each
homogeneous eigenmode is preserved and no new $a^{-2}$ component is generated
from a constant source.  This does \emph{not} hold for an arbitrary
explicitly time-dependent causal kernel: such a kernel can mix temporal modes.
Accordingly, extension of curvature locking to the full causal nonlocal renewal
response requires a separate spectral theorem for the admissible kernels.  A
$k$-independent effective component proportional to $a^{-2}$ (equivalently a
curvature-mimicking $w=-1/3$ sector) is the sharp obstruction to test for, not
something excluded by locality arguments alone.
\end{remark}

\subsection{Cumulant expansion and heat-kernel matching}

The operator basis is universal within the stated symmetry class, but its
coefficients are not.  They are determined by the normalized fourth and mixed
cumulants of the microscopic renewal channel.  We therefore separate the
basis theorem from the model-dependent matching and display explicitly where
the present numerical information enters.

\subsubsection{Covariant cumulant expansion and heat-kernel benchmark}
\label{subsec:reset-action-heat-kernel}

The normalized cumulants below can enter a bulk action only after specifying how
the reset generator is coupled to the curved continuum field.  A useful formal
starting point is the covariant Kramers--Moyal expansion.  In Riemann normal
coordinates, for a local jump kernel $W(x'|x)$ with displacement $\xi$,
\[
 \mathcal R f(x)=\sum_{n\ge1}\frac{1}{n!}
 M_n^{a_1\cdots a_n}(x)\nabla_{a_1}\cdots\nabla_{a_n}f(x),
 \qquad
 M_n^{a_1\cdots a_n}=\int \xi^{a_1}\cdots\xi^{a_n}W(x'|x)\,dx'.
\]
For a reflection-symmetric isotropic kernel,
$M_2^{ab}=Dg^{ab}$ and
$M_4^{abcd}=Q(g^{ab}g^{cd}+g^{ac}g^{bd}+g^{ad}g^{bc})$, so the principal
scalar part is
\[
       \mathcal R=\frac D2\Box+\frac Q8\Box^2+\cdots .
\]
The factor $1/8$ is the exact Kramers--Moyal contraction factor.  Commutators of
covariant derivatives then generate curvature invariants, but the resulting
bulk coefficients depend on the complete operator, measure, boundary
conditions, and subtraction convention.  Thus this expansion motivates
\Cref{hyp:renewal-induced-matching}; it does not by itself fix $A_{\ren}$.

\begin{proposition}[Minimal scalar heat-kernel benchmark]
\label{prop:minimal-scalar-heat-kernel}
For the minimal scalar Laplacian in four dimensions, modulo a total derivative,
the second local heat-kernel coefficient is
\[
 a_2=\frac{13}{1080}R^2+\frac1{180}R_{\mu\nu}R^{\mu\nu}
       +\frac1{180}C_{\mu\nu\rho\sigma}C^{\mu\nu\rho\sigma}.
\]
On FLRW, where the Bach variation vanishes and the
$R_{\mu\nu}R^{\mu\nu}$ variation equals one third of the $R^2$ variation,
this provides a universal \emph{minimal-Laplacian benchmark} for the relative
curvature channels.  It is not automatically the coefficient of the
non-minimal renewal operator $\partial_\tau+\mathcal R$.
\end{proposition}

\begin{proof}
For a minimal scalar Laplacian,
$a_2=(5R^2-2R_{\mu\nu}^2+2R_{\mu\nu\rho\sigma}^2)/360$ up to the divergence
$\Box R$.  Substituting
$R_{\mu\nu\rho\sigma}^2=C^2+2R_{\mu\nu}^2-R^2/3$ gives the displayed basis.
The FLRW reduction follows from conformal flatness and the four-dimensional
Gauss--Bonnet identity.
\end{proof}

\begin{remark}[Natural renewal correlation length]
\label{rem:renewal-correlation-length}
If the quadratic spatial generator has diffusion scale $D$ and the temporal
relaxation generator has gap $g$, the natural correlation length is
\[
             \xi_{\ren}^2:=\frac{D}{g}.
\]
For the compound-Poisson velocity-jump channel with $D=rs^2$ and $g=r$, this
gives $\xi_{\ren}=s$.  This fixes a natural candidate for the derivative
expansion scale, but identifying it with the precise proper-time cutoff in the
renormalized determinant is part of the matching problem.
\end{remark}

\begin{remark}[Six-derivative baseline and non-Gaussian excess]
\label{rem:six-derivative-baseline}
At six derivatives the coefficient $b_{\dS}$ generally contains two logically
distinct pieces: a universal contribution generated by repeated propagation of
the quadratic operator and a channel-dependent excess involving higher connected
jump cumulants.  Therefore a Gaussian reset need not make the full cubic
curvature coefficient vanish.  Only the higher-cumulant excess vanishes in a
strict Gaussian model.  This corrects the stronger, unjustified inference
``$b_{\dS}=0$ iff Gaussian.''
\end{remark}

\begin{definition}[Normalized quartic renewal cumulant]
\label{def:normalized-quartic-cumulant}
With $\langle X,Y\rangle=X_{\mu\nu\rho\sigma}Y^{\mu\nu\rho\sigma}$ on algebraic
curvature tensors and the orthogonal decomposition
$X=C[X]+\tfrac12(g\mathbin{\wedge}S[X])+\tfrac{R[X]}{12}(g\mathbin{\wedge}g)$,
$S_{\mu\nu}[X]=\Ric_{\mu\nu}[X]-\tfrac14R[X]g_{\mu\nu}$ ($\mathbin{\wedge}$ the
Kulkarni--Nomizu product), let $\delta\mathcal J$ be the centered calibrated renewal
fluctuation with cumulants $\mathcal C_2,\mathcal C_4$ and standardized
$\widehat{\mathcal C}_4=\mathcal C_2^{-1/2\otimes4}\mathcal C_4\mathcal C_2^{-1/2\otimes4}$.
Its curvature response
$\mathfrak Q_4(X)=\tfrac1{4!}\widehat{\mathcal C}_{4ABCD}\mathcal V^A(X)\cdots\mathcal V^D(X)$
is, by isotropy and parity,
$\mathfrak Q_4(X)=G\lambda^2(\kappa_R R[X]^2+\kappa_S S_{\mu\nu}[X]S^{\mu\nu}[X]
+\kappa_C C_{\mu\nu\rho\sigma}[X]C^{\mu\nu\rho\sigma}[X])$, with
$(\kappa_R,\kappa_S,\kappa_C)$ the normalized scalar, traceless-Ricci, and Weyl
cumulant coordinates.
\end{definition}

\begin{remark}[Normalization invariance]
\label{rem:cumulant-normalization-invariance}
Under $\delta\mathcal J\mapsto a\delta\mathcal J$ the three $\kappa$'s are unchanged,
as under orthogonal changes of the Gaussian-normalized coordinates; a nonorthogonal
field redefinition changes the quoted Wilson coefficients in the usual way.
\end{remark}


\begin{proposition}[Second and fourth cumulants of the worked reset stream]
\label{prop:worked-channel-cumulants}
For a compound-Poisson reset stream of rate $r$ with centered Gaussian step law
of covariance $\Sigma=s^2I_d$, the tilted cumulant generating function is
\[
 \Lambda(\theta)=r\left(e^{\frac12s^2|\theta|^2}-1\right).
\]
Its second and fourth cumulants are
\[
 K^{(2)}_{ij}=rs^2\delta_{ij},\qquad
 K^{(4)}_{ijkl}=rs^4
 (\delta_{ij}\delta_{kl}+\delta_{ik}\delta_{jl}+\delta_{il}\delta_{jk}).
\]
Hence the Gaussian propagation scale is $D=rs^2$, whereas the first isotropic
quartic scale is the independent quantity $Q=rs^4$.  The fully symmetric
$O(d)$-invariant quartic sector is one-dimensional; anisotropic traceless and
Clifford-odd channels vanish for this model.
\end{proposition}

\begin{proof}
For a compound-Poisson process, derivatives of $\Lambda$ at the origin are the
rate times the raw moments of one jump.  Isserlis' formula gives
$\mathbb E[\xi_i\xi_j\xi_k\xi_l]=s^4(\delta_{ij}\delta_{kl}+
\delta_{ik}\delta_{jl}+\delta_{il}\delta_{jk})$.  The invariant-tensor statement
follows because the fully symmetric $O(d)$-invariant fourth-order tensors are
spanned by this combination.
\end{proof}



\begin{definition}[Normalized Lorentzian diamond quartic response]
\label{def:diamond-quartic-response}
Let $D_\tau(p,u)$ be the normalized causal diamond centered at $p$ with unit
timelike axis $u$.  In diamond coordinates
$\eta=\frac{\tau}{2}x$, $x\in[-1,1]$, and
$z=\frac{\tau}{2}(1-|x|)\xi$, use the profile
\[
 \rho_{m,n}(x,\xi)\propto
 (1-|x|)^m(1-|\xi|^2)^n .
\]
Setting $\tau/2=1$, its even mixed moments factor as
\[
 \langle\eta^{2a}|z|^{2b}\rangle
 =
 \frac{B(2a+1,m+d+2b+1)}{B(1,m+d+1)}
 \frac{B(b+\frac d2,n+1)}{B(\frac d2,n+1)} .
\]
The curvature correction is computed with the normalized Lorentzian volume
factor
$\sqrt{-g}=1-\frac16R_{\mu\nu}x^\mu x^\nu+O(|x|^3)$.
\end{definition}

\begin{lemma}[Canonical diamond moments]
\label{lem:canonical-diamond-moments}
For $d=3$ and the canonical profile $(m,n)=(0,0)$,
\[
 \langle\eta^2\rangle=\frac1{15},\qquad
 \langle|z|^2\rangle=\frac25,\qquad
 \langle|z|^4\rangle=\frac3{14},\qquad
 \langle|z|^6\rangle=\frac2{15},\qquad
 \langle\eta^2|z|^4\rangle=\frac1{210}.
\]
In particular,
$\langle\eta^2|z|^4\rangle\ne
 \langle\eta^2\rangle\langle|z|^4\rangle$:
the shrinking transverse radius produces a genuine longitudinal--transverse
correlation absent from a product cylinder or a purely spatial ball.
\end{lemma}

\begin{proof}
These values follow by direct evaluation of the beta-function moments in
\Cref{def:diamond-quartic-response}.
\end{proof}

\begin{theorem}[Linear-curvature causal-diamond quartic-moment response]
\label{thm:diamond-linear-response}
For the isotropic fourth cumulant
\[
 K^{(4)}_{ijkl}
 =Q(\delta_{ij}\delta_{kl}+\delta_{ik}\delta_{jl}
      +\delta_{il}\delta_{jk}),\qquad Q=rs^4,
\]
the normalized diamond response can depend only on the two scalar contractions
selected by the diamond axis, $R$ and
$R_{uu}:=R_{\mu\nu}u^\mu u^\nu$.  Its coefficients are
\[
 \alpha_u=-\frac16
 \frac{\langle\eta^2|z|^4\rangle
 -\langle\eta^2\rangle\langle|z|^4\rangle}
 {\langle|z|^4\rangle},
 \qquad
 \alpha_\perp=-\frac16
 \frac{\langle|z|^6\rangle
 -\langle|z|^2\rangle\langle|z|^4\rangle}
 {\langle|z|^4\rangle},
\]
with
\[
 \alpha_R=\alpha_\perp,\qquad
 \alpha_{R_{uu}}=\alpha_u+\alpha_\perp .
\]
For the canonical profile $(0,0)$ in $d=3$,
\[
 \boxed{\alpha_R=-\frac1{27},\qquad
 \alpha_{R_{uu}}=-\frac4{135}.}
\]
Thus the order-$\eps^2$ diamond observable contributed by the worked isotropic
reset stream is proportional to
\[
 Q\left(-\frac1{27}R-\frac4{135}R_{\mu\nu}u^\mu u^\nu\right).
\]
The spatially traceless-Ricci response and the Weyl response vanish for this
isotropic source.
\end{theorem}

\begin{proof}
In normal coordinates,
$R_{\mu\nu}x^\mu x^\nu=R_{uu}\eta^2+R_{ij}z^iz^j$ after isotropic frame
averaging.  Spatial isotropy reduces the contraction with
$K^{(4)}$ to the two covariance ratios displayed above.  Since
$R_{ii}=R+R_{uu}$ for signature $(-,+,+,+)$, the covariant coefficients are
$\alpha_R=\alpha_\perp$ and
$\alpha_{R_{uu}}=\alpha_u+\alpha_\perp$.
Substituting \Cref{lem:canonical-diamond-moments} gives
$\alpha_u=1/135$ and $\alpha_\perp=-1/27$, hence the stated values.
An isotropic spatial fourth-order tensor has no component in the spatially
traceless symmetric or Weyl representations.
\end{proof}

\begin{remark}[Reinterpretation of the earlier ball coefficient]
\label{rem:diamond-reinterpretation}
The value $-4/135$ found from the uniform spatial-ball normalization is recovered
by the genuine diamond calculation, but it is the coefficient of
$R_{\mu\nu}u^\mu u^\nu$, not the coefficient of $R$.  The genuine scalar-curvature
coefficient is $-1/27$.  The numerical values depend on the chosen profile;
the robust statement is the two-scalar-channel selection rule and the absence of
spatially traceless-Ricci and Weyl responses for isotropic Gaussian renewal
statistics.
\end{remark}

\begin{proposition}[Conditional determinacy of the $\eps^2$ bulk coefficients]
\label{prop:conditional-wilson}
Under \Cref{hyp:four-derivative-sector}, the covariant curvature-squared
coefficients are uniquely determined once the normalized bulk cumulant
coordinates $(\kappa_R,\kappa_S,\kappa_C)$ are specified:
\[
 c_1=G\lambda^2(\kappa_R-\tfrac14\kappa_S),\qquad
 c_2=G\lambda^2\kappa_S,\qquad
 c_3=G\lambda^2\kappa_C .
\]
The map is invertible.  The diamond coefficients of
\Cref{thm:diamond-linear-response} are response coefficients in the
$(R,R_{uu})$ tomography basis; converting them into a covariant bulk Wilson
triple requires the complete response map (or an averaging prescription over
$u$) and is not identified by the moment computation alone.
\end{proposition}

\begin{proof}
Using
$S_{\mu\nu}S^{\mu\nu}=R_{\mu\nu}^2-\tfrac14R^2$ gives
\[
 \mathfrak Q_4(\Riem)=G\lambda^2
 [(\kappa_R-\tfrac14\kappa_S)R^2
 +\kappa_S R_{\mu\nu}^2+\kappa_C C^2],
\]
which yields the stated map and its inverse.  The final sentence records the
distinction between a direction-dependent local diamond response and a covariant
bulk action.
\end{proof}

\begin{remark}[What is now quantitative]
\label{rem:diamond-response-scope}
For the specified canonical causal-diamond observable, the renewal model predicts
the exact response pair
\[
 (\alpha_R,\alpha_{R_{uu}})=(-1/27,-4/135)
\]
multiplying $Q=rs^4$.  This replaces the earlier spatial-ball surrogate.
It does not by itself assert numerical, scheme-independent values for the bulk
Wilson coefficients $(c_1,c_2,c_3)$.
\end{remark}


\subsection{Causal-diamond reconstruction of Wilson data}
\label{subsec:covariant-wilson-reconstruction}

The preceding quartic-moment calculation is linear in curvature and therefore cannot by itself determine the curvature-squared Wilson triple.  We now separate these orders, give the finite timelike-axis tomography that reconstructs the quadratic invariants, reduce the isotropic worked channel to a one-parameter Wilson ray, and evaluate its volume-element contribution explicitly.

\subsubsection{Linear-curvature audit}
\label{subsec:curvature-order-audit}
% =====================================================================

We work in $n=d+1=4$ spacetime dimensions, signature $(-,+,+,+)$, with $\bar u$
the renewal rest frame (the timelike direction in which the reset statistics are
isotropic).  The quartic-response theorem computes the order-$\eps^2$
contribution of the fourth renewal cumulant to the log of the calibrated diamond
partition function, using the metric volume factor to \emph{first} order in
curvature,
\begin{equation}
\label{eq:vol-linear}
 \sqrt{-g}(x)=1-\tfrac16 R_{\mu\nu}x^\mu x^\nu+O(|x|^3),
\end{equation}
and obtains, for the isotropic channel,
\begin{equation}
\label{eq:diamond-linear}
 \Delta\mathfrak D_{\mathrm{diam}}(u)
 =\eps^2 Q\bigl(\alpha_R\,R+\alpha_{R_{uu}}\,R_{\mu\nu}u^\mu u^\nu\bigr)
 +O(\eps^2\,\|\Riem\|^2),
 \qquad
 (\alpha_R,\alpha_{R_{uu}})=\Bigl(-\tfrac1{27},-\tfrac4{135}\Bigr).
\end{equation}

\begin{proposition}[The diamond pair is a linear-curvature datum]
\label{prop:diamond-linear-datum}
The response \eqref{eq:diamond-linear} is linear in the Riemann tensor.  It is the
metric variation, in the diamond axis $u$, of the local functional
\[
 \eps^2 Q\!\int\!\sqrt{-g}\,
 \bigl(\alpha_R\,R+\alpha_{R_{uu}}\,R_{\mu\nu}u^\mu u^\nu\bigr),
\]
which contains \emph{no} curvature-squared operator.  Consequently
$(\alpha_R,\alpha_{R_{uu}})$ cannot be the Wilson triple $(c_1,c_2,c_3)$
multiplying $\{R^2,\Ric^2,\Weyl^2\}$: those live at quadratic curvature order,
first reached at $O(|x|^4)$ in \eqref{eq:vol-linear}, which the cited computation
truncates.
\end{proposition}

\begin{proof}
Each term of \eqref{eq:diamond-linear} is degree one in $\Riem$; the
curvature-squared sector enters only through the $O(|x|^4)$ coefficient of
\eqref{eq:vol-linear}, namely the Riemann-bilinear part of the Riemann
normal-coordinate volume expansion \cite{GrayVanhecke,Brewin,Vassilevich}, which
is absent at the truncation used.  A degree-one curvature scalar is, by
definition, a variation of an Einstein--Hilbert-type (linear) action, not of a
curvature-squared action.
\end{proof}

\begin{proposition}[Covariant meaning of the linear response]
\label{prop:diamond-covariant-linear}
Averaging \eqref{eq:diamond-linear} over the unit timelike directions with the
isotropic measure, $\langle u^\mu u^\nu\rangle=\tfrac1n g^{\mu\nu}$, gives the
Lorentz-scalar
\[
 \bigl\langle\Delta\mathfrak D_{\mathrm{diam}}\bigr\rangle_u
 =\eps^2 Q\Bigl(\alpha_R+\tfrac1n\alpha_{R_{uu}}\Bigr)R
 =\eps^2 Q\Bigl(-\tfrac1{27}-\tfrac1{135}\Bigr)R
 =-\tfrac{2}{45}\,\eps^2 Q\,R .
\]
Hence the isotropically averaged diamond response is a pure renormalization of the
Einstein--Hilbert coefficient, $\delta\!\bigl(\tfrac1{16\pi\Gren}\bigr)=
-\tfrac{2}{45}\eps^2 Q$.  The non-averaged, $u$-dependent part is the
renewal-frame coupling $\eps^2 Q\,\alpha_{R_{uu}}R_{\mu\nu}\bar u^\mu\bar u^\nu$,
an Einstein--aether/preferred-frame term along $\bar u$, not a higher-curvature
invariant.
\end{proposition}
\begin{proof}
The isotropic average uses $\langle u^\mu u^\nu\rangle=\tfrac1n g^{\mu\nu}$ with
$n=4$, giving $\langle R_{\mu\nu}u^\mu u^\nu\rangle=\tfrac14R$.  If no frame is
preferred, only this average is physical and the term renormalizes $\Gren$.  If
the renewal substrate selects $\bar u$ (the isotropy frame of the reset kernel),
the residual $R_{\mu\nu}\bar u^\mu\bar u^\nu$ is a genuine covariant coupling of
Einstein--aether type \cite{JacobsonMattingly,Jacobson2008aether}.
\end{proof}

\begin{remark}[Curvature-order separation]
\label{rem:gap-restated}
The linear-response calculation correctly identifies the diamond coefficients as
``response coefficients in the $(R,R_{uu})$ tomography basis'' and not as the Wilson
triple.  Proposition~\ref{prop:diamond-linear-datum} makes the obstruction precise: it is
a curvature-\emph{order} separation, not merely an averaging issue.  The parenthetical
``averaging over $u$'' recovers the covariant \emph{linear} response of
Proposition~\ref{prop:diamond-covariant-linear}, i.e.\ an $\eps^2$ shift of $\Gren$ plus a
preferred-frame term---useful in its own right---but it does not and cannot yield
$(c_1,c_2,c_3)$.  The Wilson triple requires the quadratic response of
\S\ref{subsec:u-tomography}.
\end{remark}

% =====================================================================
\subsubsection{Quadratic-curvature tomography}
\label{subsec:u-tomography}
% =====================================================================

We now define the object that does carry the Wilson triple and prove that it is
invertible by the same finite tomography used for the two-profile Ricci
separation, lifted from rank two to rank four.

\begin{definition}[Quadratic diamond response]
\label{def:quadratic-diamond-response}
Let $\Psi(u)$ denote the order-$\eps^2$, curvature-squared part of the calibrated
diamond log-partition function with timelike axis $u$, obtained by carrying the
Riemann normal-coordinate expansion of $\sqrt{-g}$ and the van~Vleck weight to
$O(|x|^4)$ and contracting the resulting Riemann-bilinear coefficient against the
fourth renewal cumulant.  $\Psi(u)$ is a Lorentz scalar built from two curvature
tensors and the unit axis $u$.
\end{definition}

\begin{lemma}[Isotropic direction moments]
\label{lem:diamond-iso-moments}
On the unit directions in $n$ dimensions, with the rotation-invariant measure,
\[
 \langle u^\mu u^\nu\rangle=\tfrac1n g^{\mu\nu},\qquad
 \langle u^\mu u^\nu u^\rho u^\sigma\rangle
 =\tfrac1{n(n+2)}\bigl(g^{\mu\nu}g^{\rho\sigma}
 +g^{\mu\rho}g^{\nu\sigma}+g^{\mu\sigma}g^{\nu\rho}\bigr).
\]
Consequently, for any algebraic curvature tensor,
\[
 \langle R_{\mu\nu}u^\mu u^\nu\rangle=\tfrac1n R,\qquad
 \langle (R_{\mu\nu}u^\mu u^\nu)^2\rangle
 =\tfrac1{n(n+2)}\bigl(R^2+2\Ric^2\bigr),
\]
and every $u$-average of a curvature-squared scalar lies in
$\spn\{R^2,\Ric^2,\Riem^2\}$.
\end{lemma}

\begin{proof}
The displayed averages are the standard isotropic spherical moments
\cite{Vassilevich}.  Contracting them with curvature bilinears and using the
first Bianchi identity expresses every $u$-average of a curvature-squared scalar
in the three-dimensional invariant space $\spn\{R^2,\Ric^2,\Riem^2\}$.
\end{proof}

\begin{theorem}[$u$-tomographic Wilson inversion]
\label{thm:u-tomography}
Define the three direction moments of the quadratic response,
\[
 \Theta_0=\langle\Psi\rangle_u,\qquad
 \Theta_2=\bigl\langle\Psi\,(u^\mu u^\nu-\tfrac1n g^{\mu\nu})\bigr\rangle_u
          \!\cdot\! R_{\mu\nu},\qquad
 \Theta_4=\bigl\langle\Psi\,(u^{\otimes4}-\mathrm{traces})\bigr\rangle_u
          \!\cdot\! \mathcal P_{\Riem^2},
\]
where $\mathcal P_{\Riem^2}$ is the projection onto the Riemann-bilinear channel.
Then there is an explicit matrix $\mathsf G\in\R^{3\times3}$, depending only on
$n$ and the diamond profile, with
\[
 \begin{pmatrix}\Theta_0\\\Theta_2\\\Theta_4\end{pmatrix}
 =\mathsf G
 \begin{pmatrix} c_{R^2}\\ c_{\Ric^2}\\ c_{\Riem^2}\end{pmatrix},
\]
where $(c_{R^2},c_{\Ric^2},c_{\Riem^2})$ are the coefficients of $\Psi$ in the
basis $\{R^2,\Ric^2,\Riem^2\}$.  If $\det\mathsf G\neq0$ the three moments
reconstruct the triple, hence---after the Gauss--Bonnet rotation to
$\{R^2,\Ric^2,\Weyl^2\}$ and the cumulant identification of the previous analysis---the
Wilson coefficients $(c_1,c_2,c_3)$ and the cumulant coordinates
$(\kR,\kS,\kC)$.  The reconstruction is stable with conditioning
$\|\mathsf G^{-1}\|_2$, the rank-four analogue of $C_{\mathrm{rec}}$ in the
two-profile separation.
\end{theorem}

\begin{proof}
By Lemma~\ref{lem:diamond-iso-moments} each direction moment $\Theta_{2k}$ is a fixed
linear functional of the three invariants, with coefficients given by the
isotropic averages; assembling them produces $\mathsf G$.  The three invariants
have distinct multipole content under the little group of $u$: $R^2$ is a pure
$\ell=0$ direction-scalar, whereas $\Ric^2$ and $\Riem^2$ acquire $\ell=2,4$
components through the contractions in Lemma~\ref{lem:diamond-iso-moments}.  Hence the
rows of $\mathsf G$ are linearly independent and $\det\mathsf G\neq0$ for any
admissible profile, exactly as the two distinct second-moment profiles made the
two-profile moment matrix invertible.  Inverting $\mathsf G$ gives the triple;
the Gauss--Bonnet rotation $\Riem^2=\Weyl^2+2\Ric^2-\tfrac23R^2$ and
Proposition~``conditional-wilson'' of the previous analysis map it to
$(\kR,\kS,\kC)$ and $(c_1,c_2,c_3)$.
\end{proof}

\begin{remark}[What the machine needs as input]
\label{rem:machine-input}
Theorem~\ref{thm:u-tomography} is the rigorous inversion the previous analysis
required: it converts a direction-resolved \emph{quadratic} diamond response into
the covariant Wilson triple.  Its only input is $\Psi(u)$ of
Definition~\ref{def:quadratic-diamond-response}---the quadratic-curvature diamond moment,
one order beyond \eqref{eq:vol-linear}.  Supplying $\Psi$ for a given channel is a
finite moment computation in the diamond profile, identical in kind to the
canonical-moment evaluation of the previous analysis, now carried to $O(|x|^4)$.  The
inversion itself, including invertibility and conditioning, is unconditional.
\end{remark}

% =====================================================================
\subsubsection{Isotropic Wilson reduction}
\label{sec:isotropic}
% =====================================================================

\begin{theorem}[One-parameter Wilson ray for the worked channel]
\label{thm:one-param}
Let the reset stream be the compound-Poisson velocity-jump channel, with second and fourth cumulants $D=rs^2$ and $Q=rs^4$ and no anisotropic
or Clifford-odd cumulant.  Then the order-$\eps^2$ Wilson triple takes the form
\[
 (c_1,c_2,c_3)=\Gren\lambda^2\,\frac{Q}{D^2}\,(\rho_1,\rho_2,\rho_3),
\]
where $(\rho_1,\rho_2,\rho_3)$ are pure numbers independent of $(r,s)$ and of the
channel within the isotropic class.  Equivalently the standardized cumulant
coordinates $(\kR,\kS,\kC)$ are fixed multiples of the single dimensionless ratio
$Q/D^2$.
\end{theorem}

\begin{proof}
The isotropic fourth cumulant is the unique fully symmetric $O(d)$-invariant
four-tensor, $K^{(4)}_{ijkl}=Q(\delta_{ij}\delta_{kl}+\delta_{ik}\delta_{jl}
+\delta_{il}\delta_{jk})$, carrying a single scale $Q$
\cite{Isserlis}.  The quadratic-curvature response $\Psi$ of
Definition~\ref{def:quadratic-diamond-response} is linear in $K^{(4)}$ (it is a single
fourth-cumulant insertion), so $\Psi=Q\,\Psi_1$ with $\Psi_1$ a fixed geometric
functional independent of $r,s$.  Standardization divides by $D^2$
(the covariance normalization $\mathcal C_2^{-1/2\otimes4}$ of the previous analysis),
giving $(\kR,\kS,\kC)\propto Q/D^2$ with universal ratios; Proposition
``conditional-wilson'' then gives the displayed $(c_i)$.  The ratios are the
fixed numbers $(\rho_i)$, computed in \S\ref{sec:induced}.
\end{proof}

\begin{remark}[Not pure $R^2$]
\label{rem:not-pure-r2}
It is tempting to expect that an isotropic source produces only $R^2$.  This is
\emph{false} at quadratic order: the isotropic four-tensor contracts nontrivially
with a traceless symmetric pair, $K^{(4)}_{ijkl}A_{ij}B_{kl}
=Q(\tr A\,\tr B+2A_{ij}B_{ij})$, so it has nonzero overlap with the
traceless-Ricci ($\St^2$) and electric-Weyl ($\Weyl^2$) bilinears.  The correct
statement is Theorem~\ref{thm:one-param}: a one-parameter ray with fixed ratios,
not a pure scalar.  Which of $(\rho_1,\rho_2,\rho_3)$ vanish is decided by the
universal contraction of \S\ref{sec:induced}, not by isotropy alone.  This
corrects the linear-order ``Weyl response vanishes'' intuition, which holds only
for the degree-one response of \S\ref{subsec:curvature-order-audit}.
\end{remark}

% =====================================================================
\subsubsection{Heat-kernel factorization}
\label{sec:induced}
% =====================================================================

The cleanest route to the numbers $(\rho_i)$ and to the matching amplitude
$\Aren$ is the covariant induced action, in the sense of the
Sakharov--Adler--Zee programme and its modern effective-field-theory formulation
\cite{Sakharov1968,Adler1982,Zee1981,VisserInducedGravity}.  In the present
renewal setting the heat-kernel factorization also turns the earlier matching
hypothesis into a theorem.

Recall the covariant Kramers--Moyal form of the reset generator,
\begin{equation}
\label{eq:KM}
 \mathcal R=\tfrac D2\Box+\tfrac Q8\Box^2+\cdots,
\end{equation}
the factor $\tfrac18$ being the exact isotropic fourth-moment contraction.  The
renewal induced gravitational action is the regularized
$W=\tfrac12\Tr\log(\partial_\tau+\mathcal R)$, whose local curvature-squared part
is the proper-time integral of the second Seeley--DeWitt coefficient of the
relevant spatial operator \cite{GilkeyHeat,AvramidiHeat,Vassilevich,
BarvinskyVilkovisky}.

\begin{theorem}[Matching factorization is a theorem, not a hypothesis]
\label{thm:induced-factorization}
Suppose the renewal generator splits as
$\partial_\tau+\mathcal R$ with $\mathcal R$ acting on the spatial factor and the
fluctuation vector $\psi$ supported on the temporal relaxation sector with
positive generator $\mathsf L$.  Then the curvature-squared part of $W$
factorizes as
\[
 c_i=\Aren\,\widehat{\mathcal C}^{(i)}_4\,\mathcal T_2,
 \qquad
 \mathcal T_2=\frac{2\langle\psi,\mathsf L^{-3}\psi\rangle}
 {\langle\psi,\mathsf L^{-1}\psi\rangle},
\]
where $\widehat{\mathcal C}^{(i)}_4$ is the normalized fourth-cumulant projection
onto the $i$-th curvature channel and $\Aren$ is the dimensionally normalized
Seeley amplitude of the spatial operator.  In particular the previous analysis's
Hypothesis ``renewal-induced-matching'' holds identically in the heat-kernel
scheme.
\end{theorem}

\begin{proof}
On the split operator the heat trace factorizes,
$\Tr e^{-\sigma(\partial_\tau+\mathcal R)}
=\Tr_\tau e^{-\sigma\partial_\tau}\cdot\Tr_{\mathrm{sp}}e^{-\sigma\mathcal R}$, so
the proper-time integral separates the temporal moment from the spatial
Seeley expansion.  The temporal factor, expanded to the order sourcing the
quartic response, is the ratio $\mathcal T_2$ of inverse-power moments of
$\mathsf L$ on $\psi$ \cite{BarvinskyVilkovisky}.  The spatial factor's
curvature-squared coefficient is the Seeley--DeWitt $a_2$ of $\mathcal R$, whose
fourth-cumulant ($Q$) part is linear in $\widehat{\mathcal C}_4$ and whose
overall normalization defines $\Aren$.  Multiplicativity of the heat trace under
the tensor split gives the displayed product.
\end{proof}

\begin{corollary}[Reduction of $\Aren$ and the ratios]
\label{cor:reduction}
The amplitude $\Aren$ and the isotropic ratios $(\rho_1,\rho_2,\rho_3)$ of
Theorem~\ref{thm:one-param} are fixed by the second Seeley--DeWitt coefficient of
the minimal fourth-order operator $\Box^2$ entering \eqref{eq:KM}:
\[
 (\rho_1,\rho_2,\rho_3)\propto
 \bigl(b_{R^2}^{(\Box^2)},\,b_{\Ric^2}^{(\Box^2)},\,b_{\Weyl^2}^{(\Box^2)}\bigr),
\]
the curvature-squared heat-kernel coefficients of $\Box^2$, with the overall
constant $\Aren$ set by the $\tfrac Q8$ normalization of \eqref{eq:KM} and the
proper-time measure.  These coefficients are the known minimal higher-derivative
Seeley data \cite{Gusynin,BarvinskyVilkovisky,Vassilevich}; they are universal
(channel-independent within the isotropic class), confirming
Theorem~\ref{thm:one-param}.
\end{corollary}

\begin{proof}
By Theorem~\ref{thm:induced-factorization} the channel structure of $c_i$ is carried by
$\widehat{\mathcal C}^{(i)}_4$, which for the isotropic channel is a single scale,
and by the Seeley $a_2$ of $\mathcal R$.  The $D$-part of $\mathcal R$ is
Laplace-type and contributes only the universal Gaussian baseline; the
fourth-cumulant excess is the $\tfrac Q8\Box^2$ insertion, whose curvature-squared
response is governed by the $\Box^2$ Seeley coefficients.  The ratios
$(\rho_i)$ are therefore those coefficients, and the normalization is $\Aren$.
\end{proof}

\begin{remark}[Scope of the rational coefficients]
\label{rem:rationals}
Corollary~\ref{cor:reduction} reduces the entire $\eps^2$ Wilson triple of the
isotropic channel to the single scale $Q/D^2$ times the minimal fourth-order
Seeley coefficients $b^{(\Box^2)}$.  Their explicit rational values are the
published minimal higher-derivative heat-kernel data
\cite{Gusynin,BarvinskyVilkovisky}; we deliberately reduce to them rather than
reproduce them from memory, to avoid a transcription error in a sign- and
combinatorics-heavy computation.  Inserting the cited $b^{(\Box^2)}$ closes
$(\rho_1,\rho_2,\rho_3)$ and $\Aren$ numerically.  This is the single remaining
quadrature, now isolated as a one-line citation rather than an open hypothesis.
\end{remark}

% =====================================================================

\subsection{Explicit isotropic volume-sector Wilson ray}
\label{sec:workedwilson-setup}
% =====================================================================

The parallel calculation establishes that the diamond response carrying the covariant
Wilson triple is the \emph{quadratic}-curvature part of the calibrated diamond
log-partition function, $\Psi(u)$, and that the pair
$(\alpha_R,\alpha_{R_{uu}})=(-\tfrac1{27},-\tfrac4{135})$ is the linear-curvature
datum one order below.  Here we evaluate $\Psi(u)$.

We isolate the same channel the linear response calculation used: the
metric \emph{volume element}.  The linear diamond response was generated by
$\sqrt{-g}=1-\tfrac16R_{\mu\nu}x^\mu x^\nu+O(|x|^3)$, i.e.\ the linear-curvature
term of the Riemann normal-coordinate (RNC) volume expansion, weighted by the
displacement moments.  The quartic renewal cumulant couples to the
\emph{fourth} displacement moment; the curvature-squared volume response is
therefore the $O(|x|^4)$ term of the RNC volume expansion, weighted by that fourth
moment.  This is the exact continuation of the linear computation, not a new
prescription.

Two structural facts make the evaluation closed-form.

\begin{enumerate}[leftmargin=1.7em]
\item \emph{Spatial isotropy.}  The compound-Poisson worked channel has fourth
cumulant $K^{(4)}_{ijkl}=Q(\delta_{ij}\delta_{kl}+\delta_{ik}\delta_{jl}
+\delta_{il}\delta_{jk})$, the unique fully symmetric $O(3)$-invariant tensor on the
reset slice.  Weighting the $O(|x|^4)$ volume coefficient by an isotropic fourth
moment is exactly the angular average that defines the \emph{volume of a small
geodesic ball}.  Hence the worked-channel volume response on the slice is the
Gray--Vanhecke ball-volume combination \cite{GrayVanhecke,Gray}.
\item \emph{Gauss projection.}  The diamond axis $u$ selects a reset slice
$\Sigma_u$.  At the diamond centre the slice is geodesic ($K_{ij}=0$), so its
intrinsic curvature equals the ambient Riemann tensor projected onto $\Sigma_u$
\cite{Besse}.  The $u$-dependence of $\Psi$ is entirely this projection.
\end{enumerate}

% =====================================================================
\subsubsection{Slice volume kernel}
\label{subsec:worked-kernel}
% =====================================================================

\begin{lemma}[Three-dimensional Gray--Vanhecke combination]
\label{lem:gv-three}
For a geodesic ball of radius $r$ in a three-dimensional Riemannian slice,
\[
 \frac{\mathrm{Vol}(B_r)}{\omega_3 r^3}
 =1-\frac{\Rth}{30}\,r^2
 +\frac{r^4}{12600}\,\bigl(8\,\Rth{}^2-4\,|\Ric^{(3)}|^2\bigr)
 +(\text{total derivative})+O(r^6),
\]
where $\Rth,\Ric^{(3)}$ are the intrinsic scalar and Ricci curvatures of the slice.
\end{lemma}

\begin{proof}
The Gray--Vanhecke expansion in dimension $n$ has $r^4$ numerator
$-3|\Riem|^2+8|\Ric|^2+5R^2-18\,\Delta R$ with prefactor
$[360(n+2)(n+4)]^{-1}$ \cite{GrayVanhecke,Gray}.  In $n=3$ the Weyl tensor
vanishes and $|\Riem^{(3)}|^2=4|\Ric^{(3)}|^2-\Rth{}^2$, so the numerator collapses
to $8\Rth{}^2-4|\Ric^{(3)}|^2$ modulo the total derivative $\Delta\Rth$, and
$360\cdot5\cdot7=12600$.
\end{proof}

The local curvature-squared Wilson data are insensitive to the $\Delta\Rth$ total
derivative, so the worked-channel volume response per unit quartic weight is the
scalar
\begin{equation}
\label{eq:psi-def}
 \Psi(u)\;\propto\;8\,\Rth(u)^2-4\,|\Ric^{(3)}(u)|^2,
\end{equation}
with $\Rth(u),\Ric^{(3)}(u)$ the curvatures of the geodesic slice orthogonal to
$u$, evaluated through the Gauss projection.

% =====================================================================
\subsubsection{Frame-resolved response}
\label{subsec:worked-frame}
% =====================================================================

Work in a Euclidean orthonormal frame at the centre (the Wick-rotated diamond, in
which the slice is a genuine geodesic ball and the axis $u$ a unit vector; the
curvature-squared ratios continue back unchanged).  Let
$h^{\mu\nu}=\delta^{\mu\nu}-u^\mu u^\nu$ be the slice projector.  At $K=0$ the Gauss
equation gives $\Rth=h^{\mu\rho}h^{\nu\sigma}R_{\mu\nu\rho\sigma}=R-2\Ruu$ and
$\Ric^{(3)}_{\mu\nu}=h_\mu^\alpha h_\nu^\beta\,(R_{\alpha\beta}-R_{\alpha u\beta u})$,
where $\Ruu=R_{\mu\nu}u^\mu u^\nu$ and $R_{\alpha u\beta u}=u^\gamma u^\delta
R_{\alpha\gamma\beta\delta}$.

\begin{lemma}[Closed form for the slice invariants]
\label{lem:slice-closed}
With $T_{\alpha\beta}:=R_{\alpha\beta}-R_{\alpha u\beta u}$,
\[
 |\Ric^{(3)}|^2=|T|^2-2|Tu|^2+(uTu)^2,
 \qquad uTu=\Ruu,\quad (Tu)_\alpha=R_{\alpha\mu}u^\mu,
\]
and hence, writing $W_1:=R^{\mu\nu}R_{\mu u\nu u}$,
$W_2:=R_{\mu u\nu u}R^{\mu u\nu u}$, $(R^2)_{uu}:=R_{\mu\alpha}R^{\alpha}{}_{\nu}
u^\mu u^\nu$,
\[
 \boxed{\;
 \Psi(u)=8R^2-32\,R\,\Ruu+28\,\Ruu^2-4\,\Ric^2
 +8\,W_1-4\,W_2+8\,(R^2)_{uu}\; }.
\]
\end{lemma}

\begin{proof}
Expanding $h^{\mu\rho}h^{\nu\sigma}\Ric^{(3)}_{\mu\nu}\Ric^{(3)}_{\rho\sigma}$ with
$h=\delta-uu$ gives $|T|^2-2|Tu|^2+(uTu)^2$.  The antisymmetry
$R_{\alpha\gamma\beta\delta}=-R_{\alpha\gamma\delta\beta}$ kills the totally
$u$-contracted terms, giving $uTu=\Ruu$ and $(Tu)_\alpha=R_{\alpha\mu}u^\mu$;
$|T|^2=\Ric^2-2W_1+W_2$.  Substituting into
$\Psi=8(R-2\Ruu)^2-4|\Ric^{(3)}|^2$ and collecting terms gives the boxed
expression.  Each contraction was checked against direct numerical projection of
random algebraic curvature tensors (Remark~\ref{rem:numerics}).
\end{proof}

Equation~\eqref{lem:slice-closed} is the explicit frame-resolved volume response.  It
exhibits the physical content directly: besides the Lorentz scalars $R^2,\Ric^2$ it
carries the preferred-frame couplings $R\,\Ruu$, $\Ruu^2$, $W_1$, $W_2$, $(R^2)_{uu}$
along the renewal axis $u$.  When the renewal substrate fixes a rest frame $\bar u$,
these are genuine Einstein--aether / Ho\v{r}ava-type higher-curvature terms
\cite{JacobsonMattingly,Horava}; $W_2$ contains the electric-Weyl square
$E_{ij}E^{ij}$ with $E_{ij}=C_{i\mu j\nu}u^\mu u^\nu$.

% =====================================================================
\subsubsection{Covariant Wilson triple}
\label{sec:wilson}
% =====================================================================

Averaging \eqref{lem:slice-closed} over the diamond axis restores Lorentz invariance.

\begin{theorem}[Worked-channel volume Wilson triple]
\label{thm:triple}
With the isotropic direction averages
$\langle u^\mu u^\nu\rangle=\tfrac14\delta^{\mu\nu}$ and
$\langle u^\mu u^\nu u^\rho u^\sigma\rangle
=\tfrac1{24}(\delta^{\mu\nu}\delta^{\rho\sigma}+\delta^{\mu\rho}\delta^{\nu\sigma}
+\delta^{\mu\sigma}\delta^{\nu\rho})$, the axis-averaged volume response is
\[
 \langle\Psi\rangle_u
 =\tfrac76\,R^2+\tfrac{13}{6}\,\Ric^2-\tfrac14\,\Riem^2
 =\tfrac53\,R^2+\tfrac53\,S^2-\tfrac14\,\Weyl^2,
\]
$S$ the traceless Ricci tensor.  Hence the standardized cumulant coordinates and the
Stelle-basis Wilson triple of the worked channel are, on the one-parameter ray
$\Gren\lambda^2\,Q/D^2$,
\[
 (\kR,\kS,\kC)=\tfrac{Q}{D^2}\,\Gren\lambda^2\,(\tfrac53,\tfrac53,-\tfrac14)
 \;\propto\;(20,20,-3),
 \qquad
 (c_1,c_2,c_3)\;\propto\;(15,20,-3),
\]
where $c_1=\kR-\tfrac14\kS$, $c_2=\kS$, $c_3=\kC$ in the basis
$\{R^2,\Ric^2,\Weyl^2\}$.
\end{theorem}

\begin{proof}
Average the boxed expression term by term.  Using
$\langle\Ruu\rangle=\tfrac14R$, $\langle\Ruu^2\rangle=\tfrac1{24}(R^2+2\Ric^2)$,
$\langle(R^2)_{uu}\rangle=\tfrac14\Ric^2$,
$\langle W_1\rangle=\tfrac14\Ric^2$,
$\langle W_2\rangle=\tfrac1{24}(\Ric^2+\tfrac32\Riem^2)$, one gets
$\langle 8(R-2\Ruu)^2\rangle=\tfrac43R^2+\tfrac83\Ric^2$ and
$\langle4|\Ric^{(3)}|^2\rangle=\tfrac16R^2+\tfrac12\Ric^2+\tfrac14\Riem^2$, whose
difference is $\tfrac76R^2+\tfrac{13}6\Ric^2-\tfrac14\Riem^2$.  The
identities $\Ric^2=S^2+\tfrac14R^2$ and $\Riem^2=\Weyl^2+2S^2+\tfrac16R^2$ in $n=4$
convert to the $(R^2,S^2,\Weyl^2)$ basis, giving $(\tfrac53,\tfrac53,-\tfrac14)$.
The cumulant standardization divides by $D^2$ and the bulk normalization supplies
$\Gren\lambda^2 Q$; Proposition ``conditional-wilson'' of the previous analysis gives the
Stelle triple.  All coefficients were independently reproduced by an exact
isotropic-moment fit over random curvature tensors with residual $<10^{-10}$
(Remark~\ref{rem:numerics}).
\end{proof}

\begin{corollary}[Structure of the worked-channel correction]
\label{cor:structure}
The worked channel produces, at order $\eps^2$, a higher-curvature action with
\begin{enumerate}[leftmargin=1.7em]
\item \emph{equal scalar and traceless-Ricci weight}, $\kR=\kS$;
\item a \emph{nonzero but subdominant Weyl term}, $\kC/\kR=-3/20$, so the response
is \emph{not} pure $R^2$ (confirming the isotropic reduction), yet is Ricci-sector dominated;
\item a \emph{positive scalar coefficient} $c_1=\tfrac54\,\Gren\lambda^2 Q/D^2>0$
(for $Q=rs^4>0$ and the positive Gray--Vanhecke measure excess), the
Starobinsky-stable sign, consistent with the $R^2$-protected
flatness branch.
\end{enumerate}
\end{corollary}

\begin{remark}[Numerical verification]
\label{rem:numerics}
Every contraction in Lemmas~\ref{lem:slice-closed} and Theorem~\ref{thm:triple} was checked
by generating generic algebraic curvature tensors as sums of Kulkarni--Nomizu
products of random symmetric matrices (which satisfy all Riemann symmetries and the
first Bianchi identity), projecting onto the geodesic slice directly, and comparing
$8\Rth{}^2-4|\Ric^{(3)}|^2$ with the boxed formula and with the averaged triple.
Frame-resolved agreement was to machine precision; the averaged coefficients
$(\tfrac76,\tfrac{13}6,-\tfrac14)$ were recovered by exact isotropic-moment
substitution with residual below $10^{-10}$ and by independent Monte-Carlo over the
axis.  No coefficient is quoted from memory.
\end{remark}

% =====================================================================

\subsection{Gaussian universality at Einstein order}

\begin{definition}[Gaussian renewal class]
\label{def:gaussian-renewal-class}
Two primitive renewal models belong to the same Gaussian renewal class if they
have the same local reset rate $\lambda$, Clifford velocity form $g$, invariant
predictive state, and projected reversible quadratic transport.  If screen
thermodynamics is compared as well, we additionally require the same stationary
screen entropy density $\eta_{\rm scr}$.
\end{definition}

\begin{theorem}[Renewal universality of Einstein order]
\label{thm:renewal-universality}
Models in the same Gaussian renewal class have the same principal band
Hamiltonian, Hamilton--Jacobi action, invariant van Vleck response, and local
Ricci reconstruction.  When they also share $\eta_{\rm scr}$ and the same modular
normalization, they produce the same local Einstein equation and Newton coupling.
Their first model-dependent difference in the centered parity-even expansion can
occur at order $\eps^2$ through the fourth renewal cumulant tensor.
\end{theorem}

\begin{proof}
The principal tilted eigenvalue and its Hamilton flow are fixed by
$(\lambda,g)$, hence so are the action and mixed-endpoint Jacobian.  Reversibility
removes the projected first-order drift up to the already recorded deficiency
potential, so the flat-normalized Gaussian response is likewise fixed.  The
Clausius assembly additionally uses the screen entropy density and modular
normalization; equality of those data gives equality of $G_{\rm ren}$ and the
Einstein-order equation.  After centering and parity, the first remaining local
statistical invariant beyond the covariance is the fourth cumulant, which enters
the finite-window quartic response.
\end{proof}

\begin{corollary}[Renewal interpretation]
\label{cor:renewal-universality-slogan}
Einstein propagation and curvature reconstruction form the Gaussian universality
class of renewal memory, while higher-curvature corrections retain information
about non-Gaussian renewal statistics.
\end{corollary}

\begin{corollary}[Corrected field equation]
\label{cor:corrected-field-equation}
The renewal field equation through $\eps^2$ is
$G_{\mu\nu}+\Lambda g_{\mu\nu}+\eps^2H^{(2)}_{\mu\nu}=8\pi G\,\Tren_{\mu\nu}+O(\eps^3)$,
$H^{(2)}_{\mu\nu}=16\pi G\sum_{i=1}^5 c_i\,(\sqrt{-g})^{-1}\delta_{g^{\mu\nu}}\!\int I_i\sqrt{-g}$,
$I_i\in\{R^2,R_{\alpha\beta}^2,C^2,(\nabla\lambda)^2,\lambda^2R\}$.
\end{corollary}

\begin{proof}
Vary \Cref{def:effective-action}; the Einstein--Hilbert part gives
$G_{\mu\nu}+\Lambda g_{\mu\nu}$ and the $\eps^2$ functional gives $H^{(2)}_{\mu\nu}$;
remaining terms are higher order, divergences, or renormalizations of $G,\Lambda$.
\end{proof}

\subsection{Further renewal observables and coarse-graining}
\label{subsec:renewal-native-observables}

The screen coupling can be computed beyond the independent-output limit.  Let
\(P\) be a primitive finite-state output kernel with stationary distribution
\(\pi P=\pi\).  Its entropy per resolved screen cell is
\[
 h(P)=-\sum_i\pi_i\sum_jP_{ij}\log P_{ij}.
\]
If \(a_\eps=\Npartial\eps^{d-1}\) is the metric area represented by one cell,
then
\[
 \eta_{\rm scr}=\frac{h(P)}{a_\eps},
 \qquad
 G_{\rm ren}=\frac{a_\eps}{4h(P)}.
\]
For
\(P=\left(\begin{smallmatrix}1-p&p\\q&1-q\end{smallmatrix}\right)\),
\(\pi=(q,p)/(p+q)\); at \(p=1/3\), \(q=1/2\) one finds
\(h(P)\simeq0.659\) nats per cell and
\(G_{\rm ren}/a_\eps\simeq0.379\).  The one-site binary entropy is recovered
only when the two rows coincide, namely \(p+q=1\).  Thus temporal renewal
correlations feed directly into the entropy-per-cell input, while the conversion
to a continuum Newton coupling also retains the geometric cell-area calibration.

The metric and homogeneous deficiency may also be encoded in one joint tilted
spectral object.  Augment the velocity-jump transfer operator by a signed count
\(\chi\) for deficiency minus completion events and denote the joint operator by
\(\mathcal T_{\theta,\chi}\).  In the embedded-step parametrization,
\[
 h(\theta)=-\log r(\mathcal T_{\theta,0}),
 \qquad
 \Bstep=\left.\partial_\chi\log r(\mathcal T_{0,\chi})\right|_{\chi=0},
 \qquad
 B=\Bphys=\frac{\Bstep}{\bartau}.
\]
For forward and backward signed Poisson processes with physical rates \(a,b\),
the mean lapse has already been absorbed into the continuous-time generator, so
\[
 B=a-b,\qquad H=\frac{a-b}{d},\qquad
 \Lambda=\frac{d-1}{2d}(a-b)^2,
 \qquad \operatorname{Var}(H_T)=\frac{a+b}{d^2T}.
\]
Geometry and cosmological deficiency are therefore different derivatives of one
renewal pressure, conjugate respectively to displacement and signed reset
imbalance; the conversion from an embedded increment to a physical rate is
always the clock division by \(\bartau\).

For irreversible renewal memory, consider the three-state ring with clockwise
rate $k$ and counter-clockwise rate $k'$.  Its stationary current and entropy
production are
\[
 j_{\rm st}=\frac{k-k'}{3},\qquad
 \sigma_{\rm ep}=(k-k')\log\frac{k}{k'}\ge0.
\]
Whenever the projected scalar WKB drift of this concrete embedding is linear in
the cycle current, $\kappa_{\rm irr}=C_0j_{\rm st}$, the elementary inequality
$\log(k/k')\ge(k-k')/\max(k,k')$ yields the controlled estimate
\[
 |\kappa_{\rm irr}|^2\le
 \frac{C_0^2}{9}\max(k,k')\,\sigma_{\rm ep}.
\]
The inequality is therefore conditional only on the model-specific linear
response coefficient $C_0$; it does not assert a universal pointwise equality
between amplitude drift and entropy production.

Finally, let $\Xi_n=\sum_{m=1}^n\xi_m$ be an additive displacement functional of
a primitive renewal chain with a spectral gap and finite cumulants.  Exponential
mixing and the Markov additive-functional central limit theorem imply
$K_n^{(k)}=O(n)$.  After covariance normalization,
\[
 \widehat K_n^{(k)}=O\!\left(n^{1-k/2}\right),\qquad
 \widehat K_n^{(3)}=O(n^{-1/2}),\qquad
 \widehat K_n^{(4)}=O(n^{-1}).
\]
Thus the Gaussian renewal class is an attractive fixed manifold under event
blocking.  This statement is invariant under whether one parametrizes the
blocked chain by event number or preserves physical time; only the bookkeeping
of the effective event rate changes.  The second-cumulant Einstein layer is
stable, while non-Gaussian higher-curvature directions become irrelevant under
coarse-graining.


\subsection{Wald entropy and derivative-expansion validity}

Once the effective action is known, the horizon entropy is no longer given only
by the leading area density.  The appropriate correction is the Wald functional
of the four-derivative action.  This provides a consistency check between the
finite-window bulk dynamics and the stationary screen thermodynamics.

\begin{theorem}[Renewal Wald entropy]
\label{thm:wald}
For a stationary horizon cross-section $\Sigma$,
$S_{\hor}(\Sigma)=\Area(\Sigma)/4G+\eps^2\Delta S_{\rm Wilson}(\Sigma)+O(\eps^3)$,
$\Delta S_{\rm Wilson}=-2\pi\int_\Sigma(\partial\Lcal_2/\partial R_{\mu\nu\rho\sigma})
\epsilon_{\mu\nu}\epsilon_{\rho\sigma}\,dA$.  If a microscopic channel yields $c_3>0$, the Weyl term gives a sign-definite
two-horizon split absent in pure Einstein gravity.
\end{theorem}

\begin{proof}
The Wald formula \cite{WaldEntropy,IyerWald} applied to
$\Lcal=\tfrac1{16\pi G}(R-2\Lambda)+\eps^2\Lcal_2$ gives $A/4G$ plus the displayed
functional derivative; whenever $c_3>0$, the Weyl term contributes with fixed sign to the
difference between cuts with different Weyl tidal data, a split absent when entropy
depends only on area.
\end{proof}

\begin{remark}[Effective, not fundamental]
The $\eps^2$ layer is a controlled effective correction, not a complete ultraviolet
quantum theory; phenomenological consequences of resumming such corrections belong to
a later observable-consequences paper \cite{PelissierPhenomenology}.
\end{remark}

\subsubsection{The trans-renewal curvature boundary}
\label{subsec:trans-renewal-boundary}

The derivative expansion is informative only while successive orders remain
parametrically separated.  We therefore conclude the section by combining the
curvature, rate-gradient, and finite-window scales into a single invariant
control parameter.  Its breakdown marks the boundary of the finite-order
continuum description, not a proposed microscopic phase beyond it.

Define
\[
 \chi_{\rm ren}(x):=\max\Big\{\eps^2|R|,\,\eps^2\sqrt{|R_{\mu\nu}R^{\mu\nu}|},\,
 \eps^2\sqrt{|R_{\mu\nu\rho\sigma}R^{\mu\nu\rho\sigma}|},\,
 \eps^2\tfrac{|\nabla\lambda|^2}{\lambda^2}\Big\}(x),
\]
the absolute values making $\chi_{\rm ren}$ a nonnegative control parameter (any
equivalent norm gives the same criterion).

\begin{proposition}[Renewal curvature boundary]
\label{prop:renewal-curvature-boundary}
If the normalized renewal cumulants remain uniformly bounded, then on every region
with $\chi_{\rm ren}\ll1$ the renewal effective dynamics is hierarchically ordered:
the Einstein term dominates, the four-derivative action gives the first relative
correction $O(\chi_{\rm ren})$, and the omitted higher-derivative terms are
parametrically smaller.  When $\chi_{\rm ren}=O(1)$, successive higher-cumulant
contributions need no longer be suppressed and no finite truncation determines the
local dynamics.
\end{proposition}

\begin{proof}
A local term with $2n$ derivatives beyond Einstein carries $\eps^{2n}$, relative size
$(\eps^2/L_{\rm curv}^2)^n$ on a background of curvature radius $L_{\rm curv}$; bounded
normalized cumulants order the expansion when $\chi_{\rm ren}\ll1$, and at
$\chi_{\rm ren}=O(1)$ the suppression is lost.
\end{proof}

\begin{corollary}[Controlled de Sitter scale separation]
\label{cor:controlled-ds-separation}
On the pure-deficiency branch $H=B/d$, $L_{\dS}=d/B$, and a parametrically controlled
de Sitter continuum requires $H\eps=B\eps/d\ll1$, i.e.\ $L_{\dS}\gg\eps$: $B$ fixes the
infrared curvature scale, $\eps$ the ultraviolet boundary.
\end{corollary}

\begin{proof}
For de Sitter curvature all quadratic norms are $O(H^2)$, so $\chi_{\rm ren}\ll1$
reduces to $H\eps\ll1$; substituting $H=B/d$ gives the equivalence.
\end{proof}

\begin{remark}[What this boundary does not imply]
\label{rem:no-bounce-from-boundary}
$\chi_{\rm ren}=O(1)$ is a breakdown criterion, not a bounce equation: it neither
produces a finite critical density nor forces $H$ to pass from contraction to
expansion.  A relation $H^2\propto\rho(1-\rho/\rho_c)$ would require separate
nonperturbative renewal dynamics, not assumed here.
\end{remark}


\begin{corollary}[Two faces of the trans-renewal boundary]
\label{cor:two-sided-trans-renewal}
The condition $\chi_{\rm ren}=O(1)$ has the same logical meaning on the outward
cosmological face and on the inward high-curvature face of a black-hole interior:
in both cases the finite-order continuum expansion loses predictive closure when
the renewal correlation length becomes comparable to the local curvature scale.
Consequently the early-universe emergence problem and the black-hole interior
problem share a boundary criterion, while the continuation across that boundary
requires nonperturbative renewal data not contained in the local Einstein action
\cite{PelissierBlackHoles}.
\end{corollary}

\begin{proof}
The proof is the scale estimate of \Cref{prop:renewal-curvature-boundary}.  It
depends only on the comparison between the renewal scale and the local curvature
radius, not on whether the high-curvature region is approached by evolving
outward from an initial boundary or inward along a collapsing/horizon interior
geometry.
\end{proof}

% =====================================================================

% =====================================================================
\section{Dynamical dark energy from renewal currents}
\label{sec:dark-energy}
% =====================================================================

The homogeneous renewal current has two sharply different macroscopic uses.
At equilibrium it fixes the vacuum term and hence the exact de~Sitter branch.
Away from equilibrium it may be promoted to a coarse balance variable, but then
one must specify its stress tensor and its exchange law.  This section keeps
three registers distinct:
\begin{enumerate}[label=\textup{(\roman*)},leftmargin=2.6em]
\item the \emph{vacuum register}, in which the deficiency enters only through
      \(\Lambda(B)g_{\mu\nu}\);
\item the \emph{effective-fluid register}, in which
      \(\rho_B=c_BB^2\) and the pressure is fixed by a continuity equation;
\item the \emph{open exchange register}, in which matter and the deficiency
      fluid exchange energy while the total stress tensor is conserved.
\end{enumerate}
The first register is rigidly constant by covariance.  The second permits
reversible fractional relaxation but, under the convex Perron potential, cannot
cross the phantom divide.  The third can cross only when an irreversible cycle
current overcomes the restoring drift.  No numerical crossing redshift is
claimed without calibrating that exchange completion.

\subsection{Perron current and rate function}
\label{subsec:dark-energy-rate-function}

The physical-depth pressure of \Cref{thm:matrix-renewal-dictionary} already
contains the homogeneous current coordinate.  We now record its convex dual,
which supplies the canonical thermodynamic potential for a coarse deficiency
fluid.

\begin{definition}[Deficiency pressure and rate function]
\label{def:deficiency-rate-function}
Set
\[
  \Lambda_B(\chi):=\Pphys(\chi,0,0),
  \qquad B_*:=\Lambda_B'(0)=\Bphys,
\]
and define the Legendre--Fenchel transform
\[
  I_B(b):=\sup_{\chi\in\R}\{\chi b-\Lambda_B(\chi)\}.
\]
The variable \(b\) is a physical-depth signed-current rate, with the same units
as the homogeneous deficiency \(B\).
\end{definition}

\begin{theorem}[Perron large-deviation potential]
\label{thm:deficiency-rate-function}
Assume the tilted matrix-renewal kernel of
\Cref{def:tilted-matrix-renewal} has exponential moments on an open interval
containing the origin and that \(\Lambda_B\) is essentially smooth and strictly
convex there.  Then the empirical physical-depth deficiency current satisfies a
large-deviation principle with good convex rate function \(I_B\).  Moreover
\[
  I_B(B_*)=I_B'(B_*)=0.
\]
If the asymptotic physical-depth variance
\(\sigma_B^2:=\Lambda_B''(0)\) is finite and strictly positive, then
\[
  I_B''(B_*)=\frac{1}{\sigma_B^2}.
\]
Thus the same tilted Perron family that fixes the de~Sitter current also fixes the
local restoring potential of the dynamical deficiency register.
\end{theorem}

\begin{proof}
Analytic Perron--Frobenius perturbation gives the differentiable scaled
cumulant-generating function \(\Lambda_B\), as in
\Cref{thm:matrix-renewal-dictionary}.  The G\"artner--Ellis theorem then gives
the large-deviation principle with Legendre--Fenchel rate function
\cite{DemboZeitouni,Touchette2009}.  Legendre duality places the unique minimum
at \(B_*=\Lambda_B'(0)\), with value zero.  Where the duality map is smooth,
\(I_B''(B_*)=[\Lambda_B''(0)]^{-1}\), proving the final identity.
\end{proof}

\begin{proposition}[Mixed current--screen and current--metric susceptibilities]
\label{prop:mixed-perron-susceptibilities}
Let \(\Pphys(\chi,k,\vartheta)\) be the physical-depth pressure of
\Cref{thm:matrix-renewal-dictionary}.  Define
\[
 \kappa_{B\eta}:=\partial_\chi\partial_\vartheta\Pphys(0,0,0),
 \qquad
 C_{ab}(\chi):=-\partial_{k_a}\partial_{k_b}\Pphys(\chi,0,0),
\]
and
\[
 \kappa_{B\eps}:=
 \left.\partial_\chi\left(\frac1d\tr C(\chi)\right)^{1/2}\right|_{\chi=0}
 =\frac{1}{2d\eps}\,\partial_\chi\tr C(0).
\]
Then \(\kappa_{B\eta}\) is the connected zero-frequency cross-cumulant per
physical depth between the deficiency increment and screen surprisal, whereas
\[
  \partial_\chi C_{ab}(0)
  =-\partial_\chi\partial_{k_a}\partial_{k_b}\Pphys(0,0,0)
\]
is the connected third cumulant with one deficiency and two displacement
insertions.  If the base kernel is reversible, the deficiency observable is odd
under time reversal, and the screen-surprisal and quadratic-displacement
observables are even, then
\[
  \kappa_{B\eta}=\kappa_{B\eps}=0.
\]
Away from that symmetry locus these derivatives quantify the response of screen
capacity and metric calibration to an open deficiency bias, but do not by
themselves supply its evolution law.
\end{proposition}

\begin{proof}
Repeated differentiation of the logarithmic Perron eigenvalue produces the
connected cumulants of the corresponding additive observables
\cite{Nagaev,KipnisVaradhan}.  The formula for \(\kappa_{B\eps}\) is the chain
rule applied to \(\eps(\chi)^2=d^{-1}\tr C(\chi)\).  Under the reversing
involution, a stationary correlation containing one odd insertion and otherwise
even insertions changes sign while its law remains invariant, hence vanishes.
Absolute convergence of the Green--Kubo series permits this argument term by
term.
\end{proof}

\subsection{Macroscopic stress-energy realizations}
\label{subsec:dark-energy-registers}

The equation-of-state language is meaningful only after the stress tensor has
been specified.  The next theorem is the covariance obstruction in its cleanest
form.

\begin{theorem}[Closed vacuum-deficiency no-go]
\label{thm:closed-source-free-deficiency-nogo}
Suppose the homogeneous deficiency contributes to the calibrated Einstein
equation only through
\[
  \Lambda(B)g_{\mu\nu},
  \qquad
  \Lambda(B)=\frac{d-1}{2d}B^2,
\]
and suppose the remaining matter stress tensor is separately conserved.  If
there is no exchange current between matter and the vacuum-deficiency sector,
then
\[
  \nabla_\nu\Lambda(B)=0.
\]
Consequently \(B\) is constant on every connected region on which \(B\ne0\).
The only nonzero closed homogeneous endpoint in this register is therefore the
cosmological-constant branch, with \(w=-1\).
\end{theorem}

\begin{proof}
The Einstein tensor is divergence-free.  Taking the divergence of
\[
  G_{\mu\nu}+\Lambda(B)g_{\mu\nu}
   =8\pi\Gren T^{\rm m}_{\mu\nu}
\]
and using \(\nabla^\mu T^{\rm m}_{\mu\nu}=0\) gives
\(\nabla_\nu\Lambda(B)=0\) \cite{WaldBook}.  Since
\(\nabla_\nu\Lambda=(d-1)B\nabla_\nu B/d\), one has
\(\nabla B=0\) wherever \(B\ne0\).
\end{proof}

\begin{proposition}[Equation-of-state identities]
\label{prop:deficiency-eos-identities}
Let a spatially flat FLRW background have Hubble rate \(H\), and let an effective
deficiency fluid have density
\[
  \rho_B=c_BB^2,
  \qquad c_B>0.
\]
\begin{enumerate}[label=\textup{(\roman*)},leftmargin=2.6em]
\item If the deficiency fluid is separately conserved, then
\[
  \boxed{
  w_B=-1-\frac{2\dot B}{dHB}.}
\]
On a deficiency-dominated calibrated branch, \(H=B/d\), this reduces to
\[
  \boxed{w_B=-1-2\frac{\dot B}{B^2}.}
\]
\item If matter transfers energy to the deficiency fluid at rate
\(\mathcal Q\), with sign convention
\[
 \dot\rho_{\mathrm{m}}+dH(\rho_{\mathrm{m}}+p_{\mathrm{m}})=-\mathcal Q,
 \qquad
 \dot\rho_B+dH(\rho_B+p_B)=\mathcal Q,
\]
then its physical pressure ratio is
\[
  \boxed{
  \frac{p_B}{\rho_B}
  =-1-\frac{2\dot B}{dHB}
    +\frac{\mathcal Q}{dHc_BB^2}.}
\]
An observer who reconstructs the same density while assuming no interaction
instead infers
\[
  \boxed{
  w_{\rm rec}:=-1-\frac1d\frac{d\log\rho_B}{d\log a}
  =-1-\frac{2\dot B}{dHB}.}
\]
Thus a reconstructed phantom-divide crossing is a turning point of \(B\), while
a crossing of the physical pressure ratio also depends on the exchange term.
\end{enumerate}
\end{proposition}

\begin{proof}
Substitute \(\dot\rho_B=2c_BB\dot B\) into the corresponding continuity
equation and divide by \(dH\rho_B\).  The reconstructed expression is the
standard logarithmic-density definition obtained when the interaction is
absorbed into an effective equation of state.
\end{proof}

\begin{theorem}[Open-budget adiabatic relation]
\label{thm:adiabatic-dark-energy}
Suppose \(B(t)>0\) is an open, slowly varying balance variable satisfying
\[
  \frac{|\dot B|}{B^2}\ll1,
  \qquad
  \frac{|\ddot B|}{B^3}\ll1,
\]
after the exchange or nonstationary renewal input responsible for its variation
has been specified.  Then local calibrated diamonds obey
\[
  H(t)=\frac{B(t)}{d}+O\!\left(\frac{\dot B}{B}\right),
\]
and in the deficiency-dominated effective-fluid register
\[
 \boxed{
 w_B(t)=-1-2\frac{\dot B}{B^2}
 +O\!\left(
 \frac{\dot B^2}{B^4},
 \frac{\ddot B}{B^3},
 \eps\right).}
\]
The theorem does not create a second closed vacuum branch: in the vacuum
register \Cref{thm:closed-source-free-deficiency-nogo} forces
\(\dot B=0\).
\end{theorem}

\begin{proof}
The homogeneous renewal calibration gives \(dH=B\) on each locally frozen
diamond.  Adiabatic variation produces the displayed derivative correction.
The deficiency-dominated equation-of-state identity is
\Cref{prop:deficiency-eos-identities}(i).  The final statement follows from the
Bianchi no-go theorem.
\end{proof}

\subsection{Fractional relaxation and reversible no-crossing}
\label{subsec:fractional-deficiency-relaxation}

The Perron potential fixes a canonical restoring force, but converting an
operational-depth gradient flow into physical cosmological time requires a clock
law.  At the critical finite-variance renewal endpoint the inverse-stable clock
has index \(\alpha_{\ren}=1/3\).

\begin{principle}[Critical fractional deficiency closure]
\label{prin:critical-deficiency-closure}
In the effective-fluid or open-exchange register, assume the coarse homogeneous
coordinate is projected onto the critical recurrent clock with no independent
heavy time scale.  Its physical-time closure is
\[
  C\,\prescript{\mathrm C}{}{D}_t^{\alpha}B(t)
  =-\Gamma I_B'(B(t))+S_{\rm exch}(t),
  \qquad C,\Gamma>0,
  \qquad \alpha=\frac13,
\]
where \(\prescript{\mathrm C}{}{D}_t^{\alpha}\) is the Caputo derivative.
This is a constitutive interface between the microscopic current large
deviations and cosmology; \(C\) and the physical-time normalization of
\(\Gamma\) remain matching data.
\end{principle}

\begin{proposition}[Subordination origin of the linearized closure]
\label{prop:deficiency-subordination}
Near the Perron minimum let
\[
 I_B(B)=\frac{m^2}{2}(B-B_*)^2+O((B-B_*)^3),
 \qquad m^2=I_B''(B_*)>0.
\]
If the operational-depth linear relaxation is
\[
  \frac{dy^{\rm op}}{d\tau}=-\lambda y^{\rm op},
  \qquad y^{\rm op}=B^{\rm op}-B_*,
  \qquad \lambda=\frac{\Gamma m^2}{C},
\]
and physical time is obtained by an inverse \(\alpha\)-stable subordinator, then
\[
  \prescript{\mathrm C}{}{D}_t^{\alpha}y(t)=-\lambda y(t),
  \qquad
  y(t)=y_0E_\alpha(-\lambda t^\alpha).
\]
At the critical endpoint \(\alpha=1/3\).
\end{proposition}

\begin{proof}
The operational semigroup is \(e^{-\lambda\tau}\).  Subordination by the density
of the inverse stable subordinator has Laplace transform
\[
  \widetilde y(s)=\frac{s^{\alpha-1}}{s^\alpha+\lambda}y_0,
\]
which is exactly the Laplace transform of the Caputo equation and of the
Mittag--Leffler solution
\cite{MetzlerKlafter,MeerschaertScheffler,Podlubny}.
\end{proof}

\begin{theorem}[Reversible fractional no-crossing]
\label{thm:source-free-no-crossing}
Assume \(S_{\rm exch}=0\), \(0<\alpha<1\), and that \(I_B\in C^1\) is convex
with a unique minimum \(B_*\), with \(I_B'\) locally Lipschitz.  In the scalar
fractional gradient flow of
\Cref{prin:critical-deficiency-closure}, a solution beginning above \(B_*\)
remains above it and relaxes monotonically toward it; the statement below the
minimum is symmetric.  Hence the reversible source-free effective-fluid branch
cannot cross the reconstructed value \(w=-1\).

For the quadratic potential the result is explicit:
\[
  B(t)-B_*=\Delta_0E_\alpha(-\lambda t^\alpha),
  \qquad \Delta_0=B(0)-B_*.
\]
For \(\Delta_0>0\),
\[
 B(t)-B_*\sim
 \frac{\Delta_0}{\lambda\Gamma(1-\alpha)}t^{-\alpha},
 \qquad
 w_B(t)+1\sim
 \frac{2\alpha\Delta_0}
 {B_*^2\lambda\Gamma(1-\alpha)}t^{-1-\alpha}.
\]
At \(\alpha=1/3\), the approach is \(B-B_*\sim t^{-1/3}\) and
\(w_B+1\sim t^{-4/3}\).
\end{theorem}

\begin{proof}
The Caputo equation is equivalent to the scalar Volterra equation
\[
 B(t)=B(0)-\frac{\Gamma}{C\Gamma(\alpha)}
 \int_0^t(t-s)^{\alpha-1}I_B'(B(s))\,ds.
\]
The fractional kernel is completely positive: its linear resolvents are the
nonnegative Mittag--Leffler functions
\(E_\alpha(-\mu t^\alpha)\) and
\(t^{\alpha-1}E_{\alpha,\alpha}(-\mu t^\alpha)\)
\cite{Pollard1948,Schneider1996}.  The scalar Volterra comparison theorem for a
monotone nonlinearity therefore preserves the order relative to the equilibrium
and gives monotone relaxation
\cite{GripenbergLondenStaffans}.  In the quadratic case the exact solution is
\Cref{prop:deficiency-subordination}; complete monotonicity gives its sign and
monotonicity.  The asymptotic expansion
\(E_\alpha(-x)\sim[x\Gamma(1-\alpha)]^{-1}\) and differentiation give the two
displayed tails.  Finally, \Cref{prop:deficiency-eos-identities} identifies the
sign of \(w+1\) with the sign of \(-\dot B\) in the reconstructed or separately
conserved register.
\end{proof}

\begin{remark}[No conflict with the Bianchi theorem]
The word ``source-free'' has different meanings in the two registers.  In
\Cref{thm:closed-source-free-deficiency-nogo} the entire contribution is a
vacuum tensor \(\Lambda(B)g_{\mu\nu}\), so covariance forces \(B\) to be
constant.  In \Cref{thm:source-free-no-crossing}, \(B\) is an effective fluid
with pressure determined by its own continuity equation; it is source-free only
in the sense of having no exchange with matter.  Its time dependence is carried
by \(p_B+\rho_B\ne0\), not by a varying vacuum term.
\end{remark}

\subsection{Irreversible exchange: threshold and entropy floor}
\label{subsec:irreversible-deficiency-exchange}

A crossing requires a source capable of reversing the sign of \(\dot B\).  The
minimal finite irreversible module is an oriented three-state ring coupled to
the scalar homogeneous budget.

\begin{construction}[Three-state irreversible exchange ring]
\label{constr:dark-energy-ring}
Let states \(0,1,2\) be arranged cyclically.  From each state the clockwise jump
rate is \(k>0\) and the counter-clockwise rate is \(k'>0\).  The stationary law
is uniform.  Define the per-edge stationary current, total oriented count rate,
and entropy-production rate by
\[
 j_{\rm st}:=\frac{k-k'}{3},
 \qquad
 \dot Q:=k-k'=3j_{\rm st},
 \qquad
 \sigma_{\rm ep}:=(k-k')\log\frac{k}{k'}\ge0.
\]
The tilted scaled cumulant-generating function for the total oriented count is
\[
  \lambda_Q(\chi)=ke^\chi+k'e^{-\chi}-(k+k'),
\]
so
\[
  \lambda_Q'(0)=\dot Q,
  \qquad
  \lambda_Q''(0)=k+k'.
\]
Reversibility \(k=k'\) is exactly the zero-current, zero-entropy-production
locus.
\end{construction}

\begin{proof}
The uniform measure solves the stationarity equations of the circulant
generator.  Tilting clockwise jumps by \(e^\chi\) and counter-clockwise jumps by
\(e^{-\chi}\) leaves a circulant matrix whose Perron eigenvalue is the displayed
\(\lambda_Q\).  Differentiation gives the current and variance.  The entropy
production is the Schnakenberg cycle affinity multiplied by the net current
\cite{Hill,Schnakenberg}.
\end{proof}

\begin{definition}[Scalar projection of the exchange current]
\label{def:exchange-projection}
Let \(\psi_B\) denote the homogeneous deficiency mode and
\(\mathcal Q_{\rm exch}\) the finite exchange observable.  Define
\[
 \kappa_{\rm ex}:=
 \frac{\langle\psi_B,\mathcal Q_{\rm exch}\rangle}
      {\langle\psi_B,\psi_B\rangle},
 \qquad 0\le\kappa_{\rm ex}\le1,
\]
after normalizing the exchange observable, and set
\[
  \sigma_0:=\kappa_{\rm ex}\dot Q.
\]
The projection form is fixed; its numerical value is a microscopic alignment
coefficient.
\end{definition}

\begin{theorem}[Linearized irreversible crossing threshold]
\label{thm:irreversible-crossing-threshold}
Let \(I_B(B)=\tfrac12m^2(B-B_*)^2\), set
\(y=B-B_*\), and use the decaying exchange profile
\(S_{\rm exch}(t)=\sigma_0e^{-\mu t}\), \(\mu>0\).  Then
\[
 \prescript{\mathrm C}{}{D}_t^\alpha y
 =-\lambda y+s_0e^{-\mu t},
 \qquad
 \lambda=\frac{\Gamma m^2}{C},
 \qquad
 s_0=\frac{\sigma_0}{C},
\]
has the exact solution
\[
 y(t)=\Delta_0E_\alpha(-\lambda t^\alpha)
 +s_0\int_0^t(t-s)^{\alpha-1}
 E_{\alpha,\alpha}[-\lambda(t-s)^\alpha]e^{-\mu s}\,ds.
\]
Its early-time expansion is
\[
  \dot y(t)=
  \frac{s_0-\lambda\Delta_0}{\Gamma(\alpha)}t^{\alpha-1}
  +o(t^{\alpha-1}).
\]
Consequently an initially nonphantom state with \(\Delta_0\ge0\) can enter a
reconstructed phantom phase only if
\[
 \boxed{
 \sigma_0>\Gamma m^2\Delta_0,}
\]
or equivalently
\[
 \boxed{
 j_{\rm st}>j_{\rm st}^{*}
 :=\frac{\Gamma m^2\Delta_0}{3\kappa_{\rm ex}}.}
\]
For a fixed source profile the exact critical amplitude is
\[
  \sigma_{0,{\rm crit}}
  =\inf_{t:\,\dot R(t)>0}
    \frac{-\dot y_{\rm hom}(t)}{\dot R(t)},
\]
where \(R\) is the unit-source response and
\(y_{\rm hom}=\Delta_0E_\alpha(-\lambda t^\alpha)\).
\end{theorem}

\begin{proof}
Laplace transformation gives
\[
 \widetilde y(s)
 =\frac{s^{\alpha-1}\Delta_0}{s^\alpha+\lambda}
 +\frac{s_0}{(s+\mu)(s^\alpha+\lambda)}.
\]
Inversion yields the displayed Mittag--Leffler convolution
\cite{Podlubny,MetzlerKlafter}.  Expanding
\(E_\alpha(-\lambda t^\alpha)
 =1-\lambda t^\alpha/\Gamma(1+\alpha)+o(t^\alpha)\)
and the source convolution
\(=s_0t^\alpha/\Gamma(1+\alpha)+o(t^\alpha)\) gives the derivative formula.
Since reconstructed \(w+1\) has the sign of \(-\dot B\), a phantom opening
requires \(s_0>\lambda\Delta_0\).  Substitution of the definitions gives the
current threshold.  Writing the solution as
\(y=y_{\rm hom}+\sigma_0R\), a turning point requires
\(\dot y_{\rm hom}+\sigma_0\dot R=0\); minimizing over times with
\(\dot R>0\) gives the exact threshold functional.
\end{proof}

\begin{corollary}[Existence of an exchange-driven crossing]
\label{cor:exchange-driven-crossing}
Under the hypotheses of \Cref{thm:irreversible-crossing-threshold}, suppose the
source is above threshold, the solution remains in the linearized basin, and the
decaying source is asymptotically negligible.  Then \(\dot B>0\) for sufficiently
small positive time, while \(B(t)\to B_*\).  Hence \(B\) has at least one later
turning point \(t_c>0\), and the reconstructed equation of state crosses
\(w=-1\) there.  A transverse crossing satisfies
\[
  \dot B(t_c)=0,
  \qquad
  \ddot B(t_c)\ne0.
\]
\end{corollary}

\begin{proof}
The threshold theorem gives \(\dot B>0\) initially.  Stability of the fractional
linear equation with an exponentially decaying source gives
\(B(t)\to B_*\).  A function that initially rises above its initial value but
converges to the lower equilibrium value must cease rising; continuity of the
derivative for \(t>0\) gives a turning point.  Nonzero second derivative is the
nondegeneracy condition for a transverse sign change rather than tangency.
\end{proof}


\subsection{Spectral one-crossing theorem}
\label{subsec:fit-free-crossing-topology}

The threshold theorem proves when a phantom interval can open, but it does not
by itself bound the number of later turning points.  For an ordinary local clock,
a decreasing source and an increasing restoring force are enough.  A hereditary
Caputo law is subtler: complete monotonicity makes each relaxation kernel a
positive mixture of exponentials, but an arbitrary signed superposition of many
exchange and restoring modes need not be variation diminishing.  The correct
fit-free statement is therefore spectral.

\begin{proposition}[Local-clock one-crossing benchmark]
\label{prop:local-clock-one-crossing}
Let
\[
  \dot B=S(t)-F(B),
\]
where \(S\in C^1\) is strictly decreasing and \(F\in C^1\) is strictly
increasing.  Then \(\dot B\) has at most one zero.  If it has one, it changes
sign from positive to negative.
\end{proposition}

\begin{proof}
At a critical point \(t_c\),
\[
  \ddot B(t_c)=S'(t_c)-F'(B(t_c))\dot B(t_c)=S'(t_c)<0.
\]
Thus every zero of \(\dot B\) is a strict positive-to-negative crossing.  Two
such zeros are impossible, since between them the derivative would have to
cross from negative back to positive at another zero.
\end{proof}

\begin{definition}[Single-crossover derivative spectrum]
\label{def:single-crossover-spectrum}
A linearized hereditary trajectory has a \emph{single-crossover derivative
spectrum} if, for \(t>0\),
\[
  \dot B(t)=P(t)-N(t),
  \qquad
  P(t)=\int_{[r_*,\infty)}e^{-rt}\,\nu_+(dr),
  \qquad
  N(t)=\int_{[0,r_*]}e^{-rt}\,\nu_-(dr),
\]
for nonzero positive Borel measures \(\nu_\pm\), with the exponentially tilted
first moments finite for every \(t>0\).  The supports are assumed strictly
ordered in the measure-theoretic sense,
\[
  \nu_+([0,r_*])=0,
  \qquad
  \nu_-([r_*,\infty))=0,
\]
so the positive exchange-dominated modes have rates \(r>r_*\), whereas the
negative restoring modes have rates \(r<r_*\).
\end{definition}

The representation is natural whenever the derivative of the linearized response
can be decomposed into a positive exchange component and a positive restoring
component, each resolved through the exponential spectral measure of a completely
monotone renewal resolvent.  Bernstein's theorem supplies the positive
exponential-mixture representation of each such component
\cite{SchillingSongVondracek}.  The nontrivial microscopic statement is the
\emph{ordering} of the positive and negative spectral supports.  It is stronger
than monotone decay of the total source and can be checked directly from the
joint matter--deficiency transfer spectrum.

\begin{theorem}[Fractional spectral one-crossing theorem]
\label{thm:fractional-one-crossing}
If \(\dot B\) has the single-crossover derivative spectrum of
\Cref{def:single-crossover-spectrum}, then \(\dot B\) has at most one zero on
\((0,\infty)\).  If \(\dot B>0\) at some early time and \(\dot B<0\) at some
later time, there is exactly one zero \(t_c\), and the sign change is
\[
  \dot B>0\quad\longrightarrow\quad\dot B<0.
\]
No amplitudes or cosmological fit parameters enter this conclusion.
\end{theorem}

\begin{proof}
Where \(P,N>0\), set \(R(t)=P(t)/N(t)\).  Differentiation under the integrals
gives
\[
 \frac{d}{dt}\log R(t)
 =-\frac{\int r e^{-rt}\nu_+(dr)}{\int e^{-rt}\nu_+(dr)}
  +\frac{\int r e^{-rt}\nu_-(dr)}{\int e^{-rt}\nu_-(dr)}.
\]
The first quotient is a weighted mean of rates strictly larger than \(r_*\),
whereas the second is a weighted mean of rates strictly smaller than \(r_*\).
Hence \(d(\log R)/dt<0\).  Since
\(\dot B=N(R-1)\), zeros of \(\dot B\) are exactly solutions of \(R=1\), and a
strictly decreasing function crosses one at most once.  If the two stated signs
occur, continuity gives existence and strict monotonicity of \(R\) gives
uniqueness and the displayed direction.
\end{proof}

\begin{corollary}[Fit-free crossing direction]
\label{cor:fit-free-crossing-direction}
Assume \(B,H>0\) and use the reconstructed equation of state
\[
  w_{\rm rec}+1=-\frac{2\dot B}{dHB}.
\]
Under \Cref{thm:fractional-one-crossing}, any crossing is unique and necessarily
\[
  \boxed{\qquad w_{\rm rec}<-1\quad\longrightarrow\quad
  w_{\rm rec}>-1.\qquad}
\]
A nonphantom-to-phantom crossing and an oscillatory sequence of crossings are
excluded within the single-crossover spectral class.
\end{corollary}

\begin{proof}
Before the unique turning point \(\dot B>0\), so
\(w_{\rm rec}+1<0\); afterwards \(\dot B<0\), so
\(w_{\rm rec}+1>0\).  The physical pressure ratio in an interacting
split contains the additional explicit exchange term of
\Cref{prop:deficiency-eos-identities}; the corollary concerns the expansion
history reconstructed from \(\rho_B\propto B^2\).
\end{proof}

\begin{remark}[Why monotone dilution alone is insufficient]
\label{rem:monotone-source-not-enough}
For the local equation of \Cref{prop:local-clock-one-crossing}, source
monotonicity directly controls the sign of \(\ddot B\) at a turning point.  A
Caputo equation retains the entire earlier history, so the instantaneous source
does not determine that sign.  Several separated exchange modes need not satisfy the ordered ratio
monotonicity, so source monotonicity alone does not exclude additional extrema.  The rigorous fit-free target is therefore to
derive \Cref{def:single-crossover-spectrum} from the microscopic joint Perron
spectrum, not to assume that complete monotonicity alone is variation
diminishing.
\end{remark}

\begin{theorem}[Entropy-production floor for phantom crossing]
\label{thm:crossing-entropy-floor}
In the long-time current regime of
\Cref{constr:dark-energy-ring}, let
\(\operatorname{Var}(\dot Q)=k+k'\).  If the linearized crossing threshold is
satisfied, then the thermodynamic uncertainty relation implies
\[
 \boxed{
 \sigma_{\rm ep}\ge
 \frac{2(\Gamma m^2\Delta_0)^2}
 {\kappa_{\rm ex}^2\operatorname{Var}(\dot Q)}.}
\]
In particular, zero entropy production forbids an exchange-driven crossing.
\end{theorem}

\begin{proof}
The steady-state thermodynamic uncertainty relation gives
\[
  \frac{\dot Q^2}{\operatorname{Var}(\dot Q)}
  \le\frac{\sigma_{\rm ep}}{2}
\]
in units with Boltzmann's constant one
\cite{BaratoSeifert2015,Gingrich2016,HorowitzGingrich2020}.  The threshold
\(\kappa_{\rm ex}|\dot Q|>\Gamma m^2\Delta_0\) then yields the displayed lower
bound.
\end{proof}

\subsection{Coupled background equations}
\label{subsec:coupled-dark-energy-cosmology}

The threshold theorem is local in the exchange amplitude.  A cosmological
trajectory requires the matter dilution, the Friedmann constraint, and the
conversion between an energy-transfer current and the scalar source in the
fractional budget law.

\begin{hypothesis}[Minimal coupled completion]
\label{hyp:coupled-deficiency-completion}
On a spatially flat FLRW background in \(d\) spatial dimensions, assume dustlike
matter and the system
\begin{align}
 \dot\rho_{\rm m}+dH\rho_{\rm m}&=-\mathcal Q,
 \\
 C\,\prescript{\mathrm C}{}{D}_t^{1/3}B
 &=-\Gamma I_B'(B)+\kappa_E\mathcal Q,
 \\
 H^2&=\frac{16\pi\Gren}{d(d-1)}
       \bigl(\rho_{\rm m}+c_BB^2\bigr),
\end{align}
where \(\kappa_E\) converts the energy-transfer current into the normalized
homogeneous deficiency source.  The physical deficiency pressure is defined by
total stress-energy conservation as in
\Cref{prop:deficiency-eos-identities}(ii).
\end{hypothesis}


\subsection{Matter depletion and growth comparison}
\label{subsec:matter-depletion-growth}

\begin{theorem}[Exact comoving-matter depletion]
\label{thm:comoving-matter-depletion}
For the matter equation in \Cref{hyp:coupled-deficiency-completion}, if
\(\mathcal Q(t)\ge0\), then
\[
  \boxed{\quad
  a(t)^d\rho_{\rm m}(t)
  =a_i^d\rho_{{\rm m},i}
   -\int_{t_i}^{t}a(s)^d\mathcal Q(s)\,ds.
  \quad}
\]
Consequently, relative to noninteracting dust with the same expansion history
and the same initial density,
\[
  \rho_{\rm m}(t)\le
  \rho_{\rm m}^{(0)}(t)
  :=\rho_{{\rm m},i}\left(\frac{a_i}{a(t)}\right)^d,
\]
with strict inequality after any interval of positive exchange.  Thus every
matter-to-deficiency episode leaves an exact, sign-definite comoving matter
deficit, independently of the exchange amplitude or its time profile.
\end{theorem}

\begin{proof}
Multiplying
\(\dot\rho_{\rm m}+dH\rho_{\rm m}=-\mathcal Q\) by \(a^d\) gives
\[
  \frac{d}{dt}(a^d\rho_{\rm m})=-a^d\mathcal Q.
\]
Integration proves the identity and the comparison.
\end{proof}

\begin{corollary}[Matter deficit accompanying an exchange-driven phantom phase]
\label{cor:phantom-matter-deficit}
Suppose the only positive pump in the homogeneous deficiency equation is
\(\kappa_E\mathcal Q\), with \(\kappa_E>0\) and \(\mathcal Q\ge0\).  Any
exchange-driven interval that opens the reconstructed phantom branch therefore
occurs on a background with
\[
  \Delta M_{\rm m}(t)
  :=a_i^d\rho_{{\rm m},i}-a(t)^d\rho_{\rm m}(t)
  =\int_{t_i}^{t}a(s)^d\mathcal Q(s)\,ds>0
\]
after the exchange has begun.  The statement is integrated: fractional memory
may keep \(\dot B>0\) after the instantaneous source has started to decay.
\end{corollary}

\begin{hypothesis}[Minimal smooth geodesic perturbation closure]
\label{hyp:minimal-growth-closure}
In four spacetime dimensions, on subhorizon linear scales, assume that the
deficiency component is smooth, the transfer four-vector is parallel to the
matter four-velocity, and the perturbation of the specific transfer rate is
homogeneous, so that there is no momentum drag and the growing matter mode
obeys
\[
  \ddot D+2H\dot D-4\pi\Gren\rho_{\rm m}D=0.
\]
This is the minimal closure used only for the growth-sign comparison below.
Interacting-fluid perturbations are otherwise sensitive to the choices of
\(\delta\mathcal Q\), momentum transfer, sound response, and gauge
\cite{ValiviitaMajerottoMaartens,ClemsonEtAl}.
\end{hypothesis}

\begin{theorem}[Growth suppression at fixed expansion history]
\label{thm:growth-sign-consistency}
Assume \Cref{hyp:minimal-growth-closure} on a compact time interval, with
\(H\) and the two matter densities continuous and nonnegative.  Compare the
exchange model with a noninteracting reference having the same \(a(t)\),
\(H(t)\), initial matter density, and nonnegative growing-mode initial data
\(D_i>0\), \(\dot D_i\ge0\).  If \(\mathcal Q\ge0\), then
\[
  0<D_{\mathcal Q}(t)\le D_0(t),
  \qquad
  0\le\dot D_{\mathcal Q}(t)\le\dot D_0(t).
\]
For equal primordial normalization this implies
\[
  \boxed{\qquad
  (f\sigma_8)_{\mathcal Q}(t)\le(f\sigma_8)_0(t)
  \qquad}
\]
at fixed expansion history, with strict inequality after a nonzero transfer on
an interval that contributes to the growing mode.
\end{theorem}

\begin{proof}
By \Cref{thm:comoving-matter-depletion},
\(0\le\rho_{\rm m}^{\mathcal Q}\le\rho_{\rm m}^{0}\).  Put
\(A(t)=4\pi\Gren\rho_{\rm m}(t)\) and
\(M(t)=\exp(2\int_{t_i}^{t}H(u)du)=[a(t)/a_i]^2\).  The growth equation is
equivalent to the positive-kernel Volterra equation
\[
 D_A(t)=D_i+\dot D_i\int_{t_i}^{t}M(u)^{-1}du
 +\int_{t_i}^{t}K(t,s)A(s)D_A(s)\,ds,
\]
where
\[
 K(t,s):=M(s)\int_s^tM(u)^{-1}du\ge0.
\]
Picard iteration from the common nonnegative free solution preserves order:
if \(A_0\ge A_{\mathcal Q}\), then every reference iterate dominates the
corresponding exchange iterate.  Passage to the unique Volterra solutions gives
\(D_0\ge D_{\mathcal Q}>0\).  Moreover
\[
 \dot D_A(t)=M(t)^{-1}\left[
 M(t_i)\dot D_i+\int_{t_i}^{t}M(s)A(s)D_A(s)\,ds\right],
\]
so the same ordering gives \(\dot D_0\ge\dot D_{\mathcal Q}\ge0\).  With common
primordial normalization, \(\sigma_8\propto D\) and
\(f\sigma_8\propto dD/d\log a=\dot D/H\); the fixed-\(H\) comparison proves the
last inequality.
\end{proof}

\begin{remark}[Sharp scope of the growth statement]
\label{rem:growth-scope}
The comoving-matter deficit is unconditional at background level.  The sign of
the perturbative growth correction is not unconditional until the joint
matter--deficiency perturbation map is derived: momentum exchange, a clustered
deficiency mode, or a nontrivial \(\delta\mathcal Q\) changes the growth
equation and can compete with the density deficit.  Thus
\Cref{thm:growth-sign-consistency} is a fit-free comparison theorem for the
minimal smooth geodesic closure, not a theorem for every interacting completion.
\end{remark}

\subsection{Expansion-current fluctuation symmetry}
\label{subsec:expansion-fluctuation-theorem}

The homogeneous deficiency is a Perron current rather than a scalar-field
coordinate.  When the microscopic exchange channel satisfies local detailed
balance, its rare positive and negative expansion fluctuations obey an exact
current symmetry.

\begin{definition}[Single-affinity renewal local detailed balance]
\label{def:renewal-local-detailed-balance}
Let \(b(x,y)=-b(y,x)\) be the signed deficiency increment and let \(\pi\) be the
stationary law of the embedded predictive states.  Assume common forward and
reverse support and a single thermodynamic affinity \(\mathcal A_B\) such that,
for every allowed depth \(n\),
\[
  \log\frac{\pi_xQ_{xy}(n)}{\pi_yQ_{yx}(n)}
  =\mathcal A_B b(x,y).
\]
This includes the direction--time-independence condition of thermodynamically
consistent Markov-renewal processes.  If several independent affinities are
present, the definition is applied vectorially; a scalar marginal symmetry
requires that the deficiency current carry the full antisymmetric action.
\end{definition}

\begin{theorem}[Gallavotti--Cohen symmetry of the physical-depth pressure]
\label{thm:renewal-gc-symmetry}
Under \Cref{def:renewal-local-detailed-balance}, the deficiency pressure of
\Cref{def:deficiency-rate-function} satisfies
\[
  \boxed{\qquad
  \Lambda_B(\chi)=\Lambda_B(-\mathcal A_B-\chi).
  \qquad}
\]
Its rate function obeys
\[
  \boxed{\qquad
  I_B(j)-I_B(-j)=-\mathcal A_B j.
  \qquad}
\]
Consequently, for the integrated signed current \(J_T\),
\[
 \lim_{T\to\infty}\frac1T
 \log\frac{\mathbb P(J_T/T\simeq j)}
          {\mathbb P(J_T/T\simeq-j)}
 =\mathcal A_Bj.
\]
If the homogeneous calibration identifies accumulated expansion with
\[
  \mathcal N_T:=\int_0^TH(t)dt=\frac{J_T}{d},
\]
then
\[
  I_{\mathcal N}(n)-I_{\mathcal N}(-n)
  =-d\mathcal A_Bn,
\]
and the mean entropy-production rate is
\[
  \sigma_{\rm ep}=\mathcal A_B B_*\ge0.
\]
\end{theorem}

\begin{proof}
Let \(D_\pi=\operatorname{diag}(\pi_x)\).  The local detailed-balance identity
and antisymmetry of \(b\) give, entry by entry,
\[
 D_\pi\Qhat(z,\chi,0,0)D_\pi^{-1}
 =\Qhat(z,-\mathcal A_B-\chi,0,0)^{\mathsf T}.
\]
The two matrices are similar up to transposition and therefore have the same
spectral radius.  The defining Perron equation
\(\rho[\Qhat(\Lambda_B(\chi),\chi,0,0)]=1\), together with uniqueness of the
physical-depth branch, gives the pressure symmetry.  Legendre duality then
implies
\[
 I_B(-j)=\mathcal A_Bj+I_B(j),
\]
which is the rate-function and probability-ratio statement.  The e-fold formula
follows from \(J_T=d\mathcal N_T\).  Finally, convexity of a pressure symmetric
about \(-\mathcal A_B/2\), with zeros at \(0\) and \(-\mathcal A_B\), makes
\(\mathcal A_B\Lambda_B'(0)=\mathcal A_BB_*\ge0\).  This is the standard
Markov and semi-Markov current fluctuation theorem expressed in the present
matrix-renewal normalization
\cite{LebowitzSpohn,AndrieuxGaspard2008,FaggionatoSemiMarkov}.
\end{proof}

\begin{corollary}[Exact expansion-current cumulant hierarchy]
\label{cor:expansion-cumulant-hierarchy}
Let
\[
  C_n:=\partial_\chi^n\Lambda_B(0)
\]
be the connected deficiency-current cumulants per physical depth.  If \(\Lambda_B\) extends holomorphically to a complex neighbourhood of the
closed disk \(\{\chi:|\chi|\le|\mathcal A_B|\}\), then
\[
  0=\Lambda_B(-\mathcal A_B)
   =\sum_{n\ge1}\frac{(-\mathcal A_B)^n}{n!}C_n,
\]
so
\[
  \boxed{
  C_1=\frac{\mathcal A_B}{2}C_2
      -\frac{\mathcal A_B^2}{6}C_3
      +\frac{\mathcal A_B^3}{24}C_4-\cdots .}
\]
Thus
\[
  \sigma_{\rm ep}=\mathcal A_BC_1
  =\frac{\mathcal A_B^2}{2}C_2
   -\frac{\mathcal A_B^3}{6}C_3+\cdots .
\]
For the accumulated e-folds, the cumulant rates are
\(C_n^{\mathcal N}=C_n/d^n\), and hence
\[
 \frac{d}{dT}\mathbb E\mathcal N_T
 =\frac{d\mathcal A_B}{2}C_2^{\mathcal N}
  -\frac{d^2\mathcal A_B^2}{6}C_3^{\mathcal N}
  +\frac{d^3\mathcal A_B^3}{24}C_4^{\mathcal N}-\cdots .
\]
At the reversible locus \(\mathcal A_B=0\), all odd current and e-fold
cumulants vanish.
\end{corollary}

\begin{proof}
The pressure symmetry gives \(\Lambda_B(-\mathcal A_B)=\Lambda_B(0)=0\).
Taylor expansion at the origin and rearrangement give the hierarchy.  Cumulants
of \(J_T/d\) scale by \(d^{-n}\).  At zero affinity the pressure is even, so all
odd derivatives at the origin vanish.
\end{proof}

\begin{remark}[What the fluctuation theorem predicts]
\label{rem:gc-scope}
The symmetry fixes probability ratios and nonlinear cumulant identities without
fitting a \(w_0\)--\(w_a\) curve.  It does not fix the affinity
\(\mathcal A_B\), the selected stationary current \(B_*\), or the present
cosmological scale.  In a multi-affinity matter sector the exact theorem is the
joint vector symmetry; projecting to the expansion current alone is valid only
when the other affinities vanish or have been retained in the joint tilt.
\end{remark}

\begin{theorem}[Renewal dark-energy decision theorem]
\label{thm:renewal-dark-energy-dichotomy}
Within the registers and closure assumptions above:
\begin{enumerate}[label=\textup{(\roman*)},leftmargin=2.6em]
\item a closed vacuum deficiency is exactly constant and has \(w=-1\);
\item a reversible source-free effective deficiency fluid relaxes toward \(B_*\)
      without crossing, with the critical algebraic tail
      \(w+1\sim t^{-4/3}\) in the quadratic basin;
\item an irreversible exchange channel can open a reconstructed phantom interval
      only above the current threshold of
      \Cref{thm:irreversible-crossing-threshold}, and every such opening requires
      the entropy-production floor of \Cref{thm:crossing-entropy-floor};
\item if the linearized derivative spectrum has the ordered single-crossover
      form of \Cref{def:single-crossover-spectrum}, then there is at most one
      crossing and its direction is necessarily
      \(w_{\rm rec}<-1\to w_{\rm rec}>-1\);
\item positive matter-to-deficiency exchange produces the exact comoving-matter
      deficit of \Cref{thm:comoving-matter-depletion}; under
      \Cref{hyp:minimal-growth-closure} it also suppresses \(D\), \(\dot D\), and
      \(f\sigma_8\) relative to a reference with the same expansion history;
\item single-affinity local detailed balance gives the fluctuation symmetry and
      cumulant hierarchy of
      \Cref{thm:renewal-gc-symmetry,cor:expansion-cumulant-hierarchy};
\item numerical values of \(w_0\), \(w_a\), \(z_{\rm cross}\), and the present
      dark-energy density still depend on the selected Perron state, the
      matter--deficiency projection, the relaxation clock in Hubble units, and
      the cosmological background.
\end{enumerate}
\end{theorem}

\begin{proof}
Items (i)--(iii) are respectively
\Cref{thm:closed-source-free-deficiency-nogo,thm:source-free-no-crossing,thm:irreversible-crossing-threshold,thm:crossing-entropy-floor}.
Item (iv) is
\Cref{thm:fractional-one-crossing,cor:fit-free-crossing-direction}; item (v) is
\Cref{thm:comoving-matter-depletion,thm:growth-sign-consistency}; and item (vi)
is \Cref{thm:renewal-gc-symmetry,cor:expansion-cumulant-hierarchy}.  Item (vii)
follows because none of those structural theorems fixes the stationary current,
the exchange projection, the physical-time clock conversion, or the initial
cosmological state.
\end{proof}

\begin{remark}[The trajectory is generically non-CPL]
The solution of \Cref{hyp:coupled-deficiency-completion} need not have the
Chevallier--Polarski--Linder form \(w(a)=w_0+w_a(1-a)\).  In the
single-crossover spectral class it has one directed transition and subsequent
relaxation, but its detailed shape is not linear in \(1-a\).  Parameters
\((w_0,w_a)\) should therefore be regarded as range-fit summaries of a renewal
trajectory, not as fundamental coordinates of the model.  Outside the ordered
spectral class, fractional memory alone does not exclude additional extrema.
\end{remark}

\begin{remark}[No numerical crossing redshift is claimed]
The present paper proves a mechanism, a threshold, and an entropy-production
floor.  It does not derive a numerical \(z_{\rm cross}\), and in particular does
not identify the finite \(K_4\) orientation angle with a crossing redshift.  A
number such as \(z\sim0.3\) can only enter after a specified exchange law is
calibrated to expansion and growth data.  Deriving \(\kappa_E\) and the
Hubble-time relaxation normalization from the joint matter--deficiency Perron
operator is the remaining microscopic step.
\end{remark}

\subsection{Scope of the dark-energy predictions}
\label{subsec:dark-energy-status}

\begin{proposition}[Local finite-window operators cannot drive late dark energy]
\label{prop:late-finite-window-null}
On a slowly varying FLRW background within the controlled derivative window,
the order-\(\eps^2\) local curvature-squared contribution is smaller than the
Einstein term by
\[
  O\!\left(\eps^2H^2,\eps^2|\dot H|\right).
\]
Consequently an order-one late-time departure from \(w=-1\) cannot be generated
by the ultraviolet finite-window operators while the renewal expansion remains
controlled.  It must arise from the homogeneous budget or its exchange law.
\end{proposition}

\begin{proof}
On FLRW, \(R=O(H^2,\dot H)\).  Variation of a term
\(\eps^2R^2\) is therefore of order
\(\eps^2(H^4,H^2\dot H,\ldots)\), whereas the Einstein tensor is of order
\(H^2,\dot H\).  Their ratio is the displayed derivative-expansion parameter.
The same estimate holds for the Ricci-squared term, while the Weyl response
vanishes on exact FLRW.  Hence a macroscopic effect from these operators would
coincide with loss of the controlled expansion rather than constitute a
late-time prediction within it.
\end{proof}

\begin{center}
\renewcommand{\arraystretch}{1.17}
\begin{tabular}{p{0.30\textwidth}p{0.42\textwidth}p{0.20\textwidth}}
\toprule
\textbf{Statement} & \textbf{Content} & \textbf{Status}\\
\midrule
Vacuum endpoint & \(B=\mathrm{constant}\), \(w=-1\) & proved by Bianchi identity\\
Perron potential & \(V_{\rm eff}=I_B\) & proved under the LDP hypotheses\\
Critical clock & Caputo order \(\alpha=1/3\) & renewal-clock constitutive input\\
Reversible branch & monotone relaxation, no crossing & proved under convexity/comparison hypotheses\\
Irreversible branch & phantom opening only above a current threshold & proved in the linearized exchange model\\
Entropy floor & nonzero minimum entropy production & proved from the TUR\\
Crossing topology & at most one, phantom-to-nonphantom & proved under ordered derivative spectrum\\
Matter background & exact comoving matter depletion for \(\mathcal Q\ge0\) & unconditional identity\\
Linear growth & suppression at fixed \(H(t)\) & proved under minimal smooth geodesic closure\\
Expansion fluctuations & Gallavotti--Cohen symmetry and cumulant hierarchy & proved under single-affinity local detailed balance\\
Crossing redshift & model/background dependent & open calibration problem\\
General perturbations and sound response & not fixed by the homogeneous closure & open\\
\bottomrule
\end{tabular}
\end{center}

\begin{remark}[Highest-leverage remaining calculation]
The decisive next step is to derive from the joint matter--deficiency tilted
transfer operator: \textup{(i)} the scalar projection \(\kappa_E\); \textup{(ii)}
the physical-time relaxation normalization; \textup{(iii)} the perturbation and
momentum-transfer kernel; and \textup{(iv)} the signed relaxation spectrum of
\(\dot B\).  The last item decides whether the single-crossover hypothesis is a
theorem of the microscopic channel.  The same operator should also compute the
affinity \(\mathcal A_B\) entering the expansion fluctuation theorem.  Closing
these quantities would turn the present fit-free topology and fluctuation
relations into a parameter-reduced trajectory for \(H(z)\), reconstructed
\(w(z)\), and growth.
\end{remark}

\section{Discussion and open problems}
\label{sec:interpretation}
% =====================================================================

The results of this paper are most naturally read as a hierarchy of responses of one microscopic transfer process.  The Lorentzian principal symbol, the proper-time action, the Ricci response, the screen entropy density, the homogeneous cosmological current, and the finite-resolution Wilson data are not introduced as independent ingredients.  They arise at different derivative orders, under different tilts, or in different asymptotic regimes of the same renewal dynamics.  This does not make the corresponding observables numerically identical; on the contrary, the construction keeps their distinct statistical meanings explicit.  Its unifying claim is that geometry, gravitational coupling, vacuum curvature, dark-energy evolution, and higher-curvature response can be traced to one common microscopic object without collapsing them into a single parameter.

\subsection{A common origin for gravitational scales}

A particularly significant consequence is the relation between quantities that are historically treated as unrelated inputs.  In general relativity Newton's constant fixes the response of geometry to stress energy, whereas the cosmological constant fixes the curvature of the vacuum.  In the standard effective description they are independent couplings, and the small observed vacuum scale is the subject of the cosmological-constant problem \cite{WeinbergCC}.  In the renewal construction the two constants remain physically distinct, but they are generated by different responses of the same tilted Perron family.  The Newton coupling is fixed by predictive entropy per unit screen area, while the de~Sitter curvature is fixed by the homogeneous signed current per physical renewal depth.  The metric scale is supplied by the displacement covariance, and the higher-curvature coefficients are supplied by higher cumulants.  Thus the usual gravitational constants occupy different components of one microscopic response hierarchy.

This common origin is stronger than a dimensional relation.  It identifies which microscopic observable controls each continuum quantity and why no direct equality between current, covariance, and entropy should be assumed.  In the normalized predictive-pure $K_4$ specialization, the technical results also show that the explicit renewal length cancels from the dimensionless product of the Newton and cosmological constants in three spatial dimensions.  What remains is the signed microscopic current and the predictive entropy.  The cancellation is therefore a structural normalization result, not yet a numerical derivation of the observed cosmological constant: the selected Perron state and its signed current still have to be determined from the microscopic cosmological sector.  Even with that limitation, placing the strength of gravity and the curvature of the vacuum inside one response object is one of the main conceptual outcomes of the construction.

The same point clarifies the status of the de~Sitter branch.  The cosmological term is not appended to an already derived Einstein equation as an arbitrary constant.  It is fixed in the homogeneous pure-deficiency sector by the physical renewal current, and the flat slicing is selected by the Hamiltonian constraint rather than chosen as an ansatz.  The four-derivative response preserves that selection throughout the controlled regime.  The local positive-curvature result is intrinsic; only the passage to complete global de~Sitter space uses the additional predictive-completion principle.

\subsection{Dark energy in relation to \texorpdfstring{$\Lambda$}{Lambda}CDM and dynamical alternatives}

At the background level, the equilibrium renewal branch reproduces the defining behavior of the cosmological constant in $\Lambda$CDM: the vacuum contribution is constant, has equation of state $w=-1$, and supports de~Sitter expansion.  The difference is interpretive and microscopic.  In $\Lambda$CDM the cosmological constant is a phenomenological parameter of the gravitational action.  Here the constant branch is the equilibrium value of a signed renewal current, and covariance explains why a closed vacuum realization cannot roll.  The model therefore does not replace the successful background structure of $\Lambda$CDM merely for the sake of dynamical freedom; it supplies a proposed microscopic origin for that structure and identifies the conditions under which departures from it can occur.

The closest conventional alternatives are scalar-field models.  Quintessence and cosmon models replace a constant vacuum term by a slowly rolling field with a chosen potential \cite{RatraPeebles1988,Wetterich1988,CaldwellDaveSteinhardt}.  K-essence instead uses nonlinear kinetic structure to obtain tracking or late acceleration \cite{ArmendarizPiconKessence}.  Renewal dark energy differs from both because the homogeneous variable is not postulated as a fundamental scalar field and its restoring law is not selected by choosing a potential or a kinetic function.  The equilibrium point and local restoring potential come from the Perron large-deviation function, while the fractional time law comes from the renewal clock.  In the reversible convex branch this produces a directed algebraic relaxation toward $w=-1$ and forbids a crossing.  The memory exponent and the no-crossing statement are therefore consequences of the renewal closure rather than model-building choices in a scalar Lagrangian.

The comparison with phantom models is sharper.  A persistent component with $w<-1$ is often represented by an exotic or wrong-sign degree of freedom and can lead to severe stability questions or a future big-rip singularity \cite{CaldwellPhantom}.  In the renewal framework a phantom value is not a new equilibrium phase and does not require a permanently phantom microscopic field.  It is a reconstructed transient produced when an irreversible exchange current temporarily drives the homogeneous deficiency away from its restoring branch.  Such an interval opens only above a current threshold and carries a nonzero entropy-production cost.  Under the ordered single-crossover spectral condition it can occur at most once and must end on the nonphantom side.  The result is therefore closer to a nonequilibrium excursion than to a phantom vacuum.

Interacting dark-energy models provide the nearest phenomenological analogue because they allow energy transfer between matter and the accelerating component \cite{Amendola2000,ValiviitaMajerottoMaartens,ClemsonEtAl}.  The renewal construction shares their basic observational logic: transfer modifies the matter abundance, the expansion history, and the growth of structure, and the perturbation sector depends on the momentum-transfer and sound-response closure.  Its distinguishing feature is that the exchange is represented by an oriented finite-cycle current inside the same Perron system that controls the homogeneous expansion.  This makes irreversibility measurable through entropy production and current fluctuations.  Positive matter-to-deficiency transfer leaves an exact comoving matter deficit, and under the minimal smooth geodesic closure it suppresses linear growth at fixed expansion history.  Single-affinity local detailed balance further yields a fluctuation symmetry and cumulant relations for accumulated expansion.  These relations have no direct counterpart in a generic phenomenological coupling chosen only at the level of background continuity equations.

The finite-resolution gravitational terms should not be confused with a modified-gravity explanation of late acceleration.  Inside the controlled derivative window their effect on an FLRW background is parametrically smaller than the Einstein term.  An order-one late-time departure from $w=-1$ must therefore come from the homogeneous current and its exchange law, not from the local curvature-squared operators derived in this paper.  The higher-curvature sector instead records how non-Gaussian renewal statistics first become visible once the continuum is probed at finite resolution.

The present dark-energy conclusions are deliberately structural.  They distinguish an equilibrium cosmological constant from reversible relaxation and irreversible exchange; they determine the allowed crossing direction under a spectral hypothesis; and they connect a phantom excursion with matter depletion, growth suppression, entropy production, and current fluctuations.  They do not yet determine the present dark-energy density, a numerical crossing redshift, or a unique trajectory in the usual $w_0$--$w_a$ plane.  Those quantities require the physical-time normalization, the matter--deficiency projection, and the selected Perron state of the joint microscopic channel.

\subsection{Gaussian universality and finite-resolution information}

The analysis also suggests a useful universality principle.  At leading order, models with the same massive-band Hamiltonian, metric covariance, screen entropy density, and modular normalization produce the same proper-time propagation, Ricci reconstruction, Einstein equation, and Newton coupling.  In this sense Einstein dynamics belongs to a Gaussian renewal universality class.  Microscopic distinctions first reappear through non-Gaussian cumulants at finite resolution.  The explicit isotropic compound-Poisson calculation illustrates this separation: its curvature-squared response is not a pure scalar correction but a fixed ray containing scalar, traceless-Ricci, and subdominant Weyl components.  The Wilson ray is a property of the specified channel, while the operator basis is universal within the stated symmetry and derivative assumptions.

This hierarchy is conceptually different from adding arbitrary higher-curvature terms to a gravitational action.  The coefficients are response data of the same process whose quadratic statistics generated the Einstein layer.  Coarse-graining suppresses normalized higher cumulants while preserving the covariance, explaining why the Einstein sector is stable and why finite-resolution corrections retain increasingly detailed information about the underlying renewal statistics.  The trans-renewal boundary marks the point at which that hierarchy ceases to be ordered.  It is a boundary of the continuum approximation, not evidence by itself for a bounce, a regular core, or any particular ultraviolet continuation.

\subsection{Scope of the construction}

Several parts of the argument are theorem-level within clearly stated classes.  The foundational Lorentzian operator limit is imported from the kinematic renewal construction, while its stability under the reset term, the isolated massive band, the matrix WKB propagation, the admissible-class subprincipal closure, the two-profile Ricci reconstruction, the primitive-screen area law, the homogeneous flat de~Sitter branch, and the finite-resolution operator basis are established here.  The modular normalization is proved for the stated quasi-free screen class, not for arbitrary interacting predictive screens.  The complete global de~Sitter interpretation remains conditional on the predictive-completion principle.  The dark-energy crossing and growth statements likewise retain their explicit spectral and perturbative closure assumptions.

The main unresolved microscopic problem is to derive the full joint matter--deficiency transfer operator in physical time.  That calculation should determine the exchange projection, the fractional-clock normalization in Hubble units, the perturbation transfer four-vector, the sound and momentum response, and the signed relaxation spectrum.  It would decide whether the ordered single-crossover condition follows from the microscopic channel rather than being an independently checkable spectral hypothesis.  The same joint tilt should determine whether a scalar expansion-current fluctuation theorem survives or must be replaced by a multi-affinity symmetry.

Further extensions concern interacting modular screens, the absolute geometric area represented by one predictive cell, the complete anisotropic and parity-odd fourth-cumulant response, and the induced-action normalization outside the factorized heat-kernel scheme.  At higher derivative order one must separate universal propagation effects from genuinely new connected cumulants and determine which nonlocal causal kernels preserve the curvature-locking argument.  A microscopic global continuation theorem would also be needed to replace the predictive-completion principle.

The resulting picture is nevertheless coherent at the level reached here.  Renewal memory supplies a Lorentzian propagation sector, reconstructs curvature from fluctuations of that propagation, converts predictive entropy into gravitational dynamics, and assigns the vacuum and dynamical dark-energy sectors to the equilibrium and nonequilibrium behavior of a homogeneous signed current.  Newton's constant, the cosmological term, and the first higher-curvature coefficients thereby become distinct but related observables of one microscopic renewal system.  The continuum Einstein--de~Sitter theory is the hydrodynamic phase of that system, while dark-energy dynamics and finite-resolution corrections record how it departs from equilibrium and Gaussianity.

\begin{thebibliography}{99}\small
\bibitem{Jacobson}
T.~Jacobson, \emph{Thermodynamics of spacetime: the Einstein equation of state},
Phys. Rev. Lett. \textbf{75} (1995), 1260--1263.

\bibitem{Unruh1976}
W.~G.~Unruh,
Phys. Rev. D \textbf{14} (1976), 870--892.

\bibitem{ElingGuedensJacobson}
C.~Eling, R.~Guedens, and T.~Jacobson,
\emph{Non-equilibrium thermodynamics of spacetime},
Phys. Rev. Lett. \textbf{96} (2006), 121301.

\bibitem{Padmanabhan}
T.~Padmanabhan, \emph{Thermodynamical aspects of gravity: new insights},
Rep. Prog. Phys. \textbf{73} (2010), 046901.

\bibitem{Verlinde2011}
E.~Verlinde, \emph{On the origin of gravity and the laws of Newton},
JHEP \textbf{04} (2011), 029.

\bibitem{Lashkari}
N.~Lashkari, M.~B.~McDermott, and M.~Van Raamsdonk,
\emph{Gravitational dynamics from entanglement ``thermodynamics''},
JHEP \textbf{2014}, 195.

\bibitem{FaulknerEtAl}
T.~Faulkner, M.~Guica, T.~Hartman, R.~C.~Myers, and M.~Van~Raamsdonk,
\emph{Gravitation from entanglement in holographic CFTs},
JHEP \textbf{03} (2014), 051.

\bibitem{Jacobson2016}
T.~Jacobson, \emph{Entanglement equilibrium and the Einstein equation},
Phys. Rev. Lett. \textbf{116} (2016), 201101.

\bibitem{BombelliLeeMeyerSorkin}
L.~Bombelli, J.~Lee, D.~Meyer, and R.~D.~Sorkin,
\emph{Space-time as a causal set},
Phys. Rev. Lett. \textbf{59} (1987), 521--524.

\bibitem{Surya}
S.~Surya, \emph{The causal set approach to quantum gravity},
Living Rev. Relativity \textbf{22} (2019), 5.

\bibitem{BenincasaDowker}
D.~M.~T.~Benincasa and F.~Dowker,
\emph{The scalar curvature of a causal set},
Phys. Rev. Lett. \textbf{104} (2010), 181301.

\bibitem{RoySinhaSurya}
M.~Roy, D.~Sinha, and S.~Surya,
\emph{The discrete geometry of a small causal diamond},
Phys. Rev. D \textbf{87} (2013), 044046.

\bibitem{AmbjornJurkiewiczLoll}
J.~Ambj{\o}rn, J.~Jurkiewicz, and R.~Loll,
\emph{Reconstructing the universe}, Phys. Rev. D \textbf{72} (2005), 064014.

\bibitem{Oriti}
D.~Oriti, \emph{The group field theory approach to quantum gravity},
in \emph{Approaches to Quantum Gravity}, Cambridge Univ. Press, 2009, 310--331.

\bibitem{Connes}
A.~Connes, \emph{Noncommutative Geometry}, Academic Press, 1994.

\bibitem{ChamseddineConnes}
A.~H.~Chamseddine and A.~Connes, \emph{The spectral action principle},
Comm. Math. Phys. \textbf{186} (1997), 731--750.

\bibitem{FrancoEckstein}
N.~Franco and M.~Eckstein,
\emph{An algebraic formulation of causality for noncommutative geometry},
Class. Quantum Grav. \textbf{30} (2013), 135007.

\bibitem{Strohmaier}
A.~Strohmaier,
\emph{On noncommutative and pseudo-Riemannian geometry},
J. Geom. Phys. \textbf{56} (2006), 175--195.

\bibitem{Barbero}
J.~F.~Barbero, \emph{Real Ashtekar variables for Lorentzian signature
space-times}, Phys. Rev. D \textbf{51} (1995), 5507--5510.

\bibitem{Ashtekar1986}
A.~Ashtekar, \emph{New variables for classical and quantum gravity},
Phys. Rev. Lett. \textbf{57} (1986), 2244--2247.

\bibitem{RovelliSmolin}
C.~Rovelli and L.~Smolin, \emph{Spin networks and quantum gravity},
Phys. Rev. D \textbf{52} (1995), 5743--5759.

\bibitem{AshtekarLewandowski}
A.~Ashtekar and J.~Lewandowski,
\emph{Projective techniques and functional integration for gauge theories},
J. Math. Phys. \textbf{36} (1995), 2170--2191.

\bibitem{EPRL}
J.~Engle, E.~Livine, R.~Pereira, and C.~Rovelli,
\emph{LQG vertex with finite Immirzi parameter},
Nucl. Phys. B \textbf{799} (2008), 136--149.

\bibitem{BarrettAsymptotics}
J.~W.~Barrett, R.~J.~Dowdall, W.~J.~Fairbairn, H.~Gomes, and F.~Hellmann,
\emph{Asymptotic analysis of the EPRL four-simplex amplitude},
J. Math. Phys. \textbf{50} (2009), 112504.

\bibitem{Immirzi}
G.~Immirzi, \emph{Real and complex connections for canonical gravity},
Class. Quantum Grav. \textbf{14} (1997), L177--L181.

\bibitem{DomagalaLewandowski}
M.~Domaga\l a and J.~Lewandowski,
\emph{Black-hole entropy from quantum geometry},
Class. Quantum Grav. \textbf{21} (2004), 5233--5243.

\bibitem{Meissner}
K.~A.~Meissner, \emph{Black-hole entropy in loop quantum gravity},
Class. Quantum Grav. \textbf{21} (2004), 5245--5251.

\bibitem{AshtekarBaezKrasnov}
A.~Ashtekar, J.~Baez, and K.~Krasnov,
\emph{Quantum geometry of isolated horizons and black hole entropy},
Adv. Theor. Math. Phys. \textbf{4} (2000), 1--94.

\bibitem{Stelle1977}
K.~S.~Stelle, \emph{Renormalization of higher-derivative quantum gravity},
Phys. Rev. D \textbf{16} (1977), 953--969.

\bibitem{Donoghue1994}
J.~F.~Donoghue,
\emph{General relativity as an effective field theory: the leading quantum
corrections}, Phys. Rev. D \textbf{50} (1994), 3874--3888.

\bibitem{Burgess2004}
C.~P.~Burgess,
\emph{Quantum gravity in everyday life: general relativity as an effective field
theory}, Living Rev. Relativity \textbf{7} (2004), 5.

\bibitem{HuVerdaguer}
B.~L.~Hu and E.~Verdaguer,
\emph{Stochastic gravity: theory and applications},
Living Rev. Relativ. \textbf{11} (2008), 3.

\bibitem{ShaliziCrutchfield}
C.~R.~Shalizi and J.~P.~Crutchfield,
\emph{Computational mechanics: pattern and prediction, structure and simplicity},
J. Stat. Phys. \textbf{104} (2001), 817--879.

\bibitem{Feller}
W.~Feller, \emph{An Introduction to Probability Theory and Its Applications,
Vol.~II}, 2nd ed., Wiley, 1971.

\bibitem{PelissierLorentzian}
A.~P\'elissier,
\emph{Renewal Spectral Geometry and the Emergence of Lorentzian Spacetime},
arXiv placeholder, 2026.

\bibitem{PelissierCurrentsK4Arxiv}
A.~P\'elissier,
\emph{Renewal Currents, Perron Response, and the
\(\mathbb Z_4\)-Lifted \(K_4\) Cycle Geometry},
arXiv placeholder, 2026.

\bibitem{Kato}
T.~Kato, \emph{Perturbation Theory for Linear Operators}, 2nd ed., Springer, 1976.

\bibitem{EvansHoeghKrohn}
D.~E.~Evans and R.~H\o egh-Krohn,
\emph{Spectral properties of positive maps on $C^*$-algebras},
J. London Math. Soc. \textbf{17} (1978), 345--355.

\bibitem{Nagaev}
S.~V.~Nagaev, \emph{Some limit theorems for stationary Markov chains}, Theory
Probab. Appl. \textbf{2} (1957), 378--406.

\bibitem{KipnisVaradhan}
C.~Kipnis and S.~R.~S.~Varadhan, \emph{Central limit theorem for additive
functionals of reversible Markov processes}, Comm. Math. Phys. \textbf{104}
(1986), 1--19.

\bibitem{LittlejohnFlynn}
R.~G.~Littlejohn and W.~G.~Flynn,
\emph{Geometric phases in the asymptotic theory of coupled wave equations},
Phys. Rev. A \textbf{44} (1991), 5239--5256.

\bibitem{WilczekZee}
F.~Wilczek and A.~Zee,
\emph{Appearance of gauge structure in simple dynamical systems},
Phys. Rev. Lett. \textbf{52} (1984), 2111--2114.

\bibitem{EmmrichWeinstein}
C.~Emmrich and A.~Weinstein,
\emph{Geometry of the transport equation in multicomponent WKB approximations},
Comm. Math. Phys. \textbf{176} (1996), 701--711.

\bibitem{Teufel}
S.~Teufel, \emph{Adiabatic Perturbation Theory in Quantum Dynamics},
Springer, 2003.

\bibitem{Berry}
M.~V.~Berry, \emph{Quantal phase factors accompanying adiabatic changes},
Proc. R. Soc. Lond. A \textbf{392} (1984), 45--57.

\bibitem{Simon}
B.~Simon, \emph{Holonomy, the quantum adiabatic theorem, and Berry's phase},
Phys. Rev. Lett. \textbf{51} (1983), 2167--2170.

\bibitem{Arnold}
V.~I.~Arnold, \emph{Mathematical Methods of Classical Mechanics}, 2nd ed.,
Springer, 1989.

\bibitem{Hartman}
P.~Hartman, \emph{Ordinary Differential Equations}, 2nd ed., SIAM, 2002.

\bibitem{WaldBook}
R.~M.~Wald, \emph{General Relativity}, University of Chicago Press, 1984.

\bibitem{Besse}
A.~L.~Besse, \emph{Einstein Manifolds}, Springer, 1987.

\bibitem{Nenciu}
G.~Nenciu, \emph{Linear adiabatic theory: exponential estimates},
Comm. Math. Phys. \textbf{152} (1993), 479--496.

\bibitem{GoldsteinKac}
S.~Goldstein, \emph{On diffusion by discontinuous movements, and on the telegraph equation},
Quart. J. Mech. Appl. Math. \textbf{4} (1951), 129--156.

\bibitem{Kac}
M.~Kac, \emph{A stochastic model related to the telegrapher's equation},
Rocky Mountain J. Math. \textbf{4} (1974), 497--509.

\bibitem{GaveauJacobsonKacSchulman}
B.~Gaveau, T.~Jacobson, M.~Kac, and L.~S.~Schulman,
\emph{Relativistic extension of the analogy between quantum mechanics and
Brownian motion}, Phys. Rev. Lett. \textbf{53} (1984), 419--422.

\bibitem{Ord1983}
G.~N.~Ord, \emph{Fractal space-time: a geometric analogue of relativistic
quantum mechanics}, J. Phys. A \textbf{16} (1983), 1869--1884.

\bibitem{McKeonOrd}
D.~G.~C.~McKeon and G.~N.~Ord,
\emph{Time reversal in stochastic processes and the Dirac equation},
Phys. Rev. Lett. \textbf{69} (1992), 3--4.

\bibitem{PST}
G.~Panati, H.~Spohn, and S.~Teufel, \emph{Space-adiabatic perturbation theory},
Adv. Theor. Math. Phys. \textbf{7} (2003), 145--204.

\bibitem{DimassiSjostrand}
M.~Dimassi and J.~Sj\"ostrand, \emph{Spectral Asymptotics in the Semi-Classical
Limit}, Cambridge Univ. Press, 1999.

\bibitem{HormanderIII}
L.~H\"ormander, \emph{The Analysis of Linear Partial Differential Operators III},
Springer, 1985.

\bibitem{KossakowskiFGV}
A.~Kossakowski, A.~Frigerio, V.~Gorini, and M.~Verri,
\emph{Quantum detailed balance and KMS condition},
Comm. Math. Phys. \textbf{57} (1977), 97--110.

\bibitem{Alicki1976}
R.~Alicki, \emph{On the detailed balance condition for non-Hamiltonian systems},
Rep. Math. Phys. \textbf{10} (1976), 249--258.

\bibitem{Cipriani}
F.~Cipriani, \emph{Dirichlet forms and Markovian semigroups on standard forms of
von Neumann algebras}, J. Funct. Anal. \textbf{147} (1997), 259--300.

\bibitem{KipnisLandim}
C.~Kipnis and C.~Landim, \emph{Scaling Limits of Interacting Particle Systems},
Springer, 1999.

\bibitem{Doob}
J.~L.~Doob, \emph{Classical Potential Theory and Its Probabilistic Counterpart},
Springer, 1984.

\bibitem{Pinsky}
R.~G.~Pinsky, \emph{Positive Harmonic Functions and Diffusion}, Cambridge
University Press, 1995.

\bibitem{Zworski}
M.~Zworski, \emph{Semiclassical Analysis}, Graduate Studies in Mathematics
\textbf{138}, American Mathematical Society, 2012.

\bibitem{CalderonVaillancourt}
A.-P.~Calder\'on and R.~Vaillancourt,
\emph{On the boundedness of pseudo-differential operators},
J. Math. Soc. Japan \textbf{23} (1971), 374--378.

\bibitem{Pavliotis}
G.~A.~Pavliotis, \emph{Stochastic Processes and Applications:
Diffusion Processes, the Fokker--Planck and Langevin Equations},
Texts in Applied Mathematics \textbf{60}, Springer, 2014.

\bibitem{LelievreStoltz}
T.~Leli\`evre and G.~Stoltz,
\emph{Partial differential equations and stochastic methods in molecular
dynamics}, Acta Numerica \textbf{25} (2016), 681--880.

\bibitem{Schrader}
R.~Schrader,
\emph{Perron--Frobenius theory for positive maps on trace ideals},
Fields Inst. Commun. \textbf{30} (2001), 361--378.

\bibitem{VisserVanVleck}
M.~Visser,
\emph{Van Vleck determinants: geodesic focusing in Lorentzian spacetimes},
Phys. Rev. D \textbf{47} (1993), 2395--2402.

\bibitem{Synge}
J.~L.~Synge, \emph{Relativity: The General Theory}, North-Holland, 1960.

\bibitem{DeWitt}
B.~S.~DeWitt, \emph{Dynamical Theory of Groups and Fields},
Gordon and Breach, 1965.

\bibitem{BuchholzLechner}
D.~Buchholz and G.~Lechner,
\emph{Modular nuclearity and localization},
Ann. Henri Poincar\'e \textbf{5} (2004), 1065--1080.

\bibitem{FannesNachtergaeleWerner}
M.~Fannes, B.~Nachtergaele, and R.~F.~Werner,
\emph{Finitely correlated states on quantum spin chains},
Comm. Math. Phys. \textbf{144} (1992), 443--490.

\bibitem{KretschmannWerner}
D.~Kretschmann and R.~F.~Werner, \emph{Quantum channels with memory},
Phys. Rev. A \textbf{72} (2005), 062323.

\bibitem{AlickiFannes}
R.~Alicki and M.~Fannes,
\emph{Continuity of quantum conditional information},
J. Phys. A \textbf{37} (2004), L55--L57.

\bibitem{CoverThomas}
T.~M.~Cover and J.~A.~Thomas, \emph{Elements of Information Theory}, 2nd ed.,
Wiley, 2006.

\bibitem{PelissierBornDisintegrationArxiv}
A.~P\'elissier,
\emph{Predictive Record Disintegration and the Quadratic Weight Rule in Renewal Memory},
arXiv placeholder, 2026.

\bibitem{WolfPerez}
M.~M.~Wolf, \emph{Quantum Channels and Operations: Guided Tour}, 2012.

\bibitem{ReehSchlieder}
H.~Reeh and S.~Schlieder, \emph{Bemerkungen zur Unit\"ar\"aquivalenz von
Lorentzinvarianten Feldern}, Nuovo Cimento \textbf{22} (1961), 1051--1068.

\bibitem{StreaterWightman}
R.~F.~Streater and A.~S.~Wightman, \emph{PCT, Spin and Statistics, and All That},
Benjamin, 1964.

\bibitem{BGL2}
D.~Guido and R.~Longo, \emph{An algebraic spin and statistics theorem},
Comm. Math. Phys. \textbf{172} (1995), 517--533.

\bibitem{Borchers}
H.-J.~Borchers,
\emph{The CPT theorem in two-dimensional theories of local observables},
Comm. Math. Phys. \textbf{143} (1992), 315--332.

\bibitem{Wiesbrock}
H.-W.~Wiesbrock, \emph{Half-sided modular inclusions of von Neumann algebras},
Comm. Math. Phys. \textbf{157} (1993), 83--92.

\bibitem{Longo}
R.~Longo, \emph{Real Hilbert subspaces, modular theory, $\\mathrm{SL}(2,\\R)$ and
CFT}, in \emph{Von Neumann Algebras in Sibiu}, Theta, 2008, 33--91.

\bibitem{BGL}
R.~Brunetti, D.~Guido, and R.~Longo, \emph{Modular localization and Wigner
particles}, Rev. Math. Phys. \textbf{14} (2002), 759--785.

\bibitem{EckmannOsterwalder}
J.-P.~Eckmann and K.~Osterwalder, \emph{An application of Tomita's theory of
modular Hilbert algebras: duality for free Bose fields}, J. Funct. Anal.
\textbf{13} (1973), 1--12.

\bibitem{FiglioliniGuido}
F.~Figliolini and D.~Guido, \emph{The Tomita operator for the free scalar field},
Ann. Inst. H. Poincar\'e Phys. Th\'eor. \textbf{51} (1989), 419--435.

\bibitem{Araki}
H.~Araki, \emph{Von Neumann algebras of local observables for free scalar field},
J. Math. Phys. \textbf{5} (1964), 1--13.

\bibitem{Driessler}
W.~Driessler, \emph{Comments on lightlike translations and applications in
relativistic quantum field theory}, Comm. Math. Phys. \textbf{44} (1975), 133--141.

\bibitem{BisognanoWichmann}
J.~J.~Bisognano and E.~H.~Wichmann,
\emph{On the duality condition for a Hermitian scalar field},
J. Math. Phys. \textbf{16} (1975), 985--1007.

\bibitem{CasiniHuertaMyers}
H.~Casini, M.~Huerta, and R.~C.~Myers,
\emph{Towards a derivation of holographic entanglement entropy},
JHEP \textbf{05} (2011), 036.

\bibitem{ArakiZsido2005}
H.~Araki and L.~Zsid\'o,
\emph{Extension of the structure theorem of Borchers and its application to
half-sided modular inclusions}, Rev. Math. Phys. \textbf{17} (2005), 491--543.

\bibitem{Raychaudhuri}
A.~Raychaudhuri, \emph{Relativistic cosmology. I},
Phys. Rev. \textbf{98} (1955), 1123--1126.

\bibitem{Poisson}
E.~Poisson, \emph{A Relativist's Toolkit: The Mathematics of Black-Hole
Mechanics}, Cambridge Univ. Press, 2004.

\bibitem{Noether}
E.~Noether, \emph{Invariante Variationsprobleme}, Nachr. Ges. Wiss. G\"ottingen
(1918), 235--257.

\bibitem{PelissierBridge}
A.~P\'elissier,
\emph{Bridge Geometry from Renewal Memory: A Predictive-Screen Criterion for
ER=EPR}, arXiv placeholder, 2026.

\bibitem{PelissierInflationBaryogenesisArxiv}
A.~P\'elissier,
\emph{Renewal Inflation and Baryogenesis: Plateau Dynamics, Reheating, and the Matter--Antimatter Asymmetry},
arXiv placeholder, 2026.

\bibitem{Starobinsky1980}
A.~A.~Starobinsky, \emph{A new type of isotropic cosmological models without
singularity}, Phys. Lett. B \textbf{91} (1980), 99--102.

\bibitem{Sakharov1968}
A.~D.~Sakharov, \emph{Vacuum quantum fluctuations in curved space and the theory
of gravitation}, Sov. Phys. Dokl. \textbf{12} (1968), 1040--1041.

\bibitem{Adler1982}
S.~L.~Adler, \emph{Einstein gravity as a symmetry-breaking effect in quantum
field theory}, Rev. Mod. Phys. \textbf{54} (1982), 729--766; Erratum
\textbf{55} (1983), 837.

\bibitem{Zee1981}
A.~Zee, \emph{Spontaneously generated gravity}, Phys. Rev. D \textbf{23}
(1981), 858--866.

\bibitem{ONeill}
B.~O'Neill, \emph{Semi-Riemannian Geometry with Applications to Relativity},
Academic Press, 1983.

\bibitem{Wolf}
J.~A.~Wolf, \emph{Spaces of Constant Curvature}, 6th ed., AMS Chelsea, 2011.

\bibitem{PelissierBlackHoles}
A.~P\'elissier,
\emph{Renewal Black Holes, Islands, and the Page Transition},
arXiv placeholder, 2026.

\bibitem{GrayVanhecke}
A.~Gray and L.~Vanhecke,
\emph{Riemannian geometry as determined by the volumes of small geodesic balls},
Acta Math. \textbf{142} (1979), 157--198.

\bibitem{Brewin}
L.~Brewin,
\emph{Riemann normal coordinate expansions using Cadabra},
Class. Quantum Grav. \textbf{26} (2009), 175017.

\bibitem{Vassilevich}
D.~V.~Vassilevich, \emph{Heat kernel expansion: user's manual},
Phys. Rep. \textbf{388} (2003), 279--360.

\bibitem{JacobsonMattingly}
T.~Jacobson and D.~Mattingly,
\emph{Gravity with a dynamical preferred frame},
Phys. Rev. D \textbf{64} (2001), 024028.

\bibitem{Jacobson2008aether}
T.~Jacobson, \emph{Einstein-aether gravity: a status report},
PoS \textbf{QG-Ph2007} (2007), 020.

\bibitem{Isserlis}
L.~Isserlis,
\emph{On a formula for the product-moment coefficient of any order of a normal
frequency distribution in any number of variables},
Biometrika \textbf{12} (1918), 134--139.

\bibitem{VisserInducedGravity}
M.~Visser,
\emph{Sakharov's induced gravity: a modern perspective},
Mod. Phys. Lett. A \textbf{17} (2002), 977--992.

\bibitem{GilkeyHeat}
P.~B.~Gilkey, \emph{Invariance Theory, the Heat Equation, and the Atiyah--Singer
Index Theorem}, 2nd ed., CRC Press, 1995.

\bibitem{AvramidiHeat}
I.~G.~Avramidi, \emph{Heat Kernel and Quantum Gravity}, Springer, 2000.

\bibitem{BarvinskyVilkovisky}
A.~O.~Barvinsky and G.~A.~Vilkovisky,
\emph{The generalized Schwinger--DeWitt technique in gauge theories and quantum
gravity}, Phys. Rep. \textbf{119} (1985), 1--74.

\bibitem{Gusynin}
V.~P.~Gusynin,
\emph{Seeley--Gilkey coefficients for fourth-order operators on a Riemannian
manifold}, Nucl. Phys. B \textbf{333} (1990), 296--316.

\bibitem{Gray}
A.~Gray, \emph{Tubes}, 2nd ed., Progress in Mathematics \textbf{221},
Birkh\"auser, 2004.

\bibitem{Horava}
P.~Ho\v{r}ava,
\emph{Quantum gravity at a Lifshitz point},
Phys. Rev. D \textbf{79} (2009), 084008.

\bibitem{WaldEntropy}
R.~M.~Wald, \emph{Black hole entropy is the Noether charge},
Phys. Rev. D \textbf{48} (1993), R3427--R3431.

\bibitem{IyerWald}
V.~Iyer and R.~M.~Wald,
\emph{Some properties of Noether charge and a proposal for dynamical black hole
entropy}, Phys. Rev. D \textbf{50} (1994), 846--864.

\bibitem{PelissierPhenomenology}
A.~P\'elissier, \emph{Observable Consequences of Renewal Memory},
arXiv placeholder, 2026.

\bibitem{DemboZeitouni}
A.~Dembo and O.~Zeitouni,
\emph{Large Deviations Techniques and Applications}, 2nd ed.,
Springer, 1998.

\bibitem{Touchette2009}
H.~Touchette,
\emph{The large deviation approach to statistical mechanics},
Phys. Rep. \textbf{478} (2009), 1--69.

\bibitem{MetzlerKlafter}
R.~Metzler and J.~Klafter,
\emph{The random walk's guide to anomalous diffusion: a fractional dynamics approach},
Phys. Rep. \textbf{339} (2000), 1--77.

\bibitem{MeerschaertScheffler}
M.~M.~Meerschaert and H.-P.~Scheffler,
\emph{Limit theorems for continuous-time random walks with infinite mean waiting times},
J. Appl. Probab. \textbf{41} (2004), 623--638.

\bibitem{Podlubny}
I.~Podlubny,
\emph{Fractional Differential Equations},
Academic Press, 1999.

\bibitem{Pollard1948}
H.~Pollard,
\emph{The completely monotonic character of the Mittag--Leffler function
\(E_\alpha(-x)\)},
Bull. Amer. Math. Soc. \textbf{54} (1948), 1115--1116.

\bibitem{Schneider1996}
W.~R.~Schneider,
\emph{Completely monotone generalized Mittag--Leffler functions},
Exposition. Math. \textbf{14} (1996), 3--16.

\bibitem{GripenbergLondenStaffans}
G.~Gripenberg, S.-O.~Londen, and O.~Staffans,
\emph{Volterra Integral and Functional Equations},
Cambridge University Press, 1990.

\bibitem{Hill}
T.~L.~Hill, \emph{Free Energy Transduction and Biochemical Cycle Kinetics},
Springer, 1989.

\bibitem{Schnakenberg}
J.~Schnakenberg, \emph{Network theory of microscopic and macroscopic behavior
of master equation systems}, Rev. Mod. Phys. \textbf{48} (1976), 571--585.

\bibitem{SchillingSongVondracek}
R.~L.~Schilling, R.~Song, and Z.~Vondra\v{c}ek,
\emph{Bernstein Functions: Theory and Applications}, 2nd ed.,
De Gruyter Studies in Mathematics \textbf{37}, De Gruyter, 2012.

\bibitem{BaratoSeifert2015}
A.~C.~Barato and U.~Seifert,
\emph{Thermodynamic uncertainty relation for biomolecular processes},
Phys. Rev. Lett. \textbf{114} (2015), 158101.

\bibitem{Gingrich2016}
T.~R.~Gingrich, J.~M.~Horowitz, N.~Perunov, and J.~L.~England,
\emph{Dissipation bounds all steady-state current fluctuations},
Phys. Rev. Lett. \textbf{116} (2016), 120601.

\bibitem{HorowitzGingrich2020}
J.~M.~Horowitz and T.~R.~Gingrich,
\emph{Thermodynamic uncertainty relations constrain non-equilibrium fluctuations},
Nature Phys. \textbf{16} (2020), 15--20.

\bibitem{ValiviitaMajerottoMaartens}
J.~V\"aliviita, E.~Majerotto, and R.~Maartens,
\emph{Large-scale instability in interacting dark energy and dark matter fluids},
J. Cosmol. Astropart. Phys. \textbf{07} (2008), 020.

\bibitem{ClemsonEtAl}
T.~Clemson, K.~Koyama, G.-B.~Zhao, R.~Maartens, and J.~V\"aliviita,
\emph{Interacting dark energy---constraints and degeneracies},
Phys. Rev. D \textbf{85} (2012), 043007.

\bibitem{LebowitzSpohn}
J.~L.~Lebowitz and H.~Spohn,
\emph{A Gallavotti--Cohen-type symmetry in the large deviation functional for
stochastic dynamics},
J. Stat. Phys. \textbf{95} (1999), 333--365.

\bibitem{AndrieuxGaspard2008}
D.~Andrieux and P.~Gaspard,
\emph{The fluctuation theorem for currents in semi-Markov processes},
J. Stat. Mech. Theory Exp. (2008), P11007.

\bibitem{FaggionatoSemiMarkov}
A.~Faggionato,
\emph{Large deviation principles and fluctuation theorems for currents in
semi-Markov processes},
arXiv:1709.05653 (2017).

\bibitem{WeinbergCC}
S.~Weinberg,
\emph{The cosmological constant problem},
Rev. Mod. Phys. \textbf{61} (1989), 1--23.

\bibitem{RatraPeebles1988}
B.~Ratra and P.~J.~E.~Peebles,
\emph{Cosmological consequences of a rolling homogeneous scalar field},
Phys. Rev. D \textbf{37} (1988), 3406--3427.

\bibitem{Wetterich1988}
C.~Wetterich,
\emph{Cosmology and the fate of dilatation symmetry},
Nucl. Phys. B \textbf{302} (1988), 668--696.

\bibitem{CaldwellDaveSteinhardt}
R.~R.~Caldwell, R.~Dave, and P.~J.~Steinhardt,
\emph{Cosmological imprint of an energy component with general equation of state},
Phys. Rev. Lett. \textbf{80} (1998), 1582--1585.

\bibitem{ArmendarizPiconKessence}
C.~Armendariz-Picon, V.~Mukhanov, and P.~J.~Steinhardt,
\emph{Essentials of k-essence},
Phys. Rev. D \textbf{63} (2001), 103510.

\bibitem{CaldwellPhantom}
R.~R.~Caldwell, M.~Kamionkowski, and N.~N.~Weinberg,
\emph{Phantom energy: dark energy with $w<-1$ causes a cosmic doomsday},
Phys. Rev. Lett. \textbf{91} (2003), 071301.

\bibitem{Amendola2000}
L.~Amendola,
\emph{Coupled quintessence},
Phys. Rev. D \textbf{62} (2000), 043511.

\end{thebibliography}
\end{document}